% Generated human-review packet template.  Source records, Lean statements,
% and raw-source-to-Spec screening fields are substituted by
% scripts/review_dashboard_packet.py.
% Do not hand-edit generated packets; annotate the PDF or reconcile notes into
% the paper-local human-review log.
\documentclass[10pt]{article}
\usepackage[margin=0.43in]{geometry}
\usepackage{fontspec}
\setmainfont{Latin Modern Roman}
\setmonofont{DejaVu Sans Mono}
\usepackage{hyperref}
\usepackage{amssymb}
\usepackage{fvextra}
\usepackage{enumitem}
\setlength{\parindent}{0pt}
\setlength{\parskip}{0.15em}
\linespread{0.92}
\hypersetup{colorlinks=true,linkcolor=blue,urlcolor=blue}
\DefineVerbatimEnvironment{ReviewVerbatim}{Verbatim}{breaklines=true,breakanywhere=true,fontsize=\scriptsize}
\newcommand{\reviewerbox}{%
  \par\noindent\fbox{\parbox{0.96\linewidth}{\rule{0pt}{0.38in}}}\par
}
\newcommand{\reviewmatch}[1]{%
  \par\noindent\textbf{Reviewer check:} $\square$\; #1\par
  \noindent\emph{If left unchecked, explain the discrepancy or uncertainty below.}\par
}

\begin{document}
\fontsize{9}{10}\selectfont

\begin{center}
{\LARGE Human review packet: DGD26AdmissionsPredictability}\\[0.35em]
\small Generated from byte-pinned source excerpts, the expanded PaperInterface
specifications, and hash-bound raw-source-to-Spec screening records.
\end{center}

\noindent\textbf{Paper:} DGD26AdmissionsPredictability\\
\textbf{Paper id:} \texttt{DGD26AdmissionsPredictability}\\
\textbf{Source version:} AAAI 2026 paper; pinned arXiv TeX source archive\\
\textbf{Source record:} \texttt{audit/paper\_statement\_map.json}\\
\textbf{Rows:} 63\\
\textbf{Generated:} 2026-08-18\\
\textbf{Source URL:} \href{https://ojs.aaai.org/index.php/AAAI/article/view/41179}{official source}

\noindent\fbox{\parbox{0.96\linewidth}{\textbf{Review status:} 0/63 source-claim screens; 0/23 library prerequisite screens. This is a review worksheet while those direct checks remain pending; it is not a final validation receipt.}}\par



\section*{How to use this packet}
For each source claim, compare the exact source excerpts with the expanded
PaperInterface specification. Mark up this PDF or fill the blank reviewer box in the TeX source;
return the annotated file for reconciliation into the paper-local human-review
log.

\section*{Open the interactive dashboard}
The interactive dashboard is an optional alternative to using this PDF; it
presents the same review surface rather than an additional review requirement.
After forking and cloning (or pulling) EconCSLib, run the following from the
repository root. The cache command is needed only the first time in a new clone.
\begin{ReviewVerbatim}
cd <your-EconCSLib-checkout>
lake exe cache get
python3 scripts/review_dashboard.py --paper DGD26AdmissionsPredictability --serve
\end{ReviewVerbatim}
Open \url{http://127.0.0.1:8765/} in a browser and keep that terminal running
while reviewing. The dashboard records saved annotations in this paper's
reviewer-owned review record.

\section*{Contents}
\begin{itemize}[leftmargin=1.2em]
\item \hyperlink{library-section-bee5354f1438f98f}{Semantic library prerequisites (23; not paper claims)}
\begin{itemize}[leftmargin=1.2em]
\item \hyperlink{library-prerequisite-aebe54633da60351}{EconCSLib.FiniteChoice.ChoiceRule}
\item \hyperlink{library-prerequisite-c51ffba76c13d204}{EconCSLib.FiniteChoice.Consistent}
\item \hyperlink{library-prerequisite-e66298c8dac25baa}{EconCSLib.FiniteChoice.choiceGainTerm}
\item \hyperlink{library-prerequisite-ba6c40008598cadf}{EconCSLib.FiniteChoice.choiceLossTerm}
\item \hyperlink{library-prerequisite-66d5b4a54f47387c}{EconCSLib.FiniteChoice.choiceDistance}
\item \hyperlink{library-prerequisite-29c4c5d3d08823af}{EconCSLib.FiniteChoice.DUnstable}
\item \hyperlink{library-prerequisite-542ade70cc681808}{EconCSLib.FiniteChoice.Feasible}
\item \hyperlink{library-prerequisite-ac076d846344765a}{EconCSLib.FiniteChoice.borderlineSet}
\item \hyperlink{library-prerequisite-84ef92453e2eeb0e}{EconCSLib.FiniteChoice.waitlistedSet}
\item \hyperlink{library-prerequisite-5f0109fb4b0fbb5d}{EconCSLib.FiniteChoice.GeneralVariabilityAtMost}
\item \hyperlink{library-prerequisite-94743536b7e3fe31}{EconCSLib.FiniteChoice.GeneralVariabilityExactly}
\item \hyperlink{library-prerequisite-0de9dc90c34c9fb6}{EconCSLib.FiniteChoice.Independent}
\item \hyperlink{library-prerequisite-8ee5c1a19c8d4e3a}{EconCSLib.FiniteChoice.Monotonic}
\item \hyperlink{library-prerequisite-f02e4b75ca5571bf}{EconCSLib.FiniteChoice.QAcceptant}
\item \hyperlink{library-prerequisite-88c2ddc1c549c716}{EconCSLib.FiniteChoice.StrictTotalOrder}
\item \hyperlink{library-prerequisite-286eac89d49c4d81}{EconCSLib.FiniteChoice.QRepresentative}
\item \hyperlink{library-prerequisite-542d5eef4a694a7d}{EconCSLib.FiniteChoice.Substitutable}
\item \hyperlink{library-prerequisite-bb61a90095edae4b}{EconCSLib.FiniteChoice.TightlyDUnstable}
\item \hyperlink{library-prerequisite-aa69285c37fbf372}{EconCSLib.FiniteChoice.VariabilityAtMost}
\item \hyperlink{library-prerequisite-4ccdc8565ec23c71}{EconCSLib.FiniteChoice.VariabilityExactly}
\item \hyperlink{library-prerequisite-e55d0830fae724cb}{EconCSLib.FiniteChoice.ZeroUnstable}
\item \hyperlink{library-prerequisite-4a3989697f894599}{EconCSLib.FiniteChoice.linearTopQChoice}
\item \hyperlink{library-prerequisite-3767c56c4f2cc69a}{EconCSLib.FiniteChoice.sequentialComposition}
\end{itemize}
\item \hyperlink{source-section-f10b5f68e4acf3dc}{Source claims (63)}
\begin{itemize}[leftmargin=1.2em]
\item \hyperlink{source-claim-03de8677e0f95903}{paper\_fixed\_threshold\_formula\_statementSpec}
\item \hyperlink{source-claim-685dcba90250bf7c}{paper\_choice\_function\_model\_definition\_statementSpec}
\item \hyperlink{source-claim-d396bf72ee1812a2}{paper\_definition\_choice\_labelSpec}
\item \hyperlink{source-claim-51b92e4a229b02a5}{paper\_q\_acceptance\_definition\_statementSpec}
\item \hyperlink{source-claim-9a3be13c89f94ffe}{paper\_total\_order\_definition\_statementSpec}
\item \hyperlink{source-claim-a21144374ad8e28c}{paper\_definition\_choice\_distanceSpec}
\item \hyperlink{source-claim-398a4d87fcced6ba}{paper\_d\_instability\_definition\_statementSpec}
\item \hyperlink{source-claim-62d00d89c6cc151a}{paper\_tight\_d\_instability\_definition\_statementSpec}
\item \hyperlink{source-claim-be4f0b5bb7e1dd6b}{paper\_variability\_exactly\_definition\_statementSpec}
\item \hyperlink{source-claim-12cc974d07b19ac2}{paper\_ml\_representation\_definition\_statementSpec}
\item \hyperlink{source-claim-1e70e813c2db1d8e}{paper\_rank\_threshold\_formula\_statementSpec}
\item \hyperlink{source-claim-d17644eda5ebe53f}{paper\_substitutability\_definition\_statementSpec}
\item \hyperlink{source-claim-69b86836fa7d4d06}{paper\_definition\_sequential\_compositionSpec}
\item \hyperlink{source-claim-1b6dc1033989d2bf}{paper\_monotonicity\_definition\_statementSpec}
\item \hyperlink{source-claim-899ae3aec5210e2c}{paper\_consistency\_definition\_statementSpec}
\item \hyperlink{source-claim-ca1d4873227a50e3}{paper\_q\_representativeness\_definition\_statementSpec}
\item \hyperlink{source-claim-3a97b71dd915eed1}{paper\_zero\_instability\_definition\_statementSpec}
\item \hyperlink{source-claim-e633cfedc7bfb493}{paper\_independence\_definition\_statementSpec}
\item \hyperlink{source-claim-8e553b3e026d261e}{paper\_definition\_borderline\_setSpec}
\item \hyperlink{source-claim-1c2136cbe341a833}{paper\_definition\_waitlisted\_setSpec}
\item \hyperlink{source-claim-d2f6a286dc057af6}{paper\_general\_variability\_exactly\_definition\_statementSpec}
\item \hyperlink{source-claim-a009d7013e95842d}{paper\_lap\_model\_definition\_statementSpec}
\item \hyperlink{source-claim-f902fe2aa64a28cb}{paper\_lap\_slot\_below\_definition\_statementSpec}
\item \hyperlink{source-claim-ac931f819ad35103}{paper\_lap\_same\_slot\_order\_definition\_statementSpec}
\item \hyperlink{source-claim-be70d394f2bee79c}{paper\_lap\_assignment\_choice\_formula\_statementSpec}
\item \hyperlink{source-claim-b05952c99a84b6af}{paper\_independent\_zero\_unstable\_statementSpec}
\item \hyperlink{source-claim-d920f9b7847b2baa}{paper\_independent\_zero\_unstable\_corollary\_statementSpec}
\item \hyperlink{source-claim-18a641b7c4365d27}{paper\_no\_consistent\_positive\_tightly\_even\_statementSpec}
\item \hyperlink{source-claim-1daa198b9d0a1bc3}{paper\_zero\_distance\_statementSpec}
\item \hyperlink{source-claim-760d01f2c31e3672}{paper\_calculating\_instability\_statementSpec}
\item \hyperlink{source-claim-72ba750e0e18812e}{paper\_corrected\_consistency\_of\_removable\_sets\_statementSpec}
\item \hyperlink{source-claim-7d8732899c3f4210}{paper\_lap\_strictly\_higher\_slot\_applicant\_assigned\_statementSpec}
\item \hyperlink{source-claim-5e3bd8769e715d3c}{paper\_lap\_assigned\_strictly\_outranks\_rejected\_statementSpec}
\item \hyperlink{source-claim-14ff35045a22028c}{paper\_monotonicity\_term\_statementSpec}
\item \hyperlink{source-claim-43c4c9b7821d369a}{paper\_monotonicity\_q\_acceptance\_incompatible\_statementSpec}
\item \hyperlink{source-claim-f926138cb2114337}{paper\_acceptant\_one\_instability\_variability\_borderline\_eq\_waitlisted\_after\_changing\_insert\_statementSpec}
\item \hyperlink{source-claim-ef4d58145103cd0c}{paper\_substitutability\_term\_statementSpec}
\item \hyperlink{source-claim-ca2a099855204aca}{paper\_non\_substitutable\_single\_add\_statementSpec}
\item \hyperlink{source-claim-bc7eeffb236e6569}{paper\_ml\_fixed\_threshold\_representation\_zero\_unstable\_statementSpec}
\item \hyperlink{source-claim-d21050cec8d233b0}{paper\_ml\_rank\_threshold\_can\_represent\_exact\_one\_statementSpec}
\item \hyperlink{source-claim-47a9f492a25c4e98}{paper\_ml\_rank\_threshold\_representation\_instability\_bound\_statementSpec}
\item \hyperlink{source-claim-98ed151e7543b81d}{paper\_ml\_rank\_threshold\_representation\_variability\_bound\_statementSpec}
\item \hyperlink{source-claim-396063f8f4209701}{paper\_program\_classes\_one\_instability\_statementSpec}
\item \hyperlink{source-claim-f595f3e6af0e1b40}{paper\_screened\_open\_variability\_one\_statementSpec}
\item \hyperlink{source-claim-8afe5dc5804dcb73}{paper\_screened\_open\_dia\_variability\_two\_statementSpec}
\item \hyperlink{source-claim-68c8be83630deedf}{paper\_educational\_option\_variability\_three\_statementSpec}
\item \hyperlink{source-claim-a241fbe5321a55d3}{paper\_educational\_option\_dia\_variability\_six\_statementSpec}
\item \hyperlink{source-claim-207dcf4700d4343b}{paper\_sequential\_additive\_variability\_bound\_statementSpec}
\item \hyperlink{source-claim-fff538b4a617e0f3}{paper\_append\_remove\_variability\_exact\_equivalence\_statementSpec}
\item \hyperlink{source-claim-e1e0f5f892109672}{paper\_choice\_distance\_triangle\_statementSpec}
\item \hyperlink{source-claim-f8de55f9c5760a87}{paper\_even\_distance\_iff\_inconsistent\_statementSpec}
\item \hyperlink{source-claim-3e732abc6f7366e9}{paper\_independent\_substitutable\_monotonic\_statementSpec}
\item \hyperlink{source-claim-dac748564d615cb6}{paper\_no\_zero\_instability\_under\_capacity\_statementSpec}
\item \hyperlink{source-claim-4efb65f58e4bc480}{paper\_substitutability\_one\_instability\_equivalence\_statementSpec}
\item \hyperlink{source-claim-d2442dd19c06f77a}{paper\_theorem1\_tight\_all\_d\_statementSpec}
\item \hyperlink{source-claim-b3ee09f1e9e6b10d}{paper\_lap\_assignment\_one\_instability\_statementSpec}
\item \hyperlink{source-claim-a0520310d6532be1}{paper\_lap\_assignment\_slot\_order\_class\_variability\_of\_unique\_global\_optima\_statementSpec}
\item \hyperlink{source-claim-a9c5e5d5c616ab58}{paper\_q\_representative\_forward\_exact\_statementSpec}
\item \hyperlink{source-claim-3d12c759820fa6e4}{paper\_q\_representative\_converse\_exact\_statementSpec}
\item \hyperlink{source-claim-5fcefe7f9430bc35}{paper\_sequential\_composition\_substitutable\_statementSpec}
\item \hyperlink{source-claim-b564e283fe9ce18d}{paper\_q\_acceptant\_substitutable\_consistent\_statementSpec}
\item \hyperlink{source-claim-7c749451f4caf111}{paper\_sequential\_q\_representative\_variability\_range\_statementSpec}
\item \hyperlink{source-claim-26bfa3509b616c5d}{paper\_variability\_one\_iff\_single\_order\_statementSpec}
\end{itemize}
\end{itemize}

\clearpage
\hypertarget{library-section-bee5354f1438f98f}{}
\section*{Material library prerequisites}
\hypertarget{library-prerequisite-aebe54633da60351}{}
\subsection*{\small\ttfamily\raggedright EconCSLib.\allowbreak{}FiniteChoice.\allowbreak{}ChoiceRule}
\textbf{Source connection:} no explicit byte-pinned paper source connection is registered.
\paragraph{Lean-expanded library semantic target}
\begin{ReviewVerbatim}
(α : Type u_2) → [DecidableEq α] → Finset α → Finset α
\end{ReviewVerbatim}

\paragraph{Exact Lean library declaration}
\begin{ReviewVerbatim}
abbrev ChoiceRule (α : Type*) [DecidableEq α] := Finset α → Finset α
\end{ReviewVerbatim}

\paragraph{Recorded source-to-library screening}
\textbf{Verdict:} \texttt{not recorded}
\reviewmatch{Matches source input}
\noindent\textbf{Reviewer annotation}\par
\reviewerbox
\clearpage
\hypertarget{library-prerequisite-c51ffba76c13d204}{}
\subsection*{\small\ttfamily\raggedright EconCSLib.\allowbreak{}FiniteChoice.\allowbreak{}Consistent}
\textbf{Source connection:} no explicit byte-pinned paper source connection is registered.
\paragraph{Lean-expanded library semantic target}
\begin{ReviewVerbatim}
∀ {α : Type u_1} [inst : DecidableEq α] (C : EconCSLib.FiniteChoice.ChoiceRule α) {X₁ X₂ : Finset α},
  C X₂ ⊆ X₁ → X₁ ⊆ X₂ → C X₂ = C X₁
\end{ReviewVerbatim}

\paragraph{Exact Lean library declaration}
\begin{ReviewVerbatim}
def Consistent (C : ChoiceRule α) : Prop :=
  ∀ {X₁ X₂}, C X₂ ⊆ X₁ → X₁ ⊆ X₂ → C X₂ = C X₁
\end{ReviewVerbatim}

\paragraph{Recorded source-to-library screening}
\textbf{Verdict:} \texttt{not recorded}
\reviewmatch{Matches source input}
\noindent\textbf{Reviewer annotation}\par
\reviewerbox
\clearpage
\hypertarget{library-prerequisite-e66298c8dac25baa}{}
\subsection*{\small\ttfamily\raggedright EconCSLib.\allowbreak{}FiniteChoice.\allowbreak{}choiceGainTerm}
\textbf{Source connection:} no explicit byte-pinned paper source connection is registered.
\paragraph{Lean-expanded library semantic target}
\begin{ReviewVerbatim}
{α : Type u_1} →
  [inst : DecidableEq α] → (C : EconCSLib.FiniteChoice.ChoiceRule α) → (X₁ X₂ : Finset α) → ((X₁ ∩ C X₂) \ C X₁).card
\end{ReviewVerbatim}

\paragraph{Exact Lean library declaration}
\begin{ReviewVerbatim}
def choiceGainTerm (C : ChoiceRule α) (X₁ X₂ : Finset α) : ℕ :=
  ((X₁ ∩ C X₂) \ C X₁).card
\end{ReviewVerbatim}

\paragraph{Recorded source-to-library screening}
\textbf{Verdict:} \texttt{not recorded}
\reviewmatch{Matches source input}
\noindent\textbf{Reviewer annotation}\par
\reviewerbox
\clearpage
\hypertarget{library-prerequisite-ba6c40008598cadf}{}
\subsection*{\small\ttfamily\raggedright EconCSLib.\allowbreak{}FiniteChoice.\allowbreak{}choiceLossTerm}
\textbf{Source connection:} no explicit byte-pinned paper source connection is registered.
\paragraph{Lean-expanded library semantic target}
\begin{ReviewVerbatim}
{α : Type u_1} →
  [inst : DecidableEq α] → (C : EconCSLib.FiniteChoice.ChoiceRule α) → (X₁ X₂ : Finset α) → (C X₁ \ C X₂).card
\end{ReviewVerbatim}

\paragraph{Exact Lean library declaration}
\begin{ReviewVerbatim}
def choiceLossTerm (C : ChoiceRule α) (X₁ X₂ : Finset α) : ℕ :=
  (C X₁ \ C X₂).card
\end{ReviewVerbatim}

\paragraph{Recorded source-to-library screening}
\textbf{Verdict:} \texttt{not recorded}
\reviewmatch{Matches source input}
\noindent\textbf{Reviewer annotation}\par
\reviewerbox
\clearpage
\hypertarget{library-prerequisite-66d5b4a54f47387c}{}
\subsection*{\small\ttfamily\raggedright EconCSLib.\allowbreak{}FiniteChoice.\allowbreak{}choiceDistance}
\textbf{Source connection:} no explicit byte-pinned paper source connection is registered.
\paragraph{Lean-expanded library semantic target}
\begin{ReviewVerbatim}
{α : Type u_1} →
  [inst : DecidableEq α] →
    (C : EconCSLib.FiniteChoice.ChoiceRule α) →
      (X₁ X₂ : Finset α) → EconCSLib.FiniteChoice.choiceGainTerm C X₁ X₂ + EconCSLib.FiniteChoice.choiceLossTerm C X₁ X₂
\end{ReviewVerbatim}

\paragraph{Exact Lean library declaration}
\begin{ReviewVerbatim}
def choiceDistance (C : ChoiceRule α) (X₁ X₂ : Finset α) : ℕ :=
  choiceGainTerm C X₁ X₂ + choiceLossTerm C X₁ X₂
\end{ReviewVerbatim}

\paragraph{Recorded source-to-library screening}
\textbf{Verdict:} \texttt{not recorded}
\reviewmatch{Matches source input}
\noindent\textbf{Reviewer annotation}\par
\reviewerbox
\clearpage
\hypertarget{library-prerequisite-29c4c5d3d08823af}{}
\subsection*{\small\ttfamily\raggedright EconCSLib.\allowbreak{}FiniteChoice.\allowbreak{}DUnstable}
\textbf{Source connection:} no explicit byte-pinned paper source connection is registered.
\paragraph{Lean-expanded library semantic target}
\begin{ReviewVerbatim}
∀ {α : Type u_1} [inst : DecidableEq α] (d : ℕ) (C : EconCSLib.FiniteChoice.ChoiceRule α) (X : Finset α),
  ∀ x ∉ X, EconCSLib.FiniteChoice.choiceDistance C X (insert x X) ≤ d
\end{ReviewVerbatim}

\paragraph{Exact Lean library declaration}
\begin{ReviewVerbatim}
def DUnstable (d : ℕ) (C : ChoiceRule α) : Prop :=
  ∀ X x, x ∉ X → choiceDistance C X (insert x X) ≤ d
\end{ReviewVerbatim}

\paragraph{Recorded source-to-library screening}
\textbf{Verdict:} \texttt{not recorded}
\reviewmatch{Matches source input}
\noindent\textbf{Reviewer annotation}\par
\reviewerbox
\clearpage
\hypertarget{library-prerequisite-542ade70cc681808}{}
\subsection*{\small\ttfamily\raggedright EconCSLib.\allowbreak{}FiniteChoice.\allowbreak{}Feasible}
\textbf{Source connection:} no explicit byte-pinned paper source connection is registered.
\paragraph{Lean-expanded library semantic target}
\begin{ReviewVerbatim}
∀ {α : Type u_1} [inst : DecidableEq α] (C : EconCSLib.FiniteChoice.ChoiceRule α) (X : Finset α), C X ⊆ X
\end{ReviewVerbatim}

\paragraph{Exact Lean library declaration}
\begin{ReviewVerbatim}
def Feasible (C : ChoiceRule α) : Prop :=
  ∀ X, C X ⊆ X
\end{ReviewVerbatim}

\paragraph{Recorded source-to-library screening}
\textbf{Verdict:} \texttt{not recorded}
\reviewmatch{Matches source input}
\noindent\textbf{Reviewer annotation}\par
\reviewerbox
\clearpage
\hypertarget{library-prerequisite-ac076d846344765a}{}
\subsection*{\small\ttfamily\raggedright EconCSLib.\allowbreak{}FiniteChoice.\allowbreak{}borderlineSet}
\textbf{Source connection:} no explicit byte-pinned paper source connection is registered.
\paragraph{Lean-expanded library semantic target}
\begin{ReviewVerbatim}
{α : Type u_1} →
  [inst : DecidableEq α] →
    [inst_1 : Fintype α] →
      (C : EconCSLib.FiniteChoice.ChoiceRule α) → (X : Finset α) → Finset.univ.biUnion fun x => C X \ C (insert x X)
\end{ReviewVerbatim}

\paragraph{Exact Lean library declaration}
\begin{ReviewVerbatim}
library declaration is not registered or discoverable
\end{ReviewVerbatim}

\paragraph{Recorded source-to-library screening}
\textbf{Verdict:} \texttt{not recorded}
\reviewmatch{Matches source input}
\noindent\textbf{Reviewer annotation}\par
\reviewerbox
\clearpage
\hypertarget{library-prerequisite-84ef92453e2eeb0e}{}
\subsection*{\small\ttfamily\raggedright EconCSLib.\allowbreak{}FiniteChoice.\allowbreak{}waitlistedSet}
\textbf{Source connection:} no explicit byte-pinned paper source connection is registered.
\paragraph{Lean-expanded library semantic target}
\begin{ReviewVerbatim}
{α : Type u_1} →
  [inst : DecidableEq α] →
    [Fintype α] → (C : EconCSLib.FiniteChoice.ChoiceRule α) → (X : Finset α) → X.biUnion fun x => C (X.erase x) \ C X
\end{ReviewVerbatim}

\paragraph{Exact Lean library declaration}
\begin{ReviewVerbatim}
def waitlistedSet [Fintype α] (C : ChoiceRule α) (X : Finset α) : Finset α :=
  X.biUnion fun x => C (X.erase x) \ C X
\end{ReviewVerbatim}

\paragraph{Recorded source-to-library screening}
\textbf{Verdict:} \texttt{not recorded}
\reviewmatch{Matches source input}
\noindent\textbf{Reviewer annotation}\par
\reviewerbox
\clearpage
\hypertarget{library-prerequisite-5f0109fb4b0fbb5d}{}
\subsection*{\small\ttfamily\raggedright EconCSLib.\allowbreak{}FiniteChoice.\allowbreak{}GeneralVariabilityAtMost}
\textbf{Source connection:} no explicit byte-pinned paper source connection is registered.
\paragraph{Lean-expanded library semantic target}
\begin{ReviewVerbatim}
∀ {α : Type u_1} [inst : DecidableEq α] [inst_1 : Fintype α] (m : ℕ) (C : EconCSLib.FiniteChoice.ChoiceRule α)
  (X : Finset α),
  (EconCSLib.FiniteChoice.borderlineSet C X).card ≤ m ∧ (EconCSLib.FiniteChoice.waitlistedSet C X).card ≤ m
\end{ReviewVerbatim}

\paragraph{Exact Lean library declaration}
\begin{ReviewVerbatim}
def GeneralVariabilityAtMost [Fintype α] (m : ℕ) (C : ChoiceRule α) : Prop :=
  ∀ X, (borderlineSet C X).card ≤ m ∧ (waitlistedSet C X).card ≤ m
\end{ReviewVerbatim}

\paragraph{Recorded source-to-library screening}
\textbf{Verdict:} \texttt{not recorded}
\reviewmatch{Matches source input}
\noindent\textbf{Reviewer annotation}\par
\reviewerbox
\clearpage
\hypertarget{library-prerequisite-94743536b7e3fe31}{}
\subsection*{\small\ttfamily\raggedright EconCSLib.\allowbreak{}FiniteChoice.\allowbreak{}GeneralVariabilityExactly}
\textbf{Source connection:} no explicit byte-pinned paper source connection is registered.
\paragraph{Lean-expanded library semantic target}
\begin{ReviewVerbatim}
∀ {α : Type u_1} [inst : DecidableEq α] [inst_1 : Fintype α] (m : ℕ) (C : EconCSLib.FiniteChoice.ChoiceRule α),
  EconCSLib.FiniteChoice.GeneralVariabilityAtMost m C ∧
    ((∃ X, (EconCSLib.FiniteChoice.borderlineSet C X).card = m) ∨
      ∃ X, (EconCSLib.FiniteChoice.waitlistedSet C X).card = m)
\end{ReviewVerbatim}

\paragraph{Exact Lean library declaration}
\begin{ReviewVerbatim}
def GeneralVariabilityExactly [Fintype α] (m : ℕ) (C : ChoiceRule α) : Prop :=
  GeneralVariabilityAtMost m C ∧
    ((∃ X, (borderlineSet C X).card = m) ∨
      (∃ X, (waitlistedSet C X).card = m))
\end{ReviewVerbatim}

\paragraph{Recorded source-to-library screening}
\textbf{Verdict:} \texttt{not recorded}
\reviewmatch{Matches source input}
\noindent\textbf{Reviewer annotation}\par
\reviewerbox
\clearpage
\hypertarget{library-prerequisite-0de9dc90c34c9fb6}{}
\subsection*{\small\ttfamily\raggedright EconCSLib.\allowbreak{}FiniteChoice.\allowbreak{}Independent}
\textbf{Source connection:} no explicit byte-pinned paper source connection is registered.
\paragraph{Lean-expanded library semantic target}
\begin{ReviewVerbatim}
∀ {α : Type u_1} [inst : DecidableEq α] (C : EconCSLib.FiniteChoice.ChoiceRule α) (x : α),
  (∀ (X : Finset α), x ∈ X → x ∈ C X) ∨ ∀ (X : Finset α), x ∈ X → x ∉ C X
\end{ReviewVerbatim}

\paragraph{Exact Lean library declaration}
\begin{ReviewVerbatim}
def Independent (C : ChoiceRule α) : Prop :=
  ∀ x, (∀ X, x ∈ X → x ∈ C X) ∨ (∀ X, x ∈ X → x ∉ C X)
\end{ReviewVerbatim}

\paragraph{Recorded source-to-library screening}
\textbf{Verdict:} \texttt{not recorded}
\reviewmatch{Matches source input}
\noindent\textbf{Reviewer annotation}\par
\reviewerbox
\clearpage
\hypertarget{library-prerequisite-8ee5c1a19c8d4e3a}{}
\subsection*{\small\ttfamily\raggedright EconCSLib.\allowbreak{}FiniteChoice.\allowbreak{}Monotonic}
\textbf{Source connection:} no explicit byte-pinned paper source connection is registered.
\paragraph{Lean-expanded library semantic target}
\begin{ReviewVerbatim}
∀ {α : Type u_1} [inst : DecidableEq α] (C : EconCSLib.FiniteChoice.ChoiceRule α) {X₁ X₂ : Finset α},
  X₁ ⊆ X₂ → C X₁ ⊆ C X₂
\end{ReviewVerbatim}

\paragraph{Exact Lean library declaration}
\begin{ReviewVerbatim}
def Monotonic (C : ChoiceRule α) : Prop :=
  ∀ {X₁ X₂}, X₁ ⊆ X₂ → C X₁ ⊆ C X₂
\end{ReviewVerbatim}

\paragraph{Recorded source-to-library screening}
\textbf{Verdict:} \texttt{not recorded}
\reviewmatch{Matches source input}
\noindent\textbf{Reviewer annotation}\par
\reviewerbox
\clearpage
\hypertarget{library-prerequisite-f02e4b75ca5571bf}{}
\subsection*{\small\ttfamily\raggedright EconCSLib.\allowbreak{}FiniteChoice.\allowbreak{}QAcceptant}
\textbf{Source connection:} no explicit byte-pinned paper source connection is registered.
\paragraph{Lean-expanded library semantic target}
\begin{ReviewVerbatim}
∀ {α : Type u_1} [inst : DecidableEq α] (q : ℕ) (C : EconCSLib.FiniteChoice.ChoiceRule α) (X : Finset α),
  (C X).card = min q X.card
\end{ReviewVerbatim}

\paragraph{Exact Lean library declaration}
\begin{ReviewVerbatim}
def QAcceptant (q : ℕ) (C : ChoiceRule α) : Prop :=
  ∀ X, (C X).card = min q X.card
\end{ReviewVerbatim}

\paragraph{Recorded source-to-library screening}
\textbf{Verdict:} \texttt{not recorded}
\reviewmatch{Matches source input}
\noindent\textbf{Reviewer annotation}\par
\reviewerbox
\clearpage
\hypertarget{library-prerequisite-88c2ddc1c549c716}{}
\subsection*{\small\ttfamily\raggedright EconCSLib.\allowbreak{}FiniteChoice.\allowbreak{}StrictTotalOrder}
\textbf{Source connection:} no explicit byte-pinned paper source connection is registered.
\paragraph{Lean-expanded library semantic target}
\begin{ReviewVerbatim}
∀ {α : Type u_1} (r : α → α → Prop),
  (∀ (x : α), ¬r x x) ∧ (∀ {x y z : α}, r x y → r y z → r x z) ∧ ∀ {x y : α}, x ≠ y → r x y ∨ r y x
\end{ReviewVerbatim}

\paragraph{Exact Lean library declaration}
\begin{ReviewVerbatim}
def StrictTotalOrder (r : α → α → Prop) : Prop :=
  (∀ x, ¬ r x x) ∧
    (∀ {x y z}, r x y → r y z → r x z) ∧
      (∀ {x y}, x ≠ y → r x y ∨ r y x)
\end{ReviewVerbatim}

\paragraph{Recorded source-to-library screening}
\textbf{Verdict:} \texttt{not recorded}
\reviewmatch{Matches source input}
\noindent\textbf{Reviewer annotation}\par
\reviewerbox
\clearpage
\hypertarget{library-prerequisite-286eac89d49c4d81}{}
\subsection*{\small\ttfamily\raggedright EconCSLib.\allowbreak{}FiniteChoice.\allowbreak{}QRepresentative}
\textbf{Source connection:} no explicit byte-pinned paper source connection is registered.
\paragraph{Lean-expanded library semantic target}
\begin{ReviewVerbatim}
∀ {α : Type u_1} [inst : DecidableEq α] (q : ℕ) (C : EconCSLib.FiniteChoice.ChoiceRule α),
  ∃ r,
    EconCSLib.FiniteChoice.StrictTotalOrder r ∧
      EconCSLib.FiniteChoice.QAcceptant q C ∧ ∀ {X : Finset α} {x y : α}, x ∈ C X → y ∈ X → y ∉ C X → r x y
\end{ReviewVerbatim}

\paragraph{Exact Lean library declaration}
\begin{ReviewVerbatim}
def QRepresentative (q : ℕ) (C : ChoiceRule α) : Prop :=
  ∃ r : α → α → Prop,
    StrictTotalOrder r ∧ QAcceptant q C ∧
      ∀ {X x y}, x ∈ C X → y ∈ X → y ∉ C X → r x y
\end{ReviewVerbatim}

\paragraph{Recorded source-to-library screening}
\textbf{Verdict:} \texttt{not recorded}
\reviewmatch{Matches source input}
\noindent\textbf{Reviewer annotation}\par
\reviewerbox
\clearpage
\hypertarget{library-prerequisite-542d5eef4a694a7d}{}
\subsection*{\small\ttfamily\raggedright EconCSLib.\allowbreak{}FiniteChoice.\allowbreak{}Substitutable}
\textbf{Source connection:} no explicit byte-pinned paper source connection is registered.
\paragraph{Lean-expanded library semantic target}
\begin{ReviewVerbatim}
∀ {α : Type u_1} [inst : DecidableEq α] (C : EconCSLib.FiniteChoice.ChoiceRule α) {X₁ X₂ : Finset α},
  X₁ ⊆ X₂ → X₁ ∩ C X₂ ⊆ C X₁
\end{ReviewVerbatim}

\paragraph{Exact Lean library declaration}
\begin{ReviewVerbatim}
def Substitutable (C : ChoiceRule α) : Prop :=
  ∀ {X₁ X₂}, X₁ ⊆ X₂ → X₁ ∩ C X₂ ⊆ C X₁
\end{ReviewVerbatim}

\paragraph{Recorded source-to-library screening}
\textbf{Verdict:} \texttt{not recorded}
\reviewmatch{Matches source input}
\noindent\textbf{Reviewer annotation}\par
\reviewerbox
\clearpage
\hypertarget{library-prerequisite-bb61a90095edae4b}{}
\subsection*{\small\ttfamily\raggedright EconCSLib.\allowbreak{}FiniteChoice.\allowbreak{}TightlyDUnstable}
\textbf{Source connection:} no explicit byte-pinned paper source connection is registered.
\paragraph{Lean-expanded library semantic target}
\begin{ReviewVerbatim}
∀ {α : Type u_1} [inst : DecidableEq α] (d : ℕ) (C : EconCSLib.FiniteChoice.ChoiceRule α),
  EconCSLib.FiniteChoice.DUnstable d C ∧ ∀ k < d, ¬EconCSLib.FiniteChoice.DUnstable k C
\end{ReviewVerbatim}

\paragraph{Exact Lean library declaration}
\begin{ReviewVerbatim}
def TightlyDUnstable (d : ℕ) (C : ChoiceRule α) : Prop :=
  DUnstable d C ∧ ∀ k, k < d → ¬ DUnstable k C
\end{ReviewVerbatim}

\paragraph{Recorded source-to-library screening}
\textbf{Verdict:} \texttt{not recorded}
\reviewmatch{Matches source input}
\noindent\textbf{Reviewer annotation}\par
\reviewerbox
\clearpage
\hypertarget{library-prerequisite-aa69285c37fbf372}{}
\subsection*{\small\ttfamily\raggedright EconCSLib.\allowbreak{}FiniteChoice.\allowbreak{}VariabilityAtMost}
\textbf{Source connection:} no explicit byte-pinned paper source connection is registered.
\paragraph{Lean-expanded library semantic target}
\begin{ReviewVerbatim}
∀ {α : Type u_1} [inst : DecidableEq α] [inst_1 : Fintype α] (m : ℕ) (C : EconCSLib.FiniteChoice.ChoiceRule α)
  (X : Finset α), (EconCSLib.FiniteChoice.borderlineSet C X).card ≤ m
\end{ReviewVerbatim}

\paragraph{Exact Lean library declaration}
\begin{ReviewVerbatim}
def VariabilityAtMost [Fintype α] (m : ℕ) (C : ChoiceRule α) : Prop :=
  ∀ X, (borderlineSet C X).card ≤ m
\end{ReviewVerbatim}

\paragraph{Recorded source-to-library screening}
\textbf{Verdict:} \texttt{not recorded}
\reviewmatch{Matches source input}
\noindent\textbf{Reviewer annotation}\par
\reviewerbox
\clearpage
\hypertarget{library-prerequisite-4ccdc8565ec23c71}{}
\subsection*{\small\ttfamily\raggedright EconCSLib.\allowbreak{}FiniteChoice.\allowbreak{}VariabilityExactly}
\textbf{Source connection:} no explicit byte-pinned paper source connection is registered.
\paragraph{Lean-expanded library semantic target}
\begin{ReviewVerbatim}
∀ {α : Type u_1} [inst : DecidableEq α] [inst_1 : Fintype α] (m : ℕ) (C : EconCSLib.FiniteChoice.ChoiceRule α),
  EconCSLib.FiniteChoice.VariabilityAtMost m C ∧ ∃ X, (EconCSLib.FiniteChoice.borderlineSet C X).card = m
\end{ReviewVerbatim}

\paragraph{Exact Lean library declaration}
\begin{ReviewVerbatim}
def VariabilityExactly [Fintype α] (m : ℕ) (C : ChoiceRule α) : Prop :=
  VariabilityAtMost m C ∧ ∃ X, (borderlineSet C X).card = m
\end{ReviewVerbatim}

\paragraph{Recorded source-to-library screening}
\textbf{Verdict:} \texttt{not recorded}
\reviewmatch{Matches source input}
\noindent\textbf{Reviewer annotation}\par
\reviewerbox
\clearpage
\hypertarget{library-prerequisite-e55d0830fae724cb}{}
\subsection*{\small\ttfamily\raggedright EconCSLib.\allowbreak{}FiniteChoice.\allowbreak{}ZeroUnstable}
\textbf{Source connection:} no explicit byte-pinned paper source connection is registered.
\paragraph{Lean-expanded library semantic target}
\begin{ReviewVerbatim}
∀ {α : Type u_1} [inst : DecidableEq α] (C : EconCSLib.FiniteChoice.ChoiceRule α) {X₁ X₂ : Finset α},
  X₁ ⊆ X₂ → EconCSLib.FiniteChoice.choiceDistance C X₁ X₂ = 0
\end{ReviewVerbatim}

\paragraph{Exact Lean library declaration}
\begin{ReviewVerbatim}
def ZeroUnstable (C : ChoiceRule α) : Prop :=
  ∀ {X₁ X₂}, X₁ ⊆ X₂ → choiceDistance C X₁ X₂ = 0
\end{ReviewVerbatim}

\paragraph{Recorded source-to-library screening}
\textbf{Verdict:} \texttt{not recorded}
\reviewmatch{Matches source input}
\noindent\textbf{Reviewer annotation}\par
\reviewerbox
\clearpage
\hypertarget{library-prerequisite-4a3989697f894599}{}
\subsection*{\small\ttfamily\raggedright EconCSLib.\allowbreak{}FiniteChoice.\allowbreak{}linearTopQChoice}
\textbf{Source connection:} no explicit byte-pinned paper source connection is registered.
\paragraph{Lean-expanded library semantic target}
\begin{ReviewVerbatim}
{α : Type u_1} →
  [inst : DecidableEq α] →
    [inst_1 : LinearOrder α] →
      (q : ℕ) →
        (a : Finset α) →
          if h : q ≤ a.card then Finset.image (fun i => ↑((a.orderIsoOfFin ⋯) ⟨↑i, ⋯⟩)) Finset.univ else a
\end{ReviewVerbatim}

\paragraph{Exact Lean library declaration}
\begin{ReviewVerbatim}
noncomputable def linearTopQChoice (q : ℕ) : ChoiceRule α :=
  fun X =>
    if h : q ≤ X.card then
      Finset.univ.image
        (fun i : Fin q =>
          ((X.orderIsoOfFin (by rfl)
            ⟨i.1, lt_of_lt_of_le i.2 h⟩ : X) : α))
    else
      X
\end{ReviewVerbatim}

\paragraph{Recorded source-to-library screening}
\textbf{Verdict:} \texttt{not recorded}
\reviewmatch{Matches source input}
\noindent\textbf{Reviewer annotation}\par
\reviewerbox
\clearpage
\hypertarget{library-prerequisite-3767c56c4f2cc69a}{}
\subsection*{\small\ttfamily\raggedright EconCSLib.\allowbreak{}FiniteChoice.\allowbreak{}sequentialComposition}
\textbf{Source connection:} no explicit byte-pinned paper source connection is registered.
\paragraph{Lean-expanded library semantic target}
\begin{ReviewVerbatim}
{α : Type u_1} →
  [inst : DecidableEq α] →
    (a : List (EconCSLib.FiniteChoice.ChoiceRule α)) →
      (a_1 : Finset α) →
        List.brecOn a
          (fun x f =>
            (match (motive :=
                (x : List (EconCSLib.FiniteChoice.ChoiceRule α)) → List.below x → EconCSLib.FiniteChoice.ChoiceRule α)
                x with
              | [] => fun x x_1 => ∅
              | C :: Cs => fun x X =>
                let chosen := C X;
                chosen ∪ x.1 (X \ chosen))
              f)
          a_1
\end{ReviewVerbatim}

\paragraph{Exact Lean library declaration}
\begin{ReviewVerbatim}
def sequentialComposition : List (ChoiceRule α) → ChoiceRule α
  | [] => fun _ => ∅
  | C :: Cs =>
      fun X =>
        let chosen := C X
        chosen ∪ sequentialComposition Cs (X \ chosen)

/-- Borderline admits displaced by some single added applicant. -/
def borderlineSet [Fintype α] (C : ChoiceRule α) (X : Finset α) : Finset α :=
  Finset.univ.biUnion fun x => C X \ C (insert x X)
\end{ReviewVerbatim}

\paragraph{Recorded source-to-library screening}
\textbf{Verdict:} \texttt{not recorded}
\reviewmatch{Matches source input}
\noindent\textbf{Reviewer annotation}\par
\reviewerbox

\clearpage
\hypertarget{source-claim-03de8677e0f95903}{}
\hypertarget{source-section-f10b5f68e4acf3dc}{}
\section*{Source claims}
\subsection*{Review row 1}
\textbf{Source locator:} {\footnotesize\raggedright cited publication, lines 162-\allowbreak{}170\par}
\paragraph{Verbatim source input}
\begin{ReviewVerbatim}
Training a machine learning model means selecting a $f$ from a set of possible models, i.e., a hypothesis class.
Generally, this is done by finding a model which achieves low error on available data (a train set).
It is typical to use continuous optimization algorithms to first train a continuous \emph{score} function $s:\mathcal X\to [0,1]$ by minimizing the error.
Then the classification model is defined as
\begin{align}\label{eq:ml-independent}
    f(x) = \mathbf 1\{s(x)\geq t\}
\end{align} based on a fixed threshold value $t$, often $t=0.5$.
Notice that models of the form~\eqref{eq:ml-independent} predict outcomes independently for each applicant.
At test time inference, they do not consider the overall applicant pool -- they implicitly model the training time pool, through its effect on the training labels.
\end{ReviewVerbatim}


\paragraph{Expanded PaperInterface specification}
\begin{ReviewVerbatim}
∀ {α : Type u_1} [inst : DecidableEq α] (score : α → ℝ) (threshold : ℝ) (X : Finset α) (x : α),
  x ∈ DGD26AdmissionsPredictability.paperFixedThresholdChoice score threshold X ↔
    x ∈ X ∧ DGD26AdmissionsPredictability.paper_definition_fixed_threshold_predictor score threshold X x
\end{ReviewVerbatim}

\paragraph{Lean proof endpoint (built separately)}\begin{ReviewVerbatim}
DGD26AdmissionsPredictability.PaperInterface.paper_fixed_threshold_formula_statement
\end{ReviewVerbatim}

\paragraph{Recorded source-to-Spec screening}
\textbf{Verdict:} \texttt{not recorded}
\\
\textbf{Reason:} not recorded
\reviewmatch{Matches source input}
\noindent\textbf{Reviewer annotation}\par
\reviewerbox
\clearpage
\hypertarget{source-claim-685dcba90250bf7c}{}
\section*{Review row 2}
\textbf{Source locator:} {\footnotesize\raggedright cited publication, lines 172-\allowbreak{}181\par}
\paragraph{Verbatim source input}
\begin{ReviewVerbatim}
\subsection{Choice Functions}

Many admissions processes have capacity constraints (either perceived or actual). This fact gives rise to what we call \textit{cohort-dependence}: an applicant's acceptance decision depends on not only their own characteristics, but also the number and characteristics of other applicants.

\textit{Choice functions} formalize admission decision processes and allow modeling such dependence.
A choice function describes the behavior of ``choosing'' elements from a set of available options, and is commonly used in economic models. Formally, a choice function over the universe $\calX$ is a function $\choice: 2^\calX \to 2^\calX$ that maps all subsets $X$ of $\calX$ to the accepted set $\choice(X)$, where $\choice(X) \subseteq X$. As our motivation focuses on real applicant datasets, we consider finite subsets.

Recalling admissions data, the binary admissions label $y$ for student $x$ is 1 if $x \in \choice(X)$ and 0 otherwise.
To aid in our discussion, we will slightly abuse notation to relate pairs $(x_i,y_i)$ to the choice function $\mathcal C(X)$ where $x_i\in X$.
Define $\mathcal C(X)_i$ as 1 if $x_i\in \mathcal C(X)$ and 0 otherwise; this allows us to write the label $y_i=\mathcal C(X)_i$.
\end{ReviewVerbatim}


\paragraph{Expanded PaperInterface specification}
\begin{ReviewVerbatim}
∀ {α : Type u_1} [inst : DecidableEq α] (C : DGD26AdmissionsPredictability.PaperChoiceRule α),
  DGD26AdmissionsPredictability.paper_feasible C ↔ ∀ (X : Finset α), C X ⊆ X
\end{ReviewVerbatim}

\paragraph{Lean proof endpoint (built separately)}\begin{ReviewVerbatim}
DGD26AdmissionsPredictability.PaperInterface.paper_choice_function_model_definition_statement
\end{ReviewVerbatim}

\paragraph{Recorded source-to-Spec screening}
\textbf{Verdict:} \texttt{not recorded}
\\
\textbf{Reason:} not recorded
\reviewmatch{Matches source input}
\noindent\textbf{Reviewer annotation}\par
\reviewerbox
\clearpage
\hypertarget{source-claim-d396bf72ee1812a2}{}
\section*{Review row 3}
\textbf{Source locator:} {\footnotesize\raggedright cited publication, lines 179-\allowbreak{}181\par}
\paragraph{Verbatim source input}
\begin{ReviewVerbatim}
Recalling admissions data, the binary admissions label $y$ for student $x$ is 1 if $x \in \choice(X)$ and 0 otherwise.
To aid in our discussion, we will slightly abuse notation to relate pairs $(x_i,y_i)$ to the choice function $\mathcal C(X)$ where $x_i\in X$.
Define $\mathcal C(X)_i$ as 1 if $x_i\in \mathcal C(X)$ and 0 otherwise; this allows us to write the label $y_i=\mathcal C(X)_i$.
\end{ReviewVerbatim}


\paragraph{Expanded PaperInterface specification}
\begin{ReviewVerbatim}
∀ {α : Type u_1} [inst : DecidableEq α] (C : DGD26AdmissionsPredictability.PaperChoiceRule α) (X : Finset α) (x : α),
  DGD26AdmissionsPredictability.paperChoiceLabel C X x = if x ∈ C X then 1 else 0
\end{ReviewVerbatim}

\paragraph{Lean proof endpoint (built separately)}\begin{ReviewVerbatim}
DGD26AdmissionsPredictability.PaperInterface.paper_definition_choice_label
\end{ReviewVerbatim}

\paragraph{Recorded source-to-Spec screening}
\textbf{Verdict:} \texttt{not recorded}
\\
\textbf{Reason:} not recorded
\reviewmatch{Matches source input}
\noindent\textbf{Reviewer annotation}\par
\reviewerbox
\clearpage
\hypertarget{source-claim-51b92e4a229b02a5}{}
\section*{Review row 4}
\textbf{Source locator:} {\footnotesize\raggedright cited publication, lines 186-\allowbreak{}189\par}
\paragraph{Verbatim source input}
\begin{ReviewVerbatim}
\begin{restatable}[$q$-Acceptance]{definition}{qacceptance}
\label{def:q-acceptance}
Choice function $\choice$ is $q$-acceptant if $|\choice(X)| = \min\{q, |X|\}$ for all input sets $X$.
\end{restatable}

[Next verbatim source excerpt]

\qacceptance*
\textit{Intuition}: Capacity of $q$.

\substitutability*
\end{ReviewVerbatim}


\paragraph{Expanded PaperInterface specification}
\begin{ReviewVerbatim}
∀ {α : Type u_1} [inst : DecidableEq α] (q : ℕ) (C : DGD26AdmissionsPredictability.PaperChoiceRule α),
  DGD26AdmissionsPredictability.paper_definition_q_acceptance q C ↔ ∀ (X : Finset α), (C X).card = min q X.card
\end{ReviewVerbatim}

\paragraph{Lean proof endpoint (built separately)}\begin{ReviewVerbatim}
DGD26AdmissionsPredictability.PaperInterface.paper_q_acceptance_definition_statement
\end{ReviewVerbatim}

\paragraph{Recorded source-to-Spec screening}
\textbf{Verdict:} \texttt{not recorded}
\\
\textbf{Reason:} not recorded
\reviewmatch{Matches source input}
\noindent\textbf{Reviewer annotation}\par
\reviewerbox
\clearpage
\hypertarget{source-claim-9a3be13c89f94ffe}{}
\section*{Review row 5}
\textbf{Source locator:} {\footnotesize\raggedright cited publication, lines 193-\allowbreak{}195\par}
\paragraph{Verbatim source input}
\begin{ReviewVerbatim}
\begin{definition}[Total order] \label{def:total_order}
    A \emph{total order} is a transitive, asymmetric, and complete binary relation $\succ$ over elements of $\calX$; that is, $x_1 \succ x_2$ implies $x_2 \not\succ x_1$, $x_1 \succ x_2$ and $x_2 \succ x_3$ implies $x_1 \succ x_3$, and all $x_1, x_2$ are comparable.
\end{definition}
\end{ReviewVerbatim}

\paragraph{Approved corrected review target}
The archival source above is not asserted equivalent to this replacement.
\begin{ReviewVerbatim}
A strict total order is asymmetric and transitive, and compares every distinct pair of applicants.
\end{ReviewVerbatim}

\textbf{Archival source anchor:} {\footnotesize\raggedright cited publication, lines 193-\allowbreak{}195\par}
\paragraph{Recorded basis}
DGD source clarification:\allowbreak{} strict-\allowbreak{}order completeness ranges over distinct applicant pairs.\allowbreak{}
\paragraph{Expanded PaperInterface specification}
\begin{ReviewVerbatim}
∀ {α : Type u_1} (r : α → α → Prop),
  DGD26AdmissionsPredictability.paper_definition_total_order r ↔
    (∀ (x : α), ¬r x x) ∧ (∀ {x y z : α}, r x y → r y z → r x z) ∧ ∀ {x y : α}, x ≠ y → r x y ∨ r y x
\end{ReviewVerbatim}

\paragraph{Lean proof endpoint (built separately)}\begin{ReviewVerbatim}
DGD26AdmissionsPredictability.PaperInterface.paper_total_order_definition_statement
\end{ReviewVerbatim}

\paragraph{Recorded source-to-Spec screening}
\textbf{Verdict:} \texttt{not recorded}
\\
\textbf{Reason:} not recorded
\reviewmatch{Matches approved corrected target}
\noindent\textbf{Reviewer annotation}\par
\reviewerbox
\clearpage
\hypertarget{source-claim-a21144374ad8e28c}{}
\section*{Review row 6}
\textbf{Source locator:} {\footnotesize\raggedright cited publication, lines 213-\allowbreak{}218\par}
\paragraph{Verbatim source input}
\begin{ReviewVerbatim}
\begin{restatable}[Choice Distance]{definition}{distance} \label{def:distance}
For a choice function $\choice$ and sets $X_1 \subseteq X_2$, we define the choice distance as:
\begin{equation}
r_\choice(X_1, X_2) := |X_1 \cap \choice(X_2) \setminus \choice(X_1)| + |\choice(X_1) \setminus \choice(X_2)|
\end{equation}
\end{restatable}
\end{ReviewVerbatim}


\paragraph{Expanded PaperInterface specification}
\begin{ReviewVerbatim}
∀ {α : Type u_1} [inst : DecidableEq α] (C : DGD26AdmissionsPredictability.PaperChoiceRule α) (X₁ X₂ : Finset α),
  EconCSLib.FiniteChoice.choiceDistance C X₁ X₂ = ((X₁ ∩ C X₂) \ C X₁).card + (C X₁ \ C X₂).card
\end{ReviewVerbatim}

\paragraph{Lean proof endpoint (built separately)}\begin{ReviewVerbatim}
DGD26AdmissionsPredictability.PaperInterface.paper_definition_choice_distance
\end{ReviewVerbatim}

\paragraph{Recorded source-to-Spec screening}
\textbf{Verdict:} \texttt{not recorded}
\\
\textbf{Reason:} not recorded
\reviewmatch{Matches source input}
\noindent\textbf{Reviewer annotation}\par
\reviewerbox
\clearpage
\hypertarget{source-claim-398a4d87fcced6ba}{}
\section*{Review row 7}
\textbf{Source locator:} {\footnotesize\raggedright cited publication, lines 222-\allowbreak{}225\par}
\paragraph{Verbatim source input}
\begin{ReviewVerbatim}
\begin{restatable}[$d$-Instability]{definition}{defstability}
\label{def:d-instability}
Choice function $\choice$ is $d$-unstable if for all $X$, and $X' = X \cup \{x'\}$, we have $r_\choice(X, X') \leq d$ for any nonnegative integer $d$. $\choice$ is called tightly $d$-unstable if it is $d$-unstable but not $(d-1)$-unstable.
\end{restatable}
\end{ReviewVerbatim}


\paragraph{Expanded PaperInterface specification}
\begin{ReviewVerbatim}
∀ {α : Type u_1} [inst : DecidableEq α] (d : ℕ) (C : DGD26AdmissionsPredictability.PaperChoiceRule α),
  DGD26AdmissionsPredictability.paper_definition_d_instability d C ↔
    ∀ (X : Finset α), ∀ x ∉ X, EconCSLib.FiniteChoice.choiceDistance C X (insert x X) ≤ d
\end{ReviewVerbatim}

\paragraph{Lean proof endpoint (built separately)}\begin{ReviewVerbatim}
DGD26AdmissionsPredictability.PaperInterface.paper_d_instability_definition_statement
\end{ReviewVerbatim}

\paragraph{Recorded source-to-Spec screening}
\textbf{Verdict:} \texttt{not recorded}
\\
\textbf{Reason:} not recorded
\reviewmatch{Matches source input}
\noindent\textbf{Reviewer annotation}\par
\reviewerbox
\clearpage
\hypertarget{source-claim-62d00d89c6cc151a}{}
\section*{Review row 8}
\textbf{Source locator:} {\footnotesize\raggedright cited publication, lines 222-\allowbreak{}225\par}
\paragraph{Verbatim source input}
\begin{ReviewVerbatim}
\begin{restatable}[$d$-Instability]{definition}{defstability}
\label{def:d-instability}
Choice function $\choice$ is $d$-unstable if for all $X$, and $X' = X \cup \{x'\}$, we have $r_\choice(X, X') \leq d$ for any nonnegative integer $d$. $\choice$ is called tightly $d$-unstable if it is $d$-unstable but not $(d-1)$-unstable.
\end{restatable}
\end{ReviewVerbatim}


\paragraph{Expanded PaperInterface specification}
\begin{ReviewVerbatim}
∀ {α : Type u_1} [inst : DecidableEq α] (d : ℕ) (C : DGD26AdmissionsPredictability.PaperChoiceRule α),
  DGD26AdmissionsPredictability.paper_definition_tight_d_instability d C ↔
    DGD26AdmissionsPredictability.paper_definition_d_instability d C ∧
      ∀ k < d, ¬DGD26AdmissionsPredictability.paper_definition_d_instability k C
\end{ReviewVerbatim}

\paragraph{Lean proof endpoint (built separately)}\begin{ReviewVerbatim}
DGD26AdmissionsPredictability.PaperInterface.paper_tight_d_instability_definition_statement
\end{ReviewVerbatim}

\paragraph{Recorded source-to-Spec screening}
\textbf{Verdict:} \texttt{not recorded}
\\
\textbf{Reason:} not recorded
\reviewmatch{Matches source input}
\noindent\textbf{Reviewer annotation}\par
\reviewerbox
\clearpage
\hypertarget{source-claim-be4f0b5bb7e1dd6b}{}
\section*{Review row 9}
\textbf{Source locator:} {\footnotesize\raggedright cited publication, lines 235-\allowbreak{}243\par}
\paragraph{Verbatim source input}
\begin{ReviewVerbatim}
We formalize\footnote{For a more general definition of variability (beyond 1-unstable and $q$-acceptant), see \Cref{sec:variability-appendix}.} this intuition with the following definition.

\begin{restatable}[Variability]{definition}{defvariability}
\label{def:variability}
A $1$-unstable, $q$-acceptant choice function $\choice$ has variability $m$ where
\begin{equation}
m \coloneqq \max_{X\subset \calX} \left|\bigcup\limits_{x'\in\calX} \choice(X) \setminus \choice(X \cup \{x'\})\right|
\end{equation}
\end{restatable}
\end{ReviewVerbatim}


\paragraph{Expanded PaperInterface specification}
\begin{ReviewVerbatim}
∀ {α : Type u_1} [inst : DecidableEq α] [inst_1 : Fintype α] (m : ℕ)
  (C : DGD26AdmissionsPredictability.PaperChoiceRule α),
  DGD26AdmissionsPredictability.paper_definition_variability_exactly m C ↔
    (∀ (X : Finset α), (DGD26AdmissionsPredictability.paper_definition_borderline_set C X).card ≤ m) ∧
      ∃ X, (DGD26AdmissionsPredictability.paper_definition_borderline_set C X).card = m
\end{ReviewVerbatim}

\paragraph{Lean proof endpoint (built separately)}\begin{ReviewVerbatim}
DGD26AdmissionsPredictability.PaperInterface.paper_variability_exactly_definition_statement
\end{ReviewVerbatim}

\paragraph{Recorded source-to-Spec screening}
\textbf{Verdict:} \texttt{not recorded}
\\
\textbf{Reason:} not recorded
\reviewmatch{Matches source input}
\noindent\textbf{Reviewer annotation}\par
\reviewerbox
\clearpage
\hypertarget{source-claim-12cc974d07b19ac2}{}
\section*{Review row 10}
\textbf{Source locator:} {\footnotesize\raggedright cited publication, lines 248-\allowbreak{}251\par}
\paragraph{Verbatim source input}
\begin{ReviewVerbatim}
We first relate the concepts of instability and variability to the practice of machine learning.
We focus on the representation capabilities of ML models: when is it possible for an ML model to faithfully represent an admissions process?
Formally, we say that an ML model can represent a choice function $\mathcal C$ if there exists a function $f$ such that for all applicant sets $X$, $f(x_i) = \mathcal C(X)_i$.
{Note that by focusing on representation, our theory is distribution agnostic, and applies to settings where applicant distributions are affected considerations such as strategic behavior and matching algorithms.}
\end{ReviewVerbatim}


\paragraph{Expanded PaperInterface specification}
\begin{ReviewVerbatim}
∀ {α : Type u_1} [inst : DecidableEq α] (predicts : DGD26AdmissionsPredictability.PaperPoolPredictor α)
  (C : DGD26AdmissionsPredictability.PaperChoiceRule α),
  DGD26AdmissionsPredictability.paper_definition_ml_representation predicts C ↔
    ∀ (X : Finset α), ∀ x ∈ X, predicts X x ↔ x ∈ C X
\end{ReviewVerbatim}

\paragraph{Lean proof endpoint (built separately)}\begin{ReviewVerbatim}
DGD26AdmissionsPredictability.PaperInterface.paper_ml_representation_definition_statement
\end{ReviewVerbatim}

\paragraph{Recorded source-to-Spec screening}
\textbf{Verdict:} \texttt{not recorded}
\\
\textbf{Reason:} not recorded
\reviewmatch{Matches source input}
\noindent\textbf{Reviewer annotation}\par
\reviewerbox
\clearpage
\hypertarget{source-claim-1e70e813c2db1d8e}{}
\section*{Review row 11}
\textbf{Source locator:} {\footnotesize\raggedright cited publication, lines 253-\allowbreak{}259\par}
\paragraph{Verbatim source input}
\begin{ReviewVerbatim}
Before stating this result, consider a straightforward way to adapt a model~\eqref{eq:ml-independent} to the capacity-constrained setting.
Given that the ML model is based on continuous scores, we can adjust the decisions to account for a known capacity constraint if we have available the set of all test-time applications.
In particular, we can rank the applicants by their scores and then select the top $q$.
This process results in decisions determined by a cohort dependent threshold $t_q(X)$ which is the score of the $q$th ranked applicant:
\begin{align}\label{eq:ml-rank}
    f(x; X) = \mathbf 1\{s(x) \geq t_q(X)\}.
\end{align}
\end{ReviewVerbatim}


\paragraph{Expanded PaperInterface specification}
\begin{ReviewVerbatim}
∀ {α : Type u_1} [inst : DecidableEq α] (q : ℕ) (score : α → ℝ) (hinjective : Function.Injective score),
  0 < q →
    ∀ (X : Finset α),
      ∃ threshold,
        ∀ x ∈ X,
          DGD26AdmissionsPredictability.paper_definition_rank_threshold_predictor q score hinjective X x ↔
            threshold ≤ score x
\end{ReviewVerbatim}

\paragraph{Lean proof endpoint (built separately)}\begin{ReviewVerbatim}
DGD26AdmissionsPredictability.PaperInterface.paper_rank_threshold_formula_statement
\end{ReviewVerbatim}

\paragraph{Recorded source-to-Spec screening}
\textbf{Verdict:} \texttt{not recorded}
\\
\textbf{Reason:} not recorded
\reviewmatch{Matches source input}
\noindent\textbf{Reviewer annotation}\par
\reviewerbox
\clearpage
\hypertarget{source-claim-d17644eda5ebe53f}{}
\section*{Review row 12}
\textbf{Source locator:} {\footnotesize\raggedright cited publication, lines 273-\allowbreak{}276\par}
\paragraph{Verbatim source input}
\begin{ReviewVerbatim}
\begin{restatable}[Substitutability]{definition}{substitutability}
\label{def:substitutability}
Choice function $\choice$ is substitutable if $X_1 \subseteq X_2$ implies $X_1 \cap \choice(X_2) \subseteq \choice(X_1)$.
\end{restatable}

[Next verbatim source excerpt]


\substitutability*
\textit{Intuition}: Students selected from a larger pool remain selected when the pool shrinks.
\end{ReviewVerbatim}


\paragraph{Expanded PaperInterface specification}
\begin{ReviewVerbatim}
∀ {α : Type u_1} [inst : DecidableEq α] (C : DGD26AdmissionsPredictability.PaperChoiceRule α),
  DGD26AdmissionsPredictability.paper_definition_substitutability C ↔ ∀ {X₁ X₂ : Finset α}, X₁ ⊆ X₂ → X₁ ∩ C X₂ ⊆ C X₁
\end{ReviewVerbatim}

\paragraph{Lean proof endpoint (built separately)}\begin{ReviewVerbatim}
DGD26AdmissionsPredictability.PaperInterface.paper_substitutability_definition_statement
\end{ReviewVerbatim}

\paragraph{Recorded source-to-Spec screening}
\textbf{Verdict:} \texttt{not recorded}
\\
\textbf{Reason:} not recorded
\reviewmatch{Matches source input}
\noindent\textbf{Reviewer annotation}\par
\reviewerbox
\clearpage
\hypertarget{source-claim-69b86836fa7d4d06}{}
\section*{Review row 13}
\textbf{Source locator:} {\footnotesize\raggedright cited publication, lines 307-\allowbreak{}311\par}
\paragraph{Verbatim source input}
\begin{ReviewVerbatim}
\begin{restatable}[Sequential Composition]{definition}{seqcomposition} \label{def:seq-composition}
A choice function $\choice$ is a \textit{sequential composition} of functions $\choice_n, \ldots, \choice_1$ if:
$$\choice(X) = \choice_1(X) \cup \choice_{2}(X_{2}) \cup \cdots \cup \choice_n(X_n)$$
where $X_{i+1}= X_{i} \setminus \choice_{i}(X_{i})$ and $X_1 = X$.
\end{restatable}

[Next verbatim source excerpt]

\subsubsection{Sequential Composition} \label{sec:seq-proofs}
A useful class of choice functions are choice functions that can be broken down or defined as a sequential composition of functions $\choice_n, \choice_{n-1}, ..., \choice_1$; that is,
$\choice(X) = \choice_n(X) \cup \choice_{n-1}(X_{n-1}) \cup ... \cup \choice_1(X_1)$ where $X_i = X_{i+1} \setminus\choice_{i+1}(X_{i+1})$ and $X_n = X$. Sequential composition is particularly useful because it preserves a variety of useful properties. Sequences of $q_i$-acceptant functions are $q$-acceptant, because no duplicate elements are selected by multiple functions, and therefore preserves $q$-acceptance when each function in the sequence has capacity $q_i$ (i.e., $q = q_n + ... + q_1$.)
\seqcomposition*
\end{ReviewVerbatim}


\paragraph{Expanded PaperInterface specification}
\begin{ReviewVerbatim}
∀ {α : Type u_1} [inst : DecidableEq α] (Cs : List (DGD26AdmissionsPredictability.PaperChoiceRule α)),
  DGD26AdmissionsPredictability.paper_sequential_composition Cs =
    DGD26AdmissionsPredictability.sequentialCompositionSource Cs
\end{ReviewVerbatim}

\paragraph{Lean proof endpoint (built separately)}\begin{ReviewVerbatim}
DGD26AdmissionsPredictability.PaperInterface.paper_definition_sequential_composition
\end{ReviewVerbatim}

\paragraph{Recorded source-to-Spec screening}
\textbf{Verdict:} \texttt{not recorded}
\\
\textbf{Reason:} not recorded
\reviewmatch{Matches source input}
\noindent\textbf{Reviewer annotation}\par
\reviewerbox
\clearpage
\hypertarget{source-claim-1b6dc1033989d2bf}{}
\section*{Review row 14}
\textbf{Source locator:} {\footnotesize\raggedright cited publication, lines 484-\allowbreak{}487\par}
\paragraph{Verbatim source input}
\begin{ReviewVerbatim}
\begin{restatable}[Monotonicity]{definition}{monotonicity}
\label{def:monotonicity}
Choice function $\choice$ is \textit{monotonic} if $X_1 \subseteq X_2$ implies $\choice(X_1) \subseteq \choice(X_2)$.
\end{restatable}
\end{ReviewVerbatim}


\paragraph{Expanded PaperInterface specification}
\begin{ReviewVerbatim}
∀ {α : Type u_1} [inst : DecidableEq α] (C : DGD26AdmissionsPredictability.PaperChoiceRule α),
  DGD26AdmissionsPredictability.paper_definition_monotonicity C ↔ ∀ {X₁ X₂ : Finset α}, X₁ ⊆ X₂ → C X₁ ⊆ C X₂
\end{ReviewVerbatim}

\paragraph{Lean proof endpoint (built separately)}\begin{ReviewVerbatim}
DGD26AdmissionsPredictability.PaperInterface.paper_monotonicity_definition_statement
\end{ReviewVerbatim}

\paragraph{Recorded source-to-Spec screening}
\textbf{Verdict:} \texttt{not recorded}
\\
\textbf{Reason:} not recorded
\reviewmatch{Matches source input}
\noindent\textbf{Reviewer annotation}\par
\reviewerbox
\clearpage
\hypertarget{source-claim-899ae3aec5210e2c}{}
\section*{Review row 15}
\textbf{Source locator:} {\footnotesize\raggedright cited publication, lines 492-\allowbreak{}495\par}
\paragraph{Verbatim source input}
\begin{ReviewVerbatim}
\begin{restatable}[Consistency]{definition}{consistency}
\label{def:consistency}
Choice function $\choice$ is \textit{consistent} if $\choice(X_2) \subseteq X_1 \subseteq X_2$ implies $\choice(X_2) = \choice(X_1)$.
\end{restatable}
\end{ReviewVerbatim}


\paragraph{Expanded PaperInterface specification}
\begin{ReviewVerbatim}
∀ {α : Type u_1} [inst : DecidableEq α] (C : DGD26AdmissionsPredictability.PaperChoiceRule α),
  DGD26AdmissionsPredictability.paper_definition_consistency C ↔ ∀ {X₁ X₂ : Finset α}, C X₂ ⊆ X₁ → X₁ ⊆ X₂ → C X₂ = C X₁
\end{ReviewVerbatim}

\paragraph{Lean proof endpoint (built separately)}\begin{ReviewVerbatim}
DGD26AdmissionsPredictability.PaperInterface.paper_consistency_definition_statement
\end{ReviewVerbatim}

\paragraph{Recorded source-to-Spec screening}
\textbf{Verdict:} \texttt{not recorded}
\\
\textbf{Reason:} not recorded
\reviewmatch{Matches source input}
\noindent\textbf{Reviewer annotation}\par
\reviewerbox
\clearpage
\hypertarget{source-claim-ca1d4873227a50e3}{}
\section*{Review row 16}
\textbf{Source locator:} {\footnotesize\raggedright cited publication, lines 500-\allowbreak{}507\par}
\paragraph{Verbatim source input}
\begin{ReviewVerbatim}
\begin{restatable}[$q$-representativeness]{definition}{qrepresentativeness}
\label{def:q-representative}
$q$-acceptant choice function $\choice$ is $q$-representative if there exists a strict total ordering $\succ$ over all $x \in \calX$ and $\choice(X)$ is the top $q$ elements of $X$ for all $X \subseteq \calX$.
\end{restatable}

\textit{Intuition}: Students are selected according to an exact ranking.

\textbf{Note:} We refer to this property in the main body as being ``characterized by a total order'', or a ``queue'', but this language --- and definition as a property in its own right --- is seen in, e.g., \citet{echenique2015control}.

[Next verbatim source excerpt]

Total orders are frequently used in modeling choice mechanisms; they are often called \textit{preference orderings} or \textit{priority orders} in economic models.
Intuitively, total orders induce a ranking over candidates and could result from, e.g., ordering applicants based on their GPA.
We will say that a $q$-acceptant choice function $\mathcal C$ is \textit{characterized by a total order} if there exists a total order $\succ$ such that for any input $X$, the accepted applicants are the top-$q$ elements of $X$ ordered by $\succ$.
In other words,  $x_1 \succ x_2$ for all $x_1 \in \choice(X), x_2 \not\in\choice(X)$.
We will refer to the process of sorting elements according to a total order $\succ$ and selecting the top $q$ as a \textit{queue}.
\end{ReviewVerbatim}


\paragraph{Expanded PaperInterface specification}
\begin{ReviewVerbatim}
∀ {α : Type u_1} [inst : DecidableEq α] (q : ℕ) (C : DGD26AdmissionsPredictability.PaperChoiceRule α),
  DGD26AdmissionsPredictability.paper_definition_q_representativeness q C ↔
    ∃ r,
      DGD26AdmissionsPredictability.paper_definition_total_order r ∧
        DGD26AdmissionsPredictability.paper_definition_q_acceptance q C ∧
          ∀ {X : Finset α} {x y : α}, x ∈ C X → y ∈ X → y ∉ C X → r x y
\end{ReviewVerbatim}

\paragraph{Lean proof endpoint (built separately)}\begin{ReviewVerbatim}
DGD26AdmissionsPredictability.PaperInterface.paper_q_representativeness_definition_statement
\end{ReviewVerbatim}

\paragraph{Recorded source-to-Spec screening}
\textbf{Verdict:} \texttt{not recorded}
\\
\textbf{Reason:} not recorded
\reviewmatch{Matches source input}
\noindent\textbf{Reviewer annotation}\par
\reviewerbox
\clearpage
\hypertarget{source-claim-3a97b71dd915eed1}{}
\section*{Review row 17}
\textbf{Source locator:} {\footnotesize\raggedright cited publication, lines 654-\allowbreak{}657\par}
\paragraph{Verbatim source input}
\begin{ReviewVerbatim}
\begin{restatable}[$0$-Instability]{definition}{zerostability}
\label{def:zero-instability}
Choice function $\choice$ is 0-unstable if and only if $r_\choice(X_1, X_2) = 0$ for all $X_1 \subseteq X_2$.
\end{restatable}
\end{ReviewVerbatim}


\paragraph{Expanded PaperInterface specification}
\begin{ReviewVerbatim}
∀ {α : Type u_1} [inst : DecidableEq α] (C : DGD26AdmissionsPredictability.PaperChoiceRule α),
  DGD26AdmissionsPredictability.paper_definition_zero_instability C ↔
    ∀ {X₁ X₂ : Finset α}, X₁ ⊆ X₂ → EconCSLib.FiniteChoice.choiceDistance C X₁ X₂ = 0
\end{ReviewVerbatim}

\paragraph{Lean proof endpoint (built separately)}\begin{ReviewVerbatim}
DGD26AdmissionsPredictability.PaperInterface.paper_zero_instability_definition_statement
\end{ReviewVerbatim}

\paragraph{Recorded source-to-Spec screening}
\textbf{Verdict:} \texttt{not recorded}
\\
\textbf{Reason:} not recorded
\reviewmatch{Matches source input}
\noindent\textbf{Reviewer annotation}\par
\reviewerbox
\clearpage
\hypertarget{source-claim-e633cfedc7bfb493}{}
\section*{Review row 18}
\textbf{Source locator:} {\footnotesize\raggedright cited publication, lines 671-\allowbreak{}679\par}
\paragraph{Verbatim source input}
\begin{ReviewVerbatim}
\begin{restatable}{definition}{Independence}
\label{def:independence}
$\choice$ is independent if for
 for all $x$,
either $x \in \choice(X)$
or $x \not\in \choice(X)$ for all $X$ where $x \in X$.
\end{restatable}

In other words, any element is either always accepted, or always rejected, when available, irrespective of the other elements available in a set $S$.
\end{ReviewVerbatim}


\paragraph{Expanded PaperInterface specification}
\begin{ReviewVerbatim}
∀ {α : Type u_1} [inst : DecidableEq α] (C : DGD26AdmissionsPredictability.PaperChoiceRule α),
  DGD26AdmissionsPredictability.paper_definition_independence C ↔
    ∀ (x : α), (∀ (X : Finset α), x ∈ X → x ∈ C X) ∨ ∀ (X : Finset α), x ∈ X → x ∉ C X
\end{ReviewVerbatim}

\paragraph{Lean proof endpoint (built separately)}\begin{ReviewVerbatim}
DGD26AdmissionsPredictability.PaperInterface.paper_independence_definition_statement
\end{ReviewVerbatim}

\paragraph{Recorded source-to-Spec screening}
\textbf{Verdict:} \texttt{not recorded}
\\
\textbf{Reason:} not recorded
\reviewmatch{Matches source input}
\noindent\textbf{Reviewer annotation}\par
\reviewerbox
\clearpage
\hypertarget{source-claim-8e553b3e026d261e}{}
\section*{Review row 19}
\textbf{Source locator:} {\footnotesize\raggedright cited publication, lines 826-\allowbreak{}830\par}
\paragraph{Verbatim source input}
\begin{ReviewVerbatim}
As per \Cref{def:variability}, when appending a new element to the input set $X$, different applicants may become rejected: let the set of ``borderline admits'' $V_\choice(X)$ be
\begin{equation}
V_\choice(X) = \bigcup\limits_{x'\in\calX} \choice(X) \setminus \choice(X \cup \{x'\})
\end{equation}
which we shall call the ``borderline set'' for short. Notice that this is bounded by $q$, as $|\choice(X)| = q$.
\end{ReviewVerbatim}


\paragraph{Expanded PaperInterface specification}
\begin{ReviewVerbatim}
∀ {α : Type u_1} [inst : DecidableEq α] [inst_1 : Fintype α] (C : DGD26AdmissionsPredictability.PaperChoiceRule α)
  (X : Finset α),
  DGD26AdmissionsPredictability.paper_borderline_set C X = Finset.univ.biUnion fun x => C X \ C (insert x X)
\end{ReviewVerbatim}

\paragraph{Lean proof endpoint (built separately)}\begin{ReviewVerbatim}
DGD26AdmissionsPredictability.PaperInterface.paper_definition_borderline_set
\end{ReviewVerbatim}

\paragraph{Recorded source-to-Spec screening}
\textbf{Verdict:} \texttt{not recorded}
\\
\textbf{Reason:} not recorded
\reviewmatch{Matches source input}
\noindent\textbf{Reviewer annotation}\par
\reviewerbox
\clearpage
\hypertarget{source-claim-1c2136cbe341a833}{}
\section*{Review row 20}
\textbf{Source locator:} {\footnotesize\raggedright cited publication, lines 832-\allowbreak{}836\par}
\paragraph{Verbatim source input}
\begin{ReviewVerbatim}
Conversely, when removing an applicant from $X$, let the set of applicants that might be newly accepted be
\begin{equation}
A_\choice(X) = \bigcup_{x^*\in X} \choice(X \setminus \{x^*\}) \setminus \choice(X)
\end{equation}
the set of ``waitlisted rejections''; likewise, we shall call this the ``waitlisted set''.
\end{ReviewVerbatim}


\paragraph{Expanded PaperInterface specification}
\begin{ReviewVerbatim}
∀ {α : Type u_1} [inst : DecidableEq α] [inst_1 : Fintype α] (C : DGD26AdmissionsPredictability.PaperChoiceRule α)
  (X : Finset α), DGD26AdmissionsPredictability.paper_waitlisted_set C X = X.biUnion fun x => C (X.erase x) \ C X
\end{ReviewVerbatim}

\paragraph{Lean proof endpoint (built separately)}\begin{ReviewVerbatim}
DGD26AdmissionsPredictability.PaperInterface.paper_definition_waitlisted_set
\end{ReviewVerbatim}

\paragraph{Recorded source-to-Spec screening}
\textbf{Verdict:} \texttt{not recorded}
\\
\textbf{Reason:} not recorded
\reviewmatch{Matches source input}
\noindent\textbf{Reviewer annotation}\par
\reviewerbox
\clearpage
\hypertarget{source-claim-d2f6a286dc057af6}{}
\section*{Review row 21}
\textbf{Source locator:} {\footnotesize\raggedright cited publication, lines 838-\allowbreak{}842\par}
\paragraph{Verbatim source input}
\begin{ReviewVerbatim}
\begin{definition}[Variability, General]
\label{def:variability-general}
A choice function $\choice$ has \textit{variability} $m$ where:
$$m := \max\left\{\max_{X\subset \calX} \left|V_\choice(X)\right|, \max_{X\subset \calX} \left|A_\choice(X)\right| \right\}$$
\end{definition}
\end{ReviewVerbatim}


\paragraph{Expanded PaperInterface specification}
\begin{ReviewVerbatim}
∀ {α : Type u_1} [inst : DecidableEq α] [inst_1 : Fintype α] (m : ℕ)
  (C : DGD26AdmissionsPredictability.PaperChoiceRule α),
  DGD26AdmissionsPredictability.paper_definition_general_variability_exactly m C ↔
    DGD26AdmissionsPredictability.paper_definition_general_variability_at_most m C ∧
      ((∃ X, (DGD26AdmissionsPredictability.paper_definition_borderline_set C X).card = m) ∨
        ∃ X, (DGD26AdmissionsPredictability.paper_definition_waitlisted_set C X).card = m)
\end{ReviewVerbatim}

\paragraph{Lean proof endpoint (built separately)}\begin{ReviewVerbatim}
DGD26AdmissionsPredictability.PaperInterface.paper_general_variability_exactly_definition_statement
\end{ReviewVerbatim}

\paragraph{Recorded source-to-Spec screening}
\textbf{Verdict:} \texttt{not recorded}
\\
\textbf{Reason:} not recorded
\reviewmatch{Matches source input}
\noindent\textbf{Reviewer annotation}\par
\reviewerbox
\clearpage
\hypertarget{source-claim-a009d7013e95842d}{}
\section*{Review row 22}
\textbf{Source locator:} {\footnotesize\raggedright cited publication, lines 1073-\allowbreak{}1084\par}
\paragraph{Verbatim source input}
\begin{ReviewVerbatim}
\subsubsection{Linear Assignment Problems} \label{sec:lap-proofs}

We highlight \textit{linear assignment problems} as a generalization of sequential composition.
The linear sum assignment problem can be represented as a bipartite graph matching.
One one side of the graph are applicants, and on the other are available admissions slots.
Edges connect these nodes according to how well suited each applicant is for each slot.
The admissions decision is made by selecting edges which have maximal sum, subject to the constraint that each applicant and each slot is selected only once.

Consider all edges connecting to slot $s$; then for each element  $x_j \in X$ there is a real number weight $w(x_j, s)$. Then sorting $x \in \calX$ by $w(x,s)$ creates a total ordering, which we call $\succ_s$. Let slot $s'$ be \textit{equivalent} to
to slot $s$ if and only if $\succ_{s'} = \succ_{s}$.
We will assume for simplicity that weights are defined such that ties do not occur.
Let assignment $A$ be the set of pairs $\{(x_{s_1}, s_1), ..., (x_{s_q}, s_q)\}$.
\end{ReviewVerbatim}


\paragraph{Expanded PaperInterface specification}
\begin{ReviewVerbatim}
∀ {α : Type u_1} [inst : DecidableEq α] {σ : Type u_2} [inst_1 : DecidableEq σ] [inst_2 : Fintype σ] (X : Finset α)
  (w : α → σ → ℝ) (A : DGD26AdmissionsPredictability.LAP.Assignment α σ),
  DGD26AdmissionsPredictability.lapModel X w A ↔
    DGD26AdmissionsPredictability.paper_definition_lap_assignment_feasible X A ∧
      DGD26AdmissionsPredictability.paper_definition_lap_capacity_filling X A ∧
        DGD26AdmissionsPredictability.paper_definition_lap_objective_optimal X w A
\end{ReviewVerbatim}

\paragraph{Lean proof endpoint (built separately)}\begin{ReviewVerbatim}
DGD26AdmissionsPredictability.PaperInterface.paper_lap_model_definition_statement
\end{ReviewVerbatim}

\paragraph{Recorded source-to-Spec screening}
\textbf{Verdict:} \texttt{not recorded}
\\
\textbf{Reason:} not recorded
\reviewmatch{Matches source input}
\noindent\textbf{Reviewer annotation}\par
\reviewerbox
\clearpage
\hypertarget{source-claim-f902fe2aa64a28cb}{}
\section*{Review row 23}
\textbf{Source locator:} {\footnotesize\raggedright cited publication, lines 1081-\allowbreak{}1083\par}
\paragraph{Verbatim source input}
\begin{ReviewVerbatim}
Consider all edges connecting to slot $s$; then for each element  $x_j \in X$ there is a real number weight $w(x_j, s)$. Then sorting $x \in \calX$ by $w(x,s)$ creates a total ordering, which we call $\succ_s$. Let slot $s'$ be \textit{equivalent} to
to slot $s$ if and only if $\succ_{s'} = \succ_{s}$.
We will assume for simplicity that weights are defined such that ties do not occur.
\end{ReviewVerbatim}


\paragraph{Expanded PaperInterface specification}
\begin{ReviewVerbatim}
∀ {α : Type u_1} {σ : Type u_2} [inst : DecidableEq σ] [inst_1 : Fintype σ] (w : α → σ → ℝ) (s : σ) (x y : α),
  DGD26AdmissionsPredictability.paper_definition_lap_slot_below w s x y ↔ w x s < w y s
\end{ReviewVerbatim}

\paragraph{Lean proof endpoint (built separately)}\begin{ReviewVerbatim}
DGD26AdmissionsPredictability.PaperInterface.paper_lap_slot_below_definition_statement
\end{ReviewVerbatim}

\paragraph{Recorded source-to-Spec screening}
\textbf{Verdict:} \texttt{not recorded}
\\
\textbf{Reason:} not recorded
\reviewmatch{Matches source input}
\noindent\textbf{Reviewer annotation}\par
\reviewerbox
\clearpage
\hypertarget{source-claim-ac931f819ad35103}{}
\section*{Review row 24}
\textbf{Source locator:} {\footnotesize\raggedright cited publication, lines 1081-\allowbreak{}1083\par}
\paragraph{Verbatim source input}
\begin{ReviewVerbatim}
Consider all edges connecting to slot $s$; then for each element  $x_j \in X$ there is a real number weight $w(x_j, s)$. Then sorting $x \in \calX$ by $w(x,s)$ creates a total ordering, which we call $\succ_s$. Let slot $s'$ be \textit{equivalent} to
to slot $s$ if and only if $\succ_{s'} = \succ_{s}$.
We will assume for simplicity that weights are defined such that ties do not occur.
\end{ReviewVerbatim}


\paragraph{Expanded PaperInterface specification}
\begin{ReviewVerbatim}
∀ {α : Type u_1} {σ : Type u_2} [DecidableEq σ] [Fintype σ] (w : α → σ → ℝ) (s t : σ),
  DGD26AdmissionsPredictability.LAP.Assignment.SameSlotOrder w s t ↔ ∀ (x y : α), w x s < w y s ↔ w x t < w y t
\end{ReviewVerbatim}

\paragraph{Lean proof endpoint (built separately)}\begin{ReviewVerbatim}
DGD26AdmissionsPredictability.PaperInterface.paper_lap_same_slot_order_definition_statement
\end{ReviewVerbatim}

\paragraph{Recorded source-to-Spec screening}
\textbf{Verdict:} \texttt{not recorded}
\\
\textbf{Reason:} not recorded
\reviewmatch{Matches source input}
\noindent\textbf{Reviewer annotation}\par
\reviewerbox
\clearpage
\hypertarget{source-claim-be70d394f2bee79c}{}
\section*{Review row 25}
\textbf{Source locator:} {\footnotesize\raggedright cited publication, lines 1084-\allowbreak{}1089\par}
\paragraph{Verbatim source input}
\begin{ReviewVerbatim}
Let assignment $A$ be the set of pairs $\{(x_{s_1}, s_1), ..., (x_{s_q}, s_q)\}$.

\begin{restatable}{lemma}{LAP Ordering}
\label{lem:lap-ordering}
Let $x_{s}$ denote the element of $\choice(X)$ assigned to slot $s$. Then for all $x \in X$, $x \in \choice(X)$ if $x \succ_s x_s$, and $x_s \succ_s x$ if $x \not\in \choice(X)$.
\end{restatable}
\end{ReviewVerbatim}


\paragraph{Expanded PaperInterface specification}
\begin{ReviewVerbatim}
∀ {α : Type u_1} [inst : DecidableEq α] {σ : Type u_2} [inst_1 : DecidableEq σ] [inst_2 : Fintype σ]
  (select : Finset α → DGD26AdmissionsPredictability.LAP.Assignment α σ) (X : Finset α) (x : α),
  x ∈ DGD26AdmissionsPredictability.paper_definition_lap_choice_rule select X ↔
    DGD26AdmissionsPredictability.paper_definition_lap_assigned (select X) x
\end{ReviewVerbatim}

\paragraph{Lean proof endpoint (built separately)}\begin{ReviewVerbatim}
DGD26AdmissionsPredictability.PaperInterface.paper_lap_assignment_choice_formula_statement
\end{ReviewVerbatim}

\paragraph{Recorded source-to-Spec screening}
\textbf{Verdict:} \texttt{not recorded}
\\
\textbf{Reason:} not recorded
\reviewmatch{Matches source input}
\noindent\textbf{Reviewer annotation}\par
\reviewerbox
\clearpage
\hypertarget{source-claim-b05952c99a84b6af}{}
\section*{Review row 26}
\textbf{Source locator:} {\footnotesize\raggedright cited publication, lines 704-\allowbreak{}707\par}
\paragraph{Verbatim source input}
\begin{ReviewVerbatim}
\begin{restatable}{corollary}{Only Independent Rules are 0-unstable}
\label{thm:independent-zero-unstable}
$\choice$ is $0$-unstable if and only if it is independent, and therefore no $q$-acceptant $\choice$ can be $0$-unstable.
\end{restatable}
\end{ReviewVerbatim}


\paragraph{Expanded PaperInterface specification}
\begin{ReviewVerbatim}
∀ {α : Type u_1} [inst : DecidableEq α] (C : DGD26AdmissionsPredictability.PaperChoiceRule α),
  DGD26AdmissionsPredictability.paper_choice_function_feasible C →
    (DGD26AdmissionsPredictability.paper_definition_independence C ↔
      DGD26AdmissionsPredictability.paper_definition_zero_instability C)
\end{ReviewVerbatim}

\paragraph{Lean proof endpoint (built separately)}\begin{ReviewVerbatim}
DGD26AdmissionsPredictability.PaperInterface.paper_independent_zero_unstable_statement
\end{ReviewVerbatim}

\paragraph{Recorded source-to-Spec screening}
\textbf{Verdict:} \texttt{not recorded}
\\
\textbf{Reason:} not recorded
\reviewmatch{Matches source input}
\noindent\textbf{Reviewer annotation}\par
\reviewerbox
\clearpage
\hypertarget{source-claim-d920f9b7847b2baa}{}
\section*{Review row 27}
\textbf{Source locator:} {\footnotesize\raggedright cited publication, lines 704-\allowbreak{}707\par}
\paragraph{Verbatim source input}
\begin{ReviewVerbatim}
\begin{restatable}{corollary}{Only Independent Rules are 0-unstable}
\label{thm:independent-zero-unstable}
$\choice$ is $0$-unstable if and only if it is independent, and therefore no $q$-acceptant $\choice$ can be $0$-unstable.
\end{restatable}
\end{ReviewVerbatim}

\paragraph{Approved corrected review target}
The archival source above is not asserted equivalent to this replacement.
\begin{ReviewVerbatim}
For a feasible q-acceptant choice rule C on a universe with an input pool strictly larger than positive capacity q, C is 0-unstable if and only if it is independent, and C is not 0-unstable.
\end{ReviewVerbatim}

\textbf{Archival source anchor:} {\footnotesize\raggedright cited publication, lines 704-\allowbreak{}707\par}
\paragraph{Recorded basis}
DGD source clarification:\allowbreak{} the no-\allowbreak{}zero conclusion is read on the positive-\allowbreak{}capacity, nontrivial applicant domain.\allowbreak{}
\paragraph{Expanded PaperInterface specification}
\begin{ReviewVerbatim}
∀ {α : Type u_1} [inst : DecidableEq α] {q : ℕ} {C : DGD26AdmissionsPredictability.PaperChoiceRule α},
  DGD26AdmissionsPredictability.paper_choice_function_feasible C →
    DGD26AdmissionsPredictability.paper_definition_q_acceptance q C →
      0 < q →
        ∀ {U : Finset α},
          q < U.card →
            (DGD26AdmissionsPredictability.paper_definition_zero_instability C ↔
                DGD26AdmissionsPredictability.paper_definition_independence C) ∧
              ¬DGD26AdmissionsPredictability.paper_definition_zero_instability C
\end{ReviewVerbatim}

\paragraph{Lean proof endpoint (built separately)}\begin{ReviewVerbatim}
DGD26AdmissionsPredictability.PaperInterface.paper_independent_zero_unstable_corollary_statement
\end{ReviewVerbatim}

\paragraph{Recorded source-to-Spec screening}
\textbf{Verdict:} \texttt{not recorded}
\\
\textbf{Reason:} not recorded
\reviewmatch{Matches approved corrected target}
\noindent\textbf{Reviewer annotation}\par
\reviewerbox
\clearpage
\hypertarget{source-claim-18a641b7c4365d27}{}
\section*{Review row 28}
\textbf{Source locator:} {\footnotesize\raggedright cited publication, lines 796-\allowbreak{}801\par}
\paragraph{Verbatim source input}
\begin{ReviewVerbatim}
\begin{restatable}{corollary}{No Consistent Tightly Even Instability}
\label{cor:even-inconsistency}
There is no $q$-acceptant and consistent $\choice$ that is tightly $dk$-unstable for even values of $d$.
\end{restatable}

This follows directly from \Cref{thm:even-inconsistency}.
\end{ReviewVerbatim}

\paragraph{Approved corrected review target}
The archival source above is not asserted equivalent to this replacement.
\begin{ReviewVerbatim}
For every positive integer k, no q-acceptant consistent choice rule is tightly 2k-unstable.
\end{ReviewVerbatim}

\textbf{Archival source anchor:} {\footnotesize\raggedright cited publication, lines 796-\allowbreak{}801\par}
\paragraph{Recorded basis}
DGD source clarification:\allowbreak{} tightly-\allowbreak{}even instability is parameterized by d=2k for positive k.\allowbreak{}
\paragraph{Expanded PaperInterface specification}
\begin{ReviewVerbatim}
∀ {α : Type u_1} [inst : DecidableEq α] {q k : ℕ} {C : DGD26AdmissionsPredictability.PaperChoiceRule α},
  0 < k →
    DGD26AdmissionsPredictability.paper_choice_function_feasible C →
      DGD26AdmissionsPredictability.paper_definition_q_acceptance q C →
        DGD26AdmissionsPredictability.paper_definition_consistency C →
          ¬DGD26AdmissionsPredictability.paper_definition_tight_d_instability (2 * k) C
\end{ReviewVerbatim}

\paragraph{Lean proof endpoint (built separately)}\begin{ReviewVerbatim}
DGD26AdmissionsPredictability.PaperInterface.paper_no_consistent_positive_tightly_even_statement
\end{ReviewVerbatim}

\paragraph{Recorded source-to-Spec screening}
\textbf{Verdict:} \texttt{not recorded}
\\
\textbf{Reason:} not recorded
\reviewmatch{Matches approved corrected target}
\noindent\textbf{Reviewer annotation}\par
\reviewerbox
\clearpage
\hypertarget{source-claim-1daa198b9d0a1bc3}{}
\section*{Review row 29}
\textbf{Source locator:} {\footnotesize\raggedright cited publication, lines 661-\allowbreak{}667\par}
\paragraph{Verbatim source input}
\begin{ReviewVerbatim}
\begin{restatable}{corollary}{Zero-Distance}
\label{thm:sub-mon-zero}
Choice function $\choice$ is 0-unstable if and only if it is both substitutable and monotonic.
\end{restatable}


This immediately follows from the two terms comprising $r_\choice$ mapping onto substitutability and monotonicity.
\end{ReviewVerbatim}


\paragraph{Expanded PaperInterface specification}
\begin{ReviewVerbatim}
∀ {α : Type u_1} [inst : DecidableEq α] (C : DGD26AdmissionsPredictability.PaperChoiceRule α),
  DGD26AdmissionsPredictability.paper_definition_zero_instability C ↔
    DGD26AdmissionsPredictability.paper_definition_substitutability C ∧
      DGD26AdmissionsPredictability.paper_definition_monotonicity C
\end{ReviewVerbatim}

\paragraph{Lean proof endpoint (built separately)}\begin{ReviewVerbatim}
DGD26AdmissionsPredictability.PaperInterface.paper_zero_distance_statement
\end{ReviewVerbatim}

\paragraph{Recorded source-to-Spec screening}
\textbf{Verdict:} \texttt{not recorded}
\\
\textbf{Reason:} not recorded
\reviewmatch{Matches source input}
\noindent\textbf{Reviewer annotation}\par
\reviewerbox
\clearpage
\hypertarget{source-claim-760d01f2c31e3672}{}
\section*{Review row 30}
\textbf{Source locator:} {\footnotesize\raggedright cited publication, lines 746-\allowbreak{}765\par}
\paragraph{Verbatim source input}
\begin{ReviewVerbatim}
\begin{restatable}[Calculating Instability]{lemma}{stablecalc}
\label{lem:unstable-calc}
For all $q$-acceptant $\choice$, let $X' = X \cup \{x'\}$ for $|X| \geq q$ and $n = |\choice(X) \setminus \choice(X')|$.  Then
$r_\choice(X, X') = 2n-\mathbf{1}_{x'\in\choice(X')}$.
\end{restatable}

\begin{proof}
\textbf{Informal}
If input sets differ by one element, but the selected sets differ by $n$ elements, then it must be that $n$ previously accepted elements were rejected and there must be $n$ new accepted elements, at least $n-1$ must have been previously rejected, because the size of the accepted set is fixed.

\textbf{Formal}
By $q$-acceptance, $|\choice(X)| = |\choice(X')| = q$. Then $|\choice(X) \setminus \choice(X')| = |\choice(X') \setminus \choice(X)| = q-|\choice(X) \cap \choice(X')|$.
Second, note that $\choice(X') \setminus \choice(X) = (\choice(X') \cap X) \setminus \choice(X)$ if $x' \not\in \choice(X')$, and otherwise $\choice(X') \setminus \choice(X) = \{x'\} \cup \left((\choice(X') \cap X) \setminus \choice(X)\right)$.
Putting the previous statements together, we get two cases:
\begin{itemize}
    \item $x' \not\in \choice(X')$: $|\choice(X) \setminus \choice(X')| = |\choice(X') \setminus \choice(X)| = |X \cap \choice(X') \setminus \choice(X)|$
    \item $x' \in \choice(X')$: $|\choice(X) \setminus \choice(X')| = |\choice(X') \setminus \choice(X)| = |X \cap \choice(X') \setminus \choice(X)| + 1$
\end{itemize}
Then $r_\choice(X, X') = |X \cap \choice(X') \setminus \choice(X)| + |\choice(X) \setminus \choice(X')| = 2|X \cap \choice(X') \setminus \choice(X)| + \mathbf{1}_{x'\in\choice(X')} = 2|\choice(X) \setminus \choice(X')| - \mathbf{1}_{x'\in\choice(X')}$
\end{proof}
\end{ReviewVerbatim}


\paragraph{Expanded PaperInterface specification}
\begin{ReviewVerbatim}
∀ {α : Type u_1} [inst : DecidableEq α] {q : ℕ} {C : DGD26AdmissionsPredictability.PaperChoiceRule α},
  DGD26AdmissionsPredictability.paper_choice_function_feasible C →
    DGD26AdmissionsPredictability.paper_definition_q_acceptance q C →
      ∀ {X : Finset α} {x : α},
        q ≤ X.card →
          x ∉ X →
            DGD26AdmissionsPredictability.paper_definition_choice_distance C X (insert x X) =
              if x ∈ C (insert x X) then 2 * (C X \ C (insert x X)).card - 1 else 2 * (C X \ C (insert x X)).card
\end{ReviewVerbatim}

\paragraph{Lean proof endpoint (built separately)}\begin{ReviewVerbatim}
DGD26AdmissionsPredictability.PaperInterface.paper_calculating_instability_statement
\end{ReviewVerbatim}

\paragraph{Recorded source-to-Spec screening}
\textbf{Verdict:} \texttt{not recorded}
\\
\textbf{Reason:} not recorded
\reviewmatch{Matches source input}
\noindent\textbf{Reviewer annotation}\par
\reviewerbox
\clearpage
\hypertarget{source-claim-72ba750e0e18812e}{}
\section*{Review row 31}
\textbf{Source locator:} {\footnotesize\raggedright cited publication, lines 885-\allowbreak{}896\par}
\paragraph{Verbatim source input}
\begin{ReviewVerbatim}
\begin{restatable}[Consistency of Removable Sets]{lemma}{consistentremove}
\label{lem:consistent-remove}
Let $\choice$ be 1-unstable and q-acceptant with variability $m$. For any $X_1, X_2$, if $\choice(X_1) = \choice(X_2)$, then $V_\choice(X_1) = V_\choice(X_1)$.
\end{restatable}

\begin{proof}
Recall from \Cref{thm:sub-accept-consistent} that $\choice$ must be consistent and so $\choice(\choice(X)) = \choice(X)$ for any $X$.

Consider any $x'$ associated with an element $x^*$ of $V_\choice(X)$, such that $\choice(X) \setminus \choice(X') = \{x^*\}$ (and even more precisely, by $q$-acceptant substitutability, that $\choice(X') = \{x'\} \cup \choice(X) \setminus \{x^*\}$ exactly). Note that $\choice(X') \subseteq \choice(X) \cup \{x'\} \subseteq X'$, and therefore $\choice(\choice(X) \cup \{x'\}) = \choice(X')$ by consistency. This applies to all possible $x'$ (including when $x' \in X$), and therefore, $V_\choice(X) = V_\choice(\choice(X))$.

As this set equality is exact, apply this logic to some $X_2$ and see $V_\choice(X_1) = V_\choice(\choice(X_1)) = V_\choice(\choice(X_2)) = V_\choice(X_2)$.
\end{proof}
\end{ReviewVerbatim}

\paragraph{Approved corrected review target}
The archival source above is not asserted equivalent to this replacement.
\begin{ReviewVerbatim}
For a feasible q-acceptant substitutable choice rule C, if C(X1)=C(X2), then the borderline sets V_C(X1) and V_C(X2) are equal.
\end{ReviewVerbatim}

\textbf{Archival source anchor:} {\footnotesize\raggedright cited publication, lines 885-\allowbreak{}896\par}
\paragraph{Recorded basis}
DGD source clarification:\allowbreak{} the removable-\allowbreak{}set equality compares the two pools X1 and X2.\allowbreak{}
\paragraph{Expanded PaperInterface specification}
\begin{ReviewVerbatim}
∀ {α : Type u_1} [inst : DecidableEq α] [inst_1 : Fintype α] {q : ℕ}
  {C : DGD26AdmissionsPredictability.PaperChoiceRule α},
  DGD26AdmissionsPredictability.paper_choice_function_feasible C →
    DGD26AdmissionsPredictability.paper_definition_q_acceptance q C →
      DGD26AdmissionsPredictability.paper_definition_substitutability C →
        ∀ {X₁ X₂ : Finset α},
          C X₁ = C X₂ →
            DGD26AdmissionsPredictability.paper_definition_borderline_set C X₁ =
              DGD26AdmissionsPredictability.paper_definition_borderline_set C X₂
\end{ReviewVerbatim}

\paragraph{Lean proof endpoint (built separately)}\begin{ReviewVerbatim}
DGD26AdmissionsPredictability.PaperInterface.paper_corrected_consistency_of_removable_sets_statement
\end{ReviewVerbatim}

\paragraph{Recorded source-to-Spec screening}
\textbf{Verdict:} \texttt{not recorded}
\\
\textbf{Reason:} not recorded
\reviewmatch{Matches approved corrected target}
\noindent\textbf{Reviewer annotation}\par
\reviewerbox
\clearpage
\hypertarget{source-claim-7d8732899c3f4210}{}
\section*{Review row 32}
\textbf{Source locator:} {\footnotesize\raggedright cited publication, lines 1086-\allowbreak{}1095\par}
\paragraph{Verbatim source input}
\begin{ReviewVerbatim}
\begin{restatable}{lemma}{LAP Ordering}
\label{lem:lap-ordering}
Let $x_{s}$ denote the element of $\choice(X)$ assigned to slot $s$. Then for all $x \in X$, $x \in \choice(X)$ if $x \succ_s x_s$, and $x_s \succ_s x$ if $x \not\in \choice(X)$.
\end{restatable}

\textit{Intuition}: Linear assignment problems have a structure that resembles a combination of rankings, even though there is not a fixed sequential order.

\begin{proof}
This comes by the optimality of the assignment solution. For contradiction, say there is $x^* \succ_s x_s$ where $x^* \not\in \choice(X)$. Then $w(x^*, s) > w(x_s, s)$, and therefore $\{x^*, s\} \cup A \setminus \{x_s, s\}$ leads to a higher sum than $\choice(X)$, contradicting the fact that $A$ is be optimal.
\end{proof}
\end{ReviewVerbatim}


\paragraph{Expanded PaperInterface specification}
\begin{ReviewVerbatim}
∀ {α : Type u_1} [inst : DecidableEq α] {σ : Type u_2} [inst_1 : DecidableEq σ] [inst_2 : Fintype σ] {X : Finset α}
  {w : α → σ → ℝ} {A : DGD26AdmissionsPredictability.LAP.Assignment α σ},
  DGD26AdmissionsPredictability.paper_definition_lap_objective_optimal X w A →
    DGD26AdmissionsPredictability.paper_definition_lap_capacity_filling X A →
      DGD26AdmissionsPredictability.paper_definition_lap_assignment_feasible X A →
        ∀ {s : σ} {y x : α},
          A.matchSlot s = some y →
            x ∈ X →
              DGD26AdmissionsPredictability.paper_definition_lap_slot_below w s y x →
                DGD26AdmissionsPredictability.paper_definition_lap_assigned A x
\end{ReviewVerbatim}

\paragraph{Lean proof endpoint (built separately)}\begin{ReviewVerbatim}
DGD26AdmissionsPredictability.PaperInterface.paper_lap_strictly_higher_slot_applicant_assigned_statement
\end{ReviewVerbatim}

\paragraph{Recorded source-to-Spec screening}
\textbf{Verdict:} \texttt{not recorded}
\\
\textbf{Reason:} not recorded
\reviewmatch{Matches source input}
\noindent\textbf{Reviewer annotation}\par
\reviewerbox
\clearpage
\hypertarget{source-claim-5e3bd8769e715d3c}{}
\section*{Review row 33}
\textbf{Source locator:} {\footnotesize\raggedright cited publication, lines 1086-\allowbreak{}1095\par}
\paragraph{Verbatim source input}
\begin{ReviewVerbatim}
\begin{restatable}{lemma}{LAP Ordering}
\label{lem:lap-ordering}
Let $x_{s}$ denote the element of $\choice(X)$ assigned to slot $s$. Then for all $x \in X$, $x \in \choice(X)$ if $x \succ_s x_s$, and $x_s \succ_s x$ if $x \not\in \choice(X)$.
\end{restatable}

\textit{Intuition}: Linear assignment problems have a structure that resembles a combination of rankings, even though there is not a fixed sequential order.

\begin{proof}
This comes by the optimality of the assignment solution. For contradiction, say there is $x^* \succ_s x_s$ where $x^* \not\in \choice(X)$. Then $w(x^*, s) > w(x_s, s)$, and therefore $\{x^*, s\} \cup A \setminus \{x_s, s\}$ leads to a higher sum than $\choice(X)$, contradicting the fact that $A$ is be optimal.
\end{proof}
\end{ReviewVerbatim}


\paragraph{Expanded PaperInterface specification}
\begin{ReviewVerbatim}
∀ {α : Type u_1} [inst : DecidableEq α] {σ : Type u_2} [inst_1 : DecidableEq σ] [inst_2 : Fintype σ] {X : Finset α}
  {w : α → σ → ℝ} {A : DGD26AdmissionsPredictability.LAP.Assignment α σ},
  DGD26AdmissionsPredictability.paper_definition_lap_objective_optimal X w A →
    DGD26AdmissionsPredictability.paper_definition_lap_capacity_filling X A →
      DGD26AdmissionsPredictability.paper_definition_lap_assignment_feasible X A →
        ∀ {s : σ} {y x : α},
          A.matchSlot s = some y →
            DGD26AdmissionsPredictability.LAP.Assignment.Rejected X A x →
              DGD26AdmissionsPredictability.LAP.Assignment.SlotNoTies w s →
                DGD26AdmissionsPredictability.paper_definition_lap_slot_below w s x y
\end{ReviewVerbatim}

\paragraph{Lean proof endpoint (built separately)}\begin{ReviewVerbatim}
DGD26AdmissionsPredictability.PaperInterface.paper_lap_assigned_strictly_outranks_rejected_statement
\end{ReviewVerbatim}

\paragraph{Recorded source-to-Spec screening}
\textbf{Verdict:} \texttt{not recorded}
\\
\textbf{Reason:} not recorded
\reviewmatch{Matches source input}
\noindent\textbf{Reviewer annotation}\par
\reviewerbox
\clearpage
\hypertarget{source-claim-14ff35045a22028c}{}
\section*{Review row 34}
\textbf{Source locator:} {\footnotesize\raggedright cited publication, lines 564-\allowbreak{}573\par}
\paragraph{Verbatim source input}
\begin{ReviewVerbatim}
\begin{restatable}[Monotonicity Term]{lemma}{mono-dist}
\label{lem:monotonicity-term}
Choice function $\choice$ is monotonic if and only if $|\choice(X_1) \setminus \choice(X_2)| = 0$ for all $X_1 \subseteq X_2$.
\end{restatable}

\textit{Intuition}:
For a monotonic choice function, previously accepted candidates are never rejected after the addition of more applicants. Thus labels cannot flip from positive to negative, which is what the second term counts.

\begin{proof}
Use the same set properties as per \Cref{lem:substitutability-term},  $|\choice(X_1) \setminus \choice(X_2)| = 0$ if and only if $\choice(X_1) \subseteq \choice(X_2)$, which is the definition of monotonicity.
\end{ReviewVerbatim}


\paragraph{Expanded PaperInterface specification}
\begin{ReviewVerbatim}
∀ {α : Type u_1} [inst : DecidableEq α] (C : DGD26AdmissionsPredictability.PaperChoiceRule α),
  DGD26AdmissionsPredictability.paper_definition_monotonicity C ↔ ∀ {X₁ X₂ : Finset α}, X₁ ⊆ X₂ → (C X₁ \ C X₂).card = 0
\end{ReviewVerbatim}

\paragraph{Lean proof endpoint (built separately)}\begin{ReviewVerbatim}
DGD26AdmissionsPredictability.PaperInterface.paper_monotonicity_term_statement
\end{ReviewVerbatim}

\paragraph{Recorded source-to-Spec screening}
\textbf{Verdict:} \texttt{not recorded}
\\
\textbf{Reason:} not recorded
\reviewmatch{Matches source input}
\noindent\textbf{Reviewer annotation}\par
\reviewerbox
\clearpage
\hypertarget{source-claim-43c4c9b7821d369a}{}
\section*{Review row 35}
\textbf{Source locator:} {\footnotesize\raggedright cited publication, lines 517-\allowbreak{}529\par}
\paragraph{Verbatim source input}
\begin{ReviewVerbatim}
\begin{restatable}[Incompatibility of Monotonicity and $q$-Acceptance]{lemma}{incompatibility}
\label{thm:mono-accept-mutex}
For $|\calX| > q \geq 1$, no choice function can be both monotonic and $q$-acceptant.
\end{restatable}

\textit{Intuition}: %
For a choice function to be monotonic, it must never reject a previously accepted applicant as the applicant pool grows. This implies that the capacity must be unbounded, which contradicts the strict capacity limit imposed by $q$-acceptance.%

\begin{proof}
Consider two different sets each of $q$ elements. If $\choice$ is $q$-acceptant, every element must be chosen in each. Then consider their union. If $\choice$ is monotonic, every element in their union would also be accepted, but that would contradict the $q$-acceptant capacity constraint.

Consider two disjoint sets $X_1, X_2 \subseteq \calX$ with $|X_1| = |X_2| = q$. By $q$-acceptance, $\choice(X_1) = X_1$ and $\choice(X_2) = X_2$. By monotonicity, both $\choice(X_1)$ and $\choice(X_2)$ must be subsets of $\choice(X_1 \cup X_2)$. But then $|\choice(X_1 \cup X_2)| \geq |X_1| + |X_2| = 2q > q$, violating $q$-acceptance.
\end{proof}
\end{ReviewVerbatim}


\paragraph{Expanded PaperInterface specification}
\begin{ReviewVerbatim}
∀ {α : Type u_1} [inst : DecidableEq α] {q : ℕ} {C : DGD26AdmissionsPredictability.PaperChoiceRule α},
  DGD26AdmissionsPredictability.paper_choice_function_feasible C →
    DGD26AdmissionsPredictability.paper_definition_monotonicity C →
      DGD26AdmissionsPredictability.paper_definition_q_acceptance q C → 0 < q → ∀ {U : Finset α}, q < U.card → False
\end{ReviewVerbatim}

\paragraph{Lean proof endpoint (built separately)}\begin{ReviewVerbatim}
DGD26AdmissionsPredictability.PaperInterface.paper_monotonicity_q_acceptance_incompatible_statement
\end{ReviewVerbatim}

\paragraph{Recorded source-to-Spec screening}
\textbf{Verdict:} \texttt{not recorded}
\\
\textbf{Reason:} not recorded
\reviewmatch{Matches source input}
\noindent\textbf{Reviewer annotation}\par
\reviewerbox
\clearpage
\hypertarget{source-claim-f926138cb2114337}{}
\section*{Review row 36}
\textbf{Source locator:} {\footnotesize\raggedright cited publication, lines 898-\allowbreak{}905\par}
\paragraph{Verbatim source input}
\begin{ReviewVerbatim}
\begin{restatable}{lemma}{Ranking m}
\label{lem:a-v-equal}
Let $q$-acceptant, $1$-unstable $\choice$ have variability $1$. Then $V_\choice(X) = A_\choice(X \cup \{x'\})$ when $\choice(X \cup \{x'\}) \neq \choice(X)$.
\end{restatable}

\begin{proof}
If $\choice(X \cup \{x'\}) \neq \choice(X)$, then $\choice(X \cup \{x'\}) = \{x'\} \cup \choice(X) \setminus x^*$ where $V_\choice(X) = \{x^*\}$ (by consistency, 1-instability, and the definition of variability). Then $X' \setminus \{x'\}) = X$ and so $A_\choice(X') = \{x^*\}$, completing the proof.
\end{proof}
\end{ReviewVerbatim}


\paragraph{Expanded PaperInterface specification}
\begin{ReviewVerbatim}
∀ {α : Type u_1} [inst : DecidableEq α] [inst_1 : Fintype α] {q : ℕ}
  {C : DGD26AdmissionsPredictability.PaperChoiceRule α},
  DGD26AdmissionsPredictability.paper_choice_function_feasible C →
    DGD26AdmissionsPredictability.paper_definition_q_acceptance q C →
      DGD26AdmissionsPredictability.paper_definition_d_instability 1 C →
        DGD26AdmissionsPredictability.paper_definition_variability_at_most 1 C →
          ∀ {X : Finset α} {x : α},
            x ∉ X →
              C (insert x X) ≠ C X →
                DGD26AdmissionsPredictability.paper_definition_borderline_set C X =
                  DGD26AdmissionsPredictability.paper_definition_waitlisted_set C (insert x X)
\end{ReviewVerbatim}

\paragraph{Lean proof endpoint (built separately)}\begin{ReviewVerbatim}
DGD26AdmissionsPredictability.PaperInterface.paper_acceptant_one_instability_variability_borderline_eq_waitlisted_after_changing_insert_statement
\end{ReviewVerbatim}

\paragraph{Recorded source-to-Spec screening}
\textbf{Verdict:} \texttt{not recorded}
\\
\textbf{Reason:} not recorded
\reviewmatch{Matches source input}
\noindent\textbf{Reviewer annotation}\par
\reviewerbox
\clearpage
\hypertarget{source-claim-ef4d58145103cd0c}{}
\section*{Review row 37}
\textbf{Source locator:} {\footnotesize\raggedright cited publication, lines 552-\allowbreak{}562\par}
\paragraph{Verbatim source input}
\begin{ReviewVerbatim}
\begin{restatable}[Substitutability Term]{lemma}{sub-dist}
\label{lem:substitutability-term}
Choice function $\choice$ is substitutable if and only if $|X_1 \cap \choice(X_2) \setminus \choice(X_1)| = 0$ for all $X_1 \subseteq X_2$.
\end{restatable}

\textit{Intuition}:
For a substitutable choice function, previously rejected candidates are never accepted after the addition of more applicants. Thus labels cannot flip from negative to positive, which is what the first term counts.

\begin{proof}
For any sets $A, B$, $|A \setminus B| = 0$ if and only if $A \setminus B = \emptyset$, and furthermore $A \setminus B = \emptyset$ if and only if $A \subseteq B$. So $|X_1 \cap \choice(X_2) \setminus \choice(X_1)| = 0$ if and only if $X_1 \cap \choice(X_2) \subseteq \choice(X_1)$, which is the definition of substitutability.
\end{proof}
\end{ReviewVerbatim}


\paragraph{Expanded PaperInterface specification}
\begin{ReviewVerbatim}
∀ {α : Type u_1} [inst : DecidableEq α] (C : DGD26AdmissionsPredictability.PaperChoiceRule α),
  DGD26AdmissionsPredictability.paper_definition_substitutability C ↔
    ∀ {X₁ X₂ : Finset α}, X₁ ⊆ X₂ → ((X₁ ∩ C X₂) \ C X₁).card = 0
\end{ReviewVerbatim}

\paragraph{Lean proof endpoint (built separately)}\begin{ReviewVerbatim}
DGD26AdmissionsPredictability.PaperInterface.paper_substitutability_term_statement
\end{ReviewVerbatim}

\paragraph{Recorded source-to-Spec screening}
\textbf{Verdict:} \texttt{not recorded}
\\
\textbf{Reason:} not recorded
\reviewmatch{Matches source input}
\noindent\textbf{Reviewer annotation}\par
\reviewerbox
\clearpage
\hypertarget{source-claim-ca2a099855204aca}{}
\section*{Review row 38}
\textbf{Source locator:} {\footnotesize\raggedright cited publication, lines 532-\allowbreak{}546\par}
\paragraph{Verbatim source input}
\begin{ReviewVerbatim}
\begin{restatable}{lemma}{unsubstitutable}
\label{lem:unsubstitutable}
If $\choice$ is not substitutable, there exists some $x^*$, $x'$, $X$ and $X' = X \cup \{x'\}$ where $x^* \in X$ and $x^* \in \choice(X')$ but $x^* \not\in \choice(X)$.
\end{restatable}

\textit{Intuition}:
If a choice function is not substitutable, then there is at least one situation where a decision ``flips'' from rejection to acceptance.
We argue that there has to be some single element $x'$ that causes this ``flip'' when added.

\begin{proof}
$\choice$ is not substitutable if there exists some  $x^*$ and $X_0 \subseteq X_k$ where $x^* \in X_0$ and $x^* \in \choice(X_k)$ but $x \not\in \choice(X_0)$.
Let $k$ be such that $k = |X_k \setminus X_0|$.%

We want to show the existence of $x'$. To do this, we identify the ``closest'' sets where this occurs. For each $x \in X_k \setminus X_0$, add elements to $X_0$, constructing $X_0 \subseteq X_1 \subseteq ... \subseteq X_j \subseteq ... \subseteq X_k$, where $X_1 = X_0 \cup \{x_1\}$ and so on. Let $X_j$ be the first in this sequence of sets such that $x^* \in \choice(X_j)$ where $x^* \in \choice(X_j)$ but $x^* \not\in \choice(X_{j-1})$. Then let $x' = x_j$, $X = X_{j-1}$, and $X' = X_j$.
\end{proof}
\end{ReviewVerbatim}


\paragraph{Expanded PaperInterface specification}
\begin{ReviewVerbatim}
∀ {α : Type u_1} [inst : DecidableEq α] (C : DGD26AdmissionsPredictability.PaperChoiceRule α),
  ¬DGD26AdmissionsPredictability.paper_definition_substitutability C →
    ∃ X x xstar, x ∉ X ∧ xstar ∈ X ∧ xstar ∈ C (insert x X) ∧ xstar ∉ C X
\end{ReviewVerbatim}

\paragraph{Lean proof endpoint (built separately)}\begin{ReviewVerbatim}
DGD26AdmissionsPredictability.PaperInterface.paper_non_substitutable_single_add_statement
\end{ReviewVerbatim}

\paragraph{Recorded source-to-Spec screening}
\textbf{Verdict:} \texttt{not recorded}
\\
\textbf{Reason:} not recorded
\reviewmatch{Matches source input}
\noindent\textbf{Reviewer annotation}\par
\reviewerbox
\clearpage
\hypertarget{source-claim-bc7eeffb236e6569}{}
\section*{Review row 39}
\textbf{Source locator:} {\footnotesize\raggedright cited publication, lines 261-\allowbreak{}265\par}
\paragraph{Verbatim source input}
\begin{ReviewVerbatim}
\begin{restatable}[ML Representation]{proposition}{mlrepresentation}
\label{prop:ml-representation}
A model of the form~\eqref{eq:ml-independent} which makes independent predictions can only represent 0-unstable choice functions.
A model of the form~\eqref{eq:ml-rank} which ranks applicants can represent a 1-variable, 1-unstable choice function. No such model can represent a choice function with variability or instability greater than one.
\end{restatable}

[Next verbatim source excerpt]

\subsection{Proof of \Cref{prop:ml-representation}} \label{sec:ml-prop-proof}

Recall \Cref{prop:ml-representation}:
\mlrepresentation*

First consider using an ML model to make decisions; then it can be seen as a choice function with various properties.
Thus we say that an ML model can represent a choice function with various properties if it can have those properties when viewed as a choice function. %

The first statement follows by noting that any rule of the form~\eqref{eq:ml-independent} must either admit or reject an applicant regardless of the applicant pool, so it must be independent and hence $0$-unstable by \Cref{thm:independent-zero-unstable}. %
The second and third statements follow by noting that scores in~\eqref{eq:ml-rank} induce a total order by embedding applicants in $\mathbb R$ and selecting accordingly. Therefore any rule of the form~\eqref{eq:ml-rank} is $q$-representative (by definition) and hence has instability and variability exactly one (by \Cref{thm:ranking-variability}), but no greater.
\end{ReviewVerbatim}


\paragraph{Expanded PaperInterface specification}
\begin{ReviewVerbatim}
∀ {α : Type u_1} [inst : DecidableEq α] {C : DGD26AdmissionsPredictability.PaperChoiceRule α} (score : α → ℝ)
  (threshold : ℝ),
  DGD26AdmissionsPredictability.paper_choice_function_feasible C →
    DGD26AdmissionsPredictability.paper_definition_ml_representation
        (DGD26AdmissionsPredictability.paper_definition_fixed_threshold_predictor score threshold) C →
      DGD26AdmissionsPredictability.paper_definition_zero_instability C
\end{ReviewVerbatim}

\paragraph{Lean proof endpoint (built separately)}\begin{ReviewVerbatim}
DGD26AdmissionsPredictability.PaperInterface.paper_ml_fixed_threshold_representation_zero_unstable_statement
\end{ReviewVerbatim}

\paragraph{Recorded source-to-Spec screening}
\textbf{Verdict:} \texttt{not recorded}
\\
\textbf{Reason:} not recorded
\reviewmatch{Matches source input}
\noindent\textbf{Reviewer annotation}\par
\reviewerbox
\clearpage
\hypertarget{source-claim-d21050cec8d233b0}{}
\section*{Review row 40}
\textbf{Source locator:} {\footnotesize\raggedright cited publication, lines 261-\allowbreak{}269\par}
\paragraph{Verbatim source input}
\begin{ReviewVerbatim}
\begin{restatable}[ML Representation]{proposition}{mlrepresentation}
\label{prop:ml-representation}
A model of the form~\eqref{eq:ml-independent} which makes independent predictions can only represent 0-unstable choice functions.
A model of the form~\eqref{eq:ml-rank} which ranks applicants can represent a 1-variable, 1-unstable choice function. No such model can represent a choice function with variability or instability greater than one.
\end{restatable}

This result shows that standard ML practice coherently applies to admissions data only in limited settings.
The intuition for this result hinges on the fact that all 1-unstable and 1-variable choice functions correspond to total preference orderings over applicants (formalized in Theorem~\ref{thm:variability}).
Then the proposition follows from noticing that total preference orderings correspond to  embedding applicants in  $\mathbb{R}$  (by their scores).

[Next verbatim source excerpt]

\subsection{Proof of \Cref{prop:ml-representation}} \label{sec:ml-prop-proof}

Recall \Cref{prop:ml-representation}:
\mlrepresentation*

First consider using an ML model to make decisions; then it can be seen as a choice function with various properties.
Thus we say that an ML model can represent a choice function with various properties if it can have those properties when viewed as a choice function. %

The first statement follows by noting that any rule of the form~\eqref{eq:ml-independent} must either admit or reject an applicant regardless of the applicant pool, so it must be independent and hence $0$-unstable by \Cref{thm:independent-zero-unstable}. %
The second and third statements follow by noting that scores in~\eqref{eq:ml-rank} induce a total order by embedding applicants in $\mathbb R$ and selecting accordingly. Therefore any rule of the form~\eqref{eq:ml-rank} is $q$-representative (by definition) and hence has instability and variability exactly one (by \Cref{thm:ranking-variability}), but no greater.
\end{ReviewVerbatim}


\paragraph{Expanded PaperInterface specification}
\begin{ReviewVerbatim}
∀ (q : ℕ),
  0 < q →
    ∃ score,
      ∃ (hinjective : Function.Injective score),
        DGD26AdmissionsPredictability.paper_definition_tight_d_instability 1
            (DGD26AdmissionsPredictability.paperRankThresholdChoice q score hinjective) ∧
          DGD26AdmissionsPredictability.paper_definition_variability_exactly 1
            (DGD26AdmissionsPredictability.paperRankThresholdChoice q score hinjective)
\end{ReviewVerbatim}

\paragraph{Lean proof endpoint (built separately)}\begin{ReviewVerbatim}
DGD26AdmissionsPredictability.PaperInterface.paper_ml_rank_threshold_can_represent_exact_one_statement
\end{ReviewVerbatim}

\paragraph{Recorded source-to-Spec screening}
\textbf{Verdict:} \texttt{not recorded}
\\
\textbf{Reason:} not recorded
\reviewmatch{Matches source input}
\noindent\textbf{Reviewer annotation}\par
\reviewerbox
\clearpage
\hypertarget{source-claim-47a9f492a25c4e98}{}
\section*{Review row 41}
\textbf{Source locator:} {\footnotesize\raggedright cited publication, lines 261-\allowbreak{}269\par}
\paragraph{Verbatim source input}
\begin{ReviewVerbatim}
\begin{restatable}[ML Representation]{proposition}{mlrepresentation}
\label{prop:ml-representation}
A model of the form~\eqref{eq:ml-independent} which makes independent predictions can only represent 0-unstable choice functions.
A model of the form~\eqref{eq:ml-rank} which ranks applicants can represent a 1-variable, 1-unstable choice function. No such model can represent a choice function with variability or instability greater than one.
\end{restatable}

This result shows that standard ML practice coherently applies to admissions data only in limited settings.
The intuition for this result hinges on the fact that all 1-unstable and 1-variable choice functions correspond to total preference orderings over applicants (formalized in Theorem~\ref{thm:variability}).
Then the proposition follows from noticing that total preference orderings correspond to  embedding applicants in  $\mathbb{R}$  (by their scores).

[Next verbatim source excerpt]

\subsection{Proof of \Cref{prop:ml-representation}} \label{sec:ml-prop-proof}

Recall \Cref{prop:ml-representation}:
\mlrepresentation*

First consider using an ML model to make decisions; then it can be seen as a choice function with various properties.
Thus we say that an ML model can represent a choice function with various properties if it can have those properties when viewed as a choice function. %

The first statement follows by noting that any rule of the form~\eqref{eq:ml-independent} must either admit or reject an applicant regardless of the applicant pool, so it must be independent and hence $0$-unstable by \Cref{thm:independent-zero-unstable}. %
The second and third statements follow by noting that scores in~\eqref{eq:ml-rank} induce a total order by embedding applicants in $\mathbb R$ and selecting accordingly. Therefore any rule of the form~\eqref{eq:ml-rank} is $q$-representative (by definition) and hence has instability and variability exactly one (by \Cref{thm:ranking-variability}), but no greater.
\end{ReviewVerbatim}


\paragraph{Expanded PaperInterface specification}
\begin{ReviewVerbatim}
∀ {α : Type u_1} [inst : DecidableEq α] [Fintype α] {C : DGD26AdmissionsPredictability.PaperChoiceRule α} (q : ℕ)
  (score : α → ℝ) (hinjective : Function.Injective score),
  DGD26AdmissionsPredictability.paper_choice_function_feasible C →
    DGD26AdmissionsPredictability.paper_definition_ml_representation
        (DGD26AdmissionsPredictability.paper_definition_rank_threshold_predictor q score hinjective) C →
      DGD26AdmissionsPredictability.paper_definition_d_instability 1 C
\end{ReviewVerbatim}

\paragraph{Lean proof endpoint (built separately)}\begin{ReviewVerbatim}
DGD26AdmissionsPredictability.PaperInterface.paper_ml_rank_threshold_representation_instability_bound_statement
\end{ReviewVerbatim}

\paragraph{Recorded source-to-Spec screening}
\textbf{Verdict:} \texttt{not recorded}
\\
\textbf{Reason:} not recorded
\reviewmatch{Matches source input}
\noindent\textbf{Reviewer annotation}\par
\reviewerbox
\clearpage
\hypertarget{source-claim-98ed151e7543b81d}{}
\section*{Review row 42}
\textbf{Source locator:} {\footnotesize\raggedright cited publication, lines 261-\allowbreak{}269\par}
\paragraph{Verbatim source input}
\begin{ReviewVerbatim}
\begin{restatable}[ML Representation]{proposition}{mlrepresentation}
\label{prop:ml-representation}
A model of the form~\eqref{eq:ml-independent} which makes independent predictions can only represent 0-unstable choice functions.
A model of the form~\eqref{eq:ml-rank} which ranks applicants can represent a 1-variable, 1-unstable choice function. No such model can represent a choice function with variability or instability greater than one.
\end{restatable}

This result shows that standard ML practice coherently applies to admissions data only in limited settings.
The intuition for this result hinges on the fact that all 1-unstable and 1-variable choice functions correspond to total preference orderings over applicants (formalized in Theorem~\ref{thm:variability}).
Then the proposition follows from noticing that total preference orderings correspond to  embedding applicants in  $\mathbb{R}$  (by their scores).

[Next verbatim source excerpt]

\subsection{Proof of \Cref{prop:ml-representation}} \label{sec:ml-prop-proof}

Recall \Cref{prop:ml-representation}:
\mlrepresentation*

First consider using an ML model to make decisions; then it can be seen as a choice function with various properties.
Thus we say that an ML model can represent a choice function with various properties if it can have those properties when viewed as a choice function. %

The first statement follows by noting that any rule of the form~\eqref{eq:ml-independent} must either admit or reject an applicant regardless of the applicant pool, so it must be independent and hence $0$-unstable by \Cref{thm:independent-zero-unstable}. %
The second and third statements follow by noting that scores in~\eqref{eq:ml-rank} induce a total order by embedding applicants in $\mathbb R$ and selecting accordingly. Therefore any rule of the form~\eqref{eq:ml-rank} is $q$-representative (by definition) and hence has instability and variability exactly one (by \Cref{thm:ranking-variability}), but no greater.
\end{ReviewVerbatim}


\paragraph{Expanded PaperInterface specification}
\begin{ReviewVerbatim}
∀ {α : Type u_1} [inst : DecidableEq α] [inst_1 : Fintype α] {C : DGD26AdmissionsPredictability.PaperChoiceRule α}
  (q : ℕ) (score : α → ℝ) (hinjective : Function.Injective score),
  DGD26AdmissionsPredictability.paper_choice_function_feasible C →
    DGD26AdmissionsPredictability.paper_definition_ml_representation
        (DGD26AdmissionsPredictability.paper_definition_rank_threshold_predictor q score hinjective) C →
      DGD26AdmissionsPredictability.paper_definition_variability_at_most 1 C
\end{ReviewVerbatim}

\paragraph{Lean proof endpoint (built separately)}\begin{ReviewVerbatim}
DGD26AdmissionsPredictability.PaperInterface.paper_ml_rank_threshold_representation_variability_bound_statement
\end{ReviewVerbatim}

\paragraph{Recorded source-to-Spec screening}
\textbf{Verdict:} \texttt{not recorded}
\\
\textbf{Reason:} not recorded
\reviewmatch{Matches source input}
\noindent\textbf{Reviewer annotation}\par
\reviewerbox
\clearpage
\hypertarget{source-claim-396063f8f4209701}{}
\section*{Review row 43}
\textbf{Source locator:} {\footnotesize\raggedright cited publication, lines 347-\allowbreak{}358\par}
\paragraph{Verbatim source input}
\begin{ReviewVerbatim}
\begin{restatable}[]{proposition}{empiricalmethods}
\label{prop:program-variability}
All considered choice functions ({Ed. Opt}, {Screened/Open}, {Ed. Opt with DIA}, and {Screened/Open with DIA}) are $1$-unstable. Furthermore, their variability is (tightly) equal to their number of queues:

\begin{itemize}[noitemsep,topsep=0pt,parsep=0pt,partopsep=0pt]
    \item Screened/Open programs are $1$-variable
    \item Screened/Open with DIA programs are $2$-variable
    \item Ed. Opt. programs are $3$-variable
    \item Ed. Opt with DIA programs are $6$-variable
\end{itemize}

\end{restatable}

[Next verbatim source excerpt]

\subsection{Proof of \Cref{prop:program-variability}} \label{sec:program-prop-proof}

Recall \Cref{prop:program-variability}:
\empiricalmethods*

We construct applicant sets to prove the variability of each empirical method we study, showing that the upper bound from \Cref{thm:variability} based on the number of queues is tight.

Note that both Screened and Open programs are explicitly a ranking, and therefore have variability $1$ by  \Cref{thm:variability}.
Screened and Open programs with DIA are upper-bounded by variability 2, as they are composed of two queues. It is easy to see that this upper bound is tight. Consider an applicant pool $X$ where the lowest-ranking students in the DIA queue is lower-ranked than the lowest-ranked student in the non-DIA queue. Then $A_\choice(X)$ contains two elements -- adding a new DIA-qualified student might lead to rejecting the lowest-ranked student in the DIA queue, while a non-DIA student would lead to a change in the non-DIA queue. Similarly, in the three-queue Ed. Opt. case, the variability $3$ bound is tight, as there are applicant pools where adding a single high-, middle-, or low-category student would displace a different student in the corresponding queue. Likewise, the six queues in Ed. Opt. with DIA programs lead to six diffrent students who might be displaced depending on what category a new applicant falls into.
\end{ReviewVerbatim}


\paragraph{Expanded PaperInterface specification}
\begin{ReviewVerbatim}
DGD26AdmissionsPredictability.paper_definition_d_instability 1
    DGD26AdmissionsPredictability.paper_definition_screened_open_program_choice ∧
  DGD26AdmissionsPredictability.paper_definition_d_instability 1
      DGD26AdmissionsPredictability.paper_definition_screened_open_dia_program_choice ∧
    DGD26AdmissionsPredictability.paper_definition_d_instability 1
        DGD26AdmissionsPredictability.paper_definition_educational_option_program_choice ∧
      DGD26AdmissionsPredictability.paper_definition_d_instability 1
        DGD26AdmissionsPredictability.paper_definition_educational_option_dia_program_choice
\end{ReviewVerbatim}

\paragraph{Lean proof endpoint (built separately)}\begin{ReviewVerbatim}
DGD26AdmissionsPredictability.PaperInterface.paper_program_classes_one_instability_statement
\end{ReviewVerbatim}

\paragraph{Recorded source-to-Spec screening}
\textbf{Verdict:} \texttt{not recorded}
\\
\textbf{Reason:} not recorded
\reviewmatch{Matches source input}
\noindent\textbf{Reviewer annotation}\par
\reviewerbox
\clearpage
\hypertarget{source-claim-f595f3e6af0e1b40}{}
\section*{Review row 44}
\textbf{Source locator:} {\footnotesize\raggedright cited publication, lines 347-\allowbreak{}358\par}
\paragraph{Verbatim source input}
\begin{ReviewVerbatim}
\begin{restatable}[]{proposition}{empiricalmethods}
\label{prop:program-variability}
All considered choice functions ({Ed. Opt}, {Screened/Open}, {Ed. Opt with DIA}, and {Screened/Open with DIA}) are $1$-unstable. Furthermore, their variability is (tightly) equal to their number of queues:

\begin{itemize}[noitemsep,topsep=0pt,parsep=0pt,partopsep=0pt]
    \item Screened/Open programs are $1$-variable
    \item Screened/Open with DIA programs are $2$-variable
    \item Ed. Opt. programs are $3$-variable
    \item Ed. Opt with DIA programs are $6$-variable
\end{itemize}

\end{restatable}

[Next verbatim source excerpt]

\subsection{Proof of \Cref{prop:program-variability}} \label{sec:program-prop-proof}

Recall \Cref{prop:program-variability}:
\empiricalmethods*

We construct applicant sets to prove the variability of each empirical method we study, showing that the upper bound from \Cref{thm:variability} based on the number of queues is tight.

Note that both Screened and Open programs are explicitly a ranking, and therefore have variability $1$ by  \Cref{thm:variability}.
Screened and Open programs with DIA are upper-bounded by variability 2, as they are composed of two queues. It is easy to see that this upper bound is tight. Consider an applicant pool $X$ where the lowest-ranking students in the DIA queue is lower-ranked than the lowest-ranked student in the non-DIA queue. Then $A_\choice(X)$ contains two elements -- adding a new DIA-qualified student might lead to rejecting the lowest-ranked student in the DIA queue, while a non-DIA student would lead to a change in the non-DIA queue. Similarly, in the three-queue Ed. Opt. case, the variability $3$ bound is tight, as there are applicant pools where adding a single high-, middle-, or low-category student would displace a different student in the corresponding queue. Likewise, the six queues in Ed. Opt. with DIA programs lead to six diffrent students who might be displaced depending on what category a new applicant falls into.
\end{ReviewVerbatim}


\paragraph{Expanded PaperInterface specification}
\begin{ReviewVerbatim}
DGD26AdmissionsPredictability.paper_definition_variability_exactly 1
  DGD26AdmissionsPredictability.paper_definition_screened_open_program_choice
\end{ReviewVerbatim}

\paragraph{Lean proof endpoint (built separately)}\begin{ReviewVerbatim}
DGD26AdmissionsPredictability.PaperInterface.paper_screened_open_variability_one_statement
\end{ReviewVerbatim}

\paragraph{Recorded source-to-Spec screening}
\textbf{Verdict:} \texttt{not recorded}
\\
\textbf{Reason:} not recorded
\reviewmatch{Matches source input}
\noindent\textbf{Reviewer annotation}\par
\reviewerbox
\clearpage
\hypertarget{source-claim-8afe5dc5804dcb73}{}
\section*{Review row 45}
\textbf{Source locator:} {\footnotesize\raggedright cited publication, lines 347-\allowbreak{}358\par}
\paragraph{Verbatim source input}
\begin{ReviewVerbatim}
\begin{restatable}[]{proposition}{empiricalmethods}
\label{prop:program-variability}
All considered choice functions ({Ed. Opt}, {Screened/Open}, {Ed. Opt with DIA}, and {Screened/Open with DIA}) are $1$-unstable. Furthermore, their variability is (tightly) equal to their number of queues:

\begin{itemize}[noitemsep,topsep=0pt,parsep=0pt,partopsep=0pt]
    \item Screened/Open programs are $1$-variable
    \item Screened/Open with DIA programs are $2$-variable
    \item Ed. Opt. programs are $3$-variable
    \item Ed. Opt with DIA programs are $6$-variable
\end{itemize}

\end{restatable}

[Next verbatim source excerpt]

\subsection{Proof of \Cref{prop:program-variability}} \label{sec:program-prop-proof}

Recall \Cref{prop:program-variability}:
\empiricalmethods*

We construct applicant sets to prove the variability of each empirical method we study, showing that the upper bound from \Cref{thm:variability} based on the number of queues is tight.

Note that both Screened and Open programs are explicitly a ranking, and therefore have variability $1$ by  \Cref{thm:variability}.
Screened and Open programs with DIA are upper-bounded by variability 2, as they are composed of two queues. It is easy to see that this upper bound is tight. Consider an applicant pool $X$ where the lowest-ranking students in the DIA queue is lower-ranked than the lowest-ranked student in the non-DIA queue. Then $A_\choice(X)$ contains two elements -- adding a new DIA-qualified student might lead to rejecting the lowest-ranked student in the DIA queue, while a non-DIA student would lead to a change in the non-DIA queue. Similarly, in the three-queue Ed. Opt. case, the variability $3$ bound is tight, as there are applicant pools where adding a single high-, middle-, or low-category student would displace a different student in the corresponding queue. Likewise, the six queues in Ed. Opt. with DIA programs lead to six diffrent students who might be displaced depending on what category a new applicant falls into.
\end{ReviewVerbatim}


\paragraph{Expanded PaperInterface specification}
\begin{ReviewVerbatim}
DGD26AdmissionsPredictability.paper_definition_variability_exactly 2
  DGD26AdmissionsPredictability.paper_definition_screened_open_dia_program_choice
\end{ReviewVerbatim}

\paragraph{Lean proof endpoint (built separately)}\begin{ReviewVerbatim}
DGD26AdmissionsPredictability.PaperInterface.paper_screened_open_dia_variability_two_statement
\end{ReviewVerbatim}

\paragraph{Recorded source-to-Spec screening}
\textbf{Verdict:} \texttt{not recorded}
\\
\textbf{Reason:} not recorded
\reviewmatch{Matches source input}
\noindent\textbf{Reviewer annotation}\par
\reviewerbox
\clearpage
\hypertarget{source-claim-68c8be83630deedf}{}
\section*{Review row 46}
\textbf{Source locator:} {\footnotesize\raggedright cited publication, lines 347-\allowbreak{}358\par}
\paragraph{Verbatim source input}
\begin{ReviewVerbatim}
\begin{restatable}[]{proposition}{empiricalmethods}
\label{prop:program-variability}
All considered choice functions ({Ed. Opt}, {Screened/Open}, {Ed. Opt with DIA}, and {Screened/Open with DIA}) are $1$-unstable. Furthermore, their variability is (tightly) equal to their number of queues:

\begin{itemize}[noitemsep,topsep=0pt,parsep=0pt,partopsep=0pt]
    \item Screened/Open programs are $1$-variable
    \item Screened/Open with DIA programs are $2$-variable
    \item Ed. Opt. programs are $3$-variable
    \item Ed. Opt with DIA programs are $6$-variable
\end{itemize}

\end{restatable}

[Next verbatim source excerpt]

\subsection{Proof of \Cref{prop:program-variability}} \label{sec:program-prop-proof}

Recall \Cref{prop:program-variability}:
\empiricalmethods*

We construct applicant sets to prove the variability of each empirical method we study, showing that the upper bound from \Cref{thm:variability} based on the number of queues is tight.

Note that both Screened and Open programs are explicitly a ranking, and therefore have variability $1$ by  \Cref{thm:variability}.
Screened and Open programs with DIA are upper-bounded by variability 2, as they are composed of two queues. It is easy to see that this upper bound is tight. Consider an applicant pool $X$ where the lowest-ranking students in the DIA queue is lower-ranked than the lowest-ranked student in the non-DIA queue. Then $A_\choice(X)$ contains two elements -- adding a new DIA-qualified student might lead to rejecting the lowest-ranked student in the DIA queue, while a non-DIA student would lead to a change in the non-DIA queue. Similarly, in the three-queue Ed. Opt. case, the variability $3$ bound is tight, as there are applicant pools where adding a single high-, middle-, or low-category student would displace a different student in the corresponding queue. Likewise, the six queues in Ed. Opt. with DIA programs lead to six diffrent students who might be displaced depending on what category a new applicant falls into.
\end{ReviewVerbatim}


\paragraph{Expanded PaperInterface specification}
\begin{ReviewVerbatim}
DGD26AdmissionsPredictability.paper_definition_variability_exactly 3
  DGD26AdmissionsPredictability.paper_definition_educational_option_program_choice
\end{ReviewVerbatim}

\paragraph{Lean proof endpoint (built separately)}\begin{ReviewVerbatim}
DGD26AdmissionsPredictability.PaperInterface.paper_educational_option_variability_three_statement
\end{ReviewVerbatim}

\paragraph{Recorded source-to-Spec screening}
\textbf{Verdict:} \texttt{not recorded}
\\
\textbf{Reason:} not recorded
\reviewmatch{Matches source input}
\noindent\textbf{Reviewer annotation}\par
\reviewerbox
\clearpage
\hypertarget{source-claim-a241fbe5321a55d3}{}
\section*{Review row 47}
\textbf{Source locator:} {\footnotesize\raggedright cited publication, lines 347-\allowbreak{}358\par}
\paragraph{Verbatim source input}
\begin{ReviewVerbatim}
\begin{restatable}[]{proposition}{empiricalmethods}
\label{prop:program-variability}
All considered choice functions ({Ed. Opt}, {Screened/Open}, {Ed. Opt with DIA}, and {Screened/Open with DIA}) are $1$-unstable. Furthermore, their variability is (tightly) equal to their number of queues:

\begin{itemize}[noitemsep,topsep=0pt,parsep=0pt,partopsep=0pt]
    \item Screened/Open programs are $1$-variable
    \item Screened/Open with DIA programs are $2$-variable
    \item Ed. Opt. programs are $3$-variable
    \item Ed. Opt with DIA programs are $6$-variable
\end{itemize}

\end{restatable}

[Next verbatim source excerpt]

\subsection{Proof of \Cref{prop:program-variability}} \label{sec:program-prop-proof}

Recall \Cref{prop:program-variability}:
\empiricalmethods*

We construct applicant sets to prove the variability of each empirical method we study, showing that the upper bound from \Cref{thm:variability} based on the number of queues is tight.

Note that both Screened and Open programs are explicitly a ranking, and therefore have variability $1$ by  \Cref{thm:variability}.
Screened and Open programs with DIA are upper-bounded by variability 2, as they are composed of two queues. It is easy to see that this upper bound is tight. Consider an applicant pool $X$ where the lowest-ranking students in the DIA queue is lower-ranked than the lowest-ranked student in the non-DIA queue. Then $A_\choice(X)$ contains two elements -- adding a new DIA-qualified student might lead to rejecting the lowest-ranked student in the DIA queue, while a non-DIA student would lead to a change in the non-DIA queue. Similarly, in the three-queue Ed. Opt. case, the variability $3$ bound is tight, as there are applicant pools where adding a single high-, middle-, or low-category student would displace a different student in the corresponding queue. Likewise, the six queues in Ed. Opt. with DIA programs lead to six diffrent students who might be displaced depending on what category a new applicant falls into.
\end{ReviewVerbatim}


\paragraph{Expanded PaperInterface specification}
\begin{ReviewVerbatim}
DGD26AdmissionsPredictability.paper_definition_variability_exactly 6
  DGD26AdmissionsPredictability.paper_definition_educational_option_dia_program_choice
\end{ReviewVerbatim}

\paragraph{Lean proof endpoint (built separately)}\begin{ReviewVerbatim}
DGD26AdmissionsPredictability.PaperInterface.paper_educational_option_dia_variability_six_statement
\end{ReviewVerbatim}

\paragraph{Recorded source-to-Spec screening}
\textbf{Verdict:} \texttt{not recorded}
\\
\textbf{Reason:} not recorded
\reviewmatch{Matches source input}
\noindent\textbf{Reviewer annotation}\par
\reviewerbox
\clearpage
\hypertarget{source-claim-207dcf4700d4343b}{}
\section*{Review row 48}
\textbf{Source locator:} {\footnotesize\raggedright cited publication, lines 998-\allowbreak{}1003\par}
\paragraph{Verbatim source input}
\begin{ReviewVerbatim}
\begin{restatable}[Additive Variability Bound]{theorem}{addvar}
\label{thm:additive-variability}
If $\choice$ is a sequential composition of $q$-acceptant, $1$-unstable functions with variabilities $m_1, \ldots, m_n$, then $\choice$ has variability at most $\sum_{i=1}^n m_i$.
\end{restatable}

Note that this is only an upper bound; we could trivially construct cases where this is not tight. For example, compose a $\choice$ of two $q$-representative choice functions with the same underlying ordering. On the other hand, it is also easy to construct instances in which this upper bound is tight. Consider the composition of $m$ queues which each rank a student by a different subject area exam. Then the addition of a new student to the applicant pool could displace any of $m$ students, depending on the subject area.
\end{ReviewVerbatim}


\paragraph{Expanded PaperInterface specification}
\begin{ReviewVerbatim}
∀ {α : Type u_1} [inst : DecidableEq α] [inst_1 : Fintype α] {qs ms : List ℕ}
  {Cs : List (DGD26AdmissionsPredictability.PaperChoiceRule α)},
  List.Forall₂
      (fun q C =>
        DGD26AdmissionsPredictability.paper_choice_function_feasible C ∧
          DGD26AdmissionsPredictability.paper_definition_q_acceptance q C ∧
            DGD26AdmissionsPredictability.paper_definition_d_instability 1 C)
      qs Cs →
    List.Forall₂ (fun m C => DGD26AdmissionsPredictability.paper_definition_variability_at_most m C) ms Cs →
      DGD26AdmissionsPredictability.paper_definition_variability_at_most ms.sum
        (DGD26AdmissionsPredictability.paper_definition_sequential_composition Cs)
\end{ReviewVerbatim}

\paragraph{Lean proof endpoint (built separately)}\begin{ReviewVerbatim}
DGD26AdmissionsPredictability.PaperInterface.paper_sequential_additive_variability_bound_statement
\end{ReviewVerbatim}

\paragraph{Recorded source-to-Spec screening}
\textbf{Verdict:} \texttt{not recorded}
\\
\textbf{Reason:} not recorded
\reviewmatch{Matches source input}
\noindent\textbf{Reviewer annotation}\par
\reviewerbox
\clearpage
\hypertarget{source-claim-fff538b4a617e0f3}{}
\section*{Review row 49}
\textbf{Source locator:} {\footnotesize\raggedright cited publication, lines 850-\allowbreak{}879\par}
\paragraph{Verbatim source input}
\begin{ReviewVerbatim}
\begin{restatable}{theorem}{appendremove}
\label{thm:append-remove-size}
For $q$-acceptant, 1-unstable $\choice$, and $|\calX| \geq 2q$, the maximum size of the waitlisted set is equal to the maximum size of the borderline set; that is,
\begin{equation}
\max_{X\subset \calX} \left| V_\choice(X)\right| = \max_{X\subset \calX} \left| A_\choice(X)\right|
\end{equation}
\end{restatable}

In short, if there is a set $X$ with a waitlisted set of size $u$, there must be a set $X_0$ with a borderline set of size (at least) $u$, and vice versa. %

\begin{proof}
To prove equality, we will prove that 1) $\max_{X\subset \calX} \left| V_\choice(X)\right| \geq \max_{X\subset \calX} \left| A_\choice(X)\right|$ and then that 2) $\max_{X\subset \calX} \left| V_\choice(X)\right| \leq  \max_{X\subset \calX} \left| A_\choice(X)\right|$.

\textbf{Part 1).}
To prove the first statement, assume that there is some $X$ where the waitlisted set is of cardinality at least $u$, so $x^*_1, ..., x^*_u \in A_\choice(X)$.

Then (recall that $\choice$ is consistent) there are $X_1, ..., X_u \subseteq X$ where $X_i = X \setminus \{x'_i\}$ such that $\choice(X_i) = x^*_i \cup \choice(X) \setminus x'_i$ for $i \in \{1, ..., u\}$. Then consider $X_0 = \bigcap_{i=1}^u X_i = X \setminus \{x'_1, ..., x'_u\}$, and $\choice(X_0)$. It must be that $\choice(X_0) = \{x^*_1, ..., x^*_u\} \cup \choice(X) \setminus \{x'_1, ..., x'_u\}$ by substitutability. I claim that $\{x^*_1, ..., x^*_u\} \subseteq V_\choice(X_0)$, which means $|V_c(X_0)| \geq u$, which would complete the proof.

Assume for contradiction that this is untrue; $\{x^*_1, ..., x^*_u\} \not\subseteq V_\choice(X_0)$. Then some $x^*_i$, say $x^*_u$ (without loss of generality), is not in $V_\choice(X_0)$, which means that $x^*_u \in \choice(\{x'_u\} \cup X_0)$.
Consider the makeup of $\choice(\{x'_u\} \cup X_0)$. Recall that $\{x'_u\} \cup X_0 = \bigcap_{i=1}^{u-1} X_i = X \setminus \{x'_1, ..., x'_{u-1}\}$. Because of substitutability, as a subset of $X$, we have $\left(\choice(X) \setminus \{x'_1, ..., x'_{u-1}\}\right) \subseteq \{x'_u\} \cup X_0$. From substitutability, as a subset of from $X_1$ through $X_{u-1}$, $\{x^*_2, ..., x^*_{u-1}\}$ are likewise chosen.
However, consider the size of $\choice(\{x'_u\} \cup X_0)$. $|\choice(X) \setminus \{x'_1, ..., x'_{u-1}\}|$ is of size $q-(u-1)$, because $x'_i \in \choice(X)$ for all $i$. We also know that $|\{x^*_1, ..., x^*_{u-1}\}| = u-1$, and as no  $x^*_i$ is in $\choice(X)$ by definition, all of these elements are distinct. As $x^*_u$ remains by assumption, this adds up to $|\choice(\{x'_u\} \cup X_0)| = q+1$, which violates $q$-acceptance and completes the proof.

\textbf{Part 2).}
The proof for 2) is extremely similar to the former. Begin with $X_0$ with $V_\choice(X_0) = \{x^*_1, ..., x^*_m\}$, with $x'_1, ..., x'_m$ such that $\choice(X'_1) = \{x'_1\} \cup \choice(X_0) \setminus \{x^*_1\}$. Let $X = \bigcup_{i=1}^m /X'_i = X_0 \cup \{x'_1,...,x'_m\}$ and $X_i = X \setminus \{x'_i\}$. Note that $X_i = X'_1 \cup ... X'_{i-1} \cup X'_{i+1} \cup ... \cup X_m$. %
Because of $1$-instability, $|\choice(X_0) \setminus \choice(X)| \leq |X_0 \setminus X| = m$. Because of substitutability, any $x \in \choice(X)$ is either in $\choice(X_0)$ or not in $X_0$ and therefore in $\{x'_1, ..., x'_m\}$.

I claim $\choice(X) = \{x'_1,...,x'_m\} \cup \choice(X_0) \setminus \{x^*_1,..., x^*_m\}$. Firstly, $x^*_i \not\in \choice(X)$ because  $x^*_i \not\in \choice(X'_i)$ and $X'_i \subseteq X$, so that would be a substitutability violation. As $\choice(X_0) \setminus \{x^*_1,..., x^*_m\}$ has $q-m$ elements, the only possible elements to fill $\choice(X)$ must be the $m$ elements of $\{x'_1,...,x'_m\}$.

The same argument holds for any $X_i$, but use $X_m$ without loss of generality for ease of notation; $X'_j \subseteq X_m$ for $m \neq j$, so $x^*_j \not\in\choice(X_m)$ by substitutability. With $\choice(X_0) \setminus \{x^*_1,..., x^*_{m-1}\}$ being the only $q-m+1$ elements of $X_0$ that can be in $\choice(X_m)$, the other $m-1$ elements must be $\{x'_1,. .., x'_{m-1}\}$. Then consider the definition of $A_\choice(X)$ and see that $\choice(X \setminus x'_i) \setminus \choice(X) = \{x^*_i\}$ for any $i$, with each $x^*_i$ being unique.
\end{proof}
\end{ReviewVerbatim}


\paragraph{Expanded PaperInterface specification}
\begin{ReviewVerbatim}
∀ {α : Type u_1} [inst : DecidableEq α] [inst_1 : Fintype α] {m q : ℕ}
  {C : DGD26AdmissionsPredictability.PaperChoiceRule α},
  DGD26AdmissionsPredictability.paper_choice_function_feasible C →
    DGD26AdmissionsPredictability.paper_definition_q_acceptance q C →
      DGD26AdmissionsPredictability.paper_definition_d_instability 1 C →
        (DGD26AdmissionsPredictability.paper_definition_general_variability_exactly m C ↔
          DGD26AdmissionsPredictability.paper_definition_variability_exactly m C)
\end{ReviewVerbatim}

\paragraph{Lean proof endpoint (built separately)}\begin{ReviewVerbatim}
DGD26AdmissionsPredictability.PaperInterface.paper_append_remove_variability_exact_equivalence_statement
\end{ReviewVerbatim}

\paragraph{Recorded source-to-Spec screening}
\textbf{Verdict:} \texttt{not recorded}
\\
\textbf{Reason:} not recorded
\reviewmatch{Matches source input}
\noindent\textbf{Reviewer annotation}\par
\reviewerbox
\clearpage
\hypertarget{source-claim-e1e0f5f892109672}{}
\section*{Review row 50}
\textbf{Source locator:} {\footnotesize\raggedright cited publication, lines 576-\allowbreak{}580\par}
\paragraph{Verbatim source input}
\begin{ReviewVerbatim}
\begin{restatable}[Triangle Inequality]{theorem}{triangle}
\label{thm:triangle-ineq}
For any choice function $\choice$ and sets $X_1 \subseteq X_2 \subseteq X_3$, $r_\choice$ obeys the triangle inequality;
$$r_\choice(X_1, X_3) \leq r_\choice(X_1, X_2) + r_\choice(X_2, X_3)$$
\end{restatable}
\end{ReviewVerbatim}


\paragraph{Expanded PaperInterface specification}
\begin{ReviewVerbatim}
∀ {α : Type u_1} [inst : DecidableEq α] (C : DGD26AdmissionsPredictability.PaperChoiceRule α) {X₁ X₂ X₃ : Finset α},
  X₁ ⊆ X₂ →
    X₂ ⊆ X₃ →
      DGD26AdmissionsPredictability.paper_definition_choice_distance C X₁ X₃ ≤
        DGD26AdmissionsPredictability.paper_definition_choice_distance C X₁ X₂ +
          DGD26AdmissionsPredictability.paper_definition_choice_distance C X₂ X₃
\end{ReviewVerbatim}

\paragraph{Lean proof endpoint (built separately)}\begin{ReviewVerbatim}
DGD26AdmissionsPredictability.PaperInterface.paper_choice_distance_triangle_statement
\end{ReviewVerbatim}

\paragraph{Recorded source-to-Spec screening}
\textbf{Verdict:} \texttt{not recorded}
\\
\textbf{Reason:} not recorded
\reviewmatch{Matches source input}
\noindent\textbf{Reviewer annotation}\par
\reviewerbox
\clearpage
\hypertarget{source-claim-f8de55f9c5760a87}{}
\section*{Review row 51}
\textbf{Source locator:} {\footnotesize\raggedright cited publication, lines 777-\allowbreak{}794\par}
\paragraph{Verbatim source input}
\begin{ReviewVerbatim}
\begin{restatable}[Even Instability is Inconsistent]{theorem}{inconsistentstable}
\label{thm:even-inconsistency}
For any $q$-acceptant $\choice$, $\choice$ is inconsistent if and only if there exist $X, X' = X \cup \{x'\}$ where $r_\choice(X, X')$ is positive and even.
\end{restatable}

\textit{Intuition}: for $\choice$ to have an even but nonzero choice distance $r_\choice$, the indicator term in  \Cref{lem:unstable-calc} must be zero; the added $x'$ is not selected. But that means $x'$ is a rejected element that still has an impact, which violates consistency.
\begin{proof}
\textbf{Informal}
First, we prove that an even and positive value of $r_\choice$ implies inconsistency. For $\choice$ to have an even but nonzero choice distance $r_\choice$, the indicator term in  \Cref{lem:unstable-calc} must be zero; the added $x'$ is not selected. But that means $x'$ affects selection without actually being chosen, which violates consistency.

Next, prove that consistency limits the values of $r_\choice$. If $\choice$ is not consistent, then there is some $X, X'$ where the new element $x'$ is not chosen (so the indicator term in \Cref{lem:unstable-calc} is zero and $r_\choice$ is even), but changes the chosen elements. This makes  $r_\choice$ is positive under a fixed capacity constraint.

\textbf{Formal}
Consider \Cref{lem:unstable-calc}. See that for $r_\choice \neq 0$, $r_\choice$ is odd if and only if $x' \in \choice(X')$ and even if and only if $x' \not\in \choice(X')$. If $x' \not\in \choice(X')$, then $\choice(X') \subseteq X$. %
So if $x' \not\in \choice(X')$ but $n > 0$, $\choice(X') \subseteq X$ but $\choice(X') \neq \choice(X)$, which proves inconsistency.
Similarly, if $\choice$ is inconsistent, there is $\choice(X') \subseteq X$ but $\choice(X') \neq \choice(X)$. %
As $\choice(X') \neq \choice(X)$, $n > 0$, but for $\choice(X') \subseteq X$ to be true, $x' \neq \choice(X')$, and so $r_\choice(X, X')$ is even.
\end{proof}
\end{ReviewVerbatim}


\paragraph{Expanded PaperInterface specification}
\begin{ReviewVerbatim}
∀ {α : Type u_1} [inst : DecidableEq α] {q : ℕ} {C : DGD26AdmissionsPredictability.PaperChoiceRule α},
  DGD26AdmissionsPredictability.paper_choice_function_feasible C →
    DGD26AdmissionsPredictability.paper_definition_q_acceptance q C →
      (¬DGD26AdmissionsPredictability.paper_definition_consistency C ↔
        ∃ X,
          ∃ x ∉ X,
            0 < DGD26AdmissionsPredictability.paper_definition_choice_distance C X (insert x X) ∧
              ∃ k, DGD26AdmissionsPredictability.paper_definition_choice_distance C X (insert x X) = 2 * k)
\end{ReviewVerbatim}

\paragraph{Lean proof endpoint (built separately)}\begin{ReviewVerbatim}
DGD26AdmissionsPredictability.PaperInterface.paper_even_distance_iff_inconsistent_statement
\end{ReviewVerbatim}

\paragraph{Recorded source-to-Spec screening}
\textbf{Verdict:} \texttt{not recorded}
\\
\textbf{Reason:} not recorded
\reviewmatch{Matches source input}
\noindent\textbf{Reviewer annotation}\par
\reviewerbox
\clearpage
\hypertarget{source-claim-3e732abc6f7366e9}{}
\section*{Review row 52}
\textbf{Source locator:} {\footnotesize\raggedright cited publication, lines 683-\allowbreak{}707\par}
\paragraph{Verbatim source input}
\begin{ReviewVerbatim}
\begin{restatable}{theorem}{Independent Rules are Substitutable and Monotonic}
\label{thm:independent-sub-mon}
$\choice$ is independent if and only if it is both substitutable and monotonic.
\end{restatable}

\textit{Intuition}: The only way for adding applicants to neither hurt those already accepted or help those already rejected is to not affect them at all.
\begin{proof}
    Assume $\choice$ is independent.

    \textbf{Informal}
    Consider how independence implies substitutability and monotonicity. Independence means any element is always accepted or always rejected regardless of applicant pool. Then any $x$ that is accepted in $X_1$ is still accepted in any subset of $X_1$, which satisfies the definition of substitutability, and is accepted in any superset of $X_1$, which satisfies monotonicity.

    To prove that substitutability and monotonicity combine to form independence, consider the behavior they cover. Recall the intuition behind substitutability --- removing applicants does not lead to rejecting those accepted, and the intuition behind monotonicity --- adding applicants does not lead to rejecting those accepted. Given both properties, nothing can affect an acceptance decision at all.

    \textbf{Formal}
    First, we assume independence and prove $0$-stablility.
    Recall \Cref{lem:substitutability-term}: $\choice$ is substitutable if and only if $|X_1 \cap \choice(X_2) \setminus \choice(X_1)| = 0$ for all $X_1 \subseteq X_2$. Then for any $x \in X_1$, either $x \in \choice(X_1),\choice(X_2)$ or $x \not\in \choice(X_1),\choice(X_2)$, so $X_1 \cap \choice(X_2) = \choice(X_1)$. Then $|X_1 \cap \choice(X_2) \setminus \choice(X_1)| = 0$ and so $\choice$ is substitutable. Similarly, recall \Cref{lem:monotonicity-term}: $\choice$ is monotonic if and only if $|\choice(X_1) \setminus \choice(X_2)| = 0$ for all $X_1 \subseteq X_2$. Any element $x \in \choice(X_1)$ is in $X_2$ because $\choice(X_1) \subseteq X_1 \subseteq X_2$. By definition of independence, $x \in \choice(X)$ for all $X$ where $x \in X$, so if $x \in \choice(X_1)$, $x \in \choice(X_2)$, so $\choice(X_1) \subseteq \choice(X_2)$.

    To prove that substitutability and monotonicity together imply independence, consider any $x \in X$. If $x \in \choice(X)$, we can show that $x \in \choice(X')$ for any $X'$ where $x \in X'$. By monotonicity, $x \in \choice(X' \cup X)$, and so by substitutability, $x \in \choice(X')$. If $x \not\in \choice(X)$, we can prove $x \not\in \choice(X')$ for any $X'$. If we assumed for contradiction that $x \in \choice(X')$, then the same argument as above proves $x \in \choice(X)$, leading to contradiction.
\end{proof}

\begin{restatable}{corollary}{Only Independent Rules are 0-unstable}
\label{thm:independent-zero-unstable}
$\choice$ is $0$-unstable if and only if it is independent, and therefore no $q$-acceptant $\choice$ can be $0$-unstable.
\end{restatable}
\end{ReviewVerbatim}


\paragraph{Expanded PaperInterface specification}
\begin{ReviewVerbatim}
∀ {α : Type u_1} [inst : DecidableEq α] (C : DGD26AdmissionsPredictability.PaperChoiceRule α),
  DGD26AdmissionsPredictability.paper_choice_function_feasible C →
    (DGD26AdmissionsPredictability.paper_definition_independence C ↔
      DGD26AdmissionsPredictability.paper_definition_substitutability C ∧
        DGD26AdmissionsPredictability.paper_definition_monotonicity C)
\end{ReviewVerbatim}

\paragraph{Lean proof endpoint (built separately)}\begin{ReviewVerbatim}
DGD26AdmissionsPredictability.PaperInterface.paper_independent_substitutable_monotonic_statement
\end{ReviewVerbatim}

\paragraph{Recorded source-to-Spec screening}
\textbf{Verdict:} \texttt{not recorded}
\\
\textbf{Reason:} not recorded
\reviewmatch{Matches source input}
\noindent\textbf{Reviewer annotation}\par
\reviewerbox
\clearpage
\hypertarget{source-claim-dac748564d615cb6}{}
\section*{Review row 53}
\textbf{Source locator:} {\footnotesize\raggedright cited publication, lines 280-\allowbreak{}289\par}
\paragraph{Verbatim source input}
\begin{ReviewVerbatim}
\begin{restatable}{theorem}{thminstability}
\label{thm:instability}
A $q$-acceptant choice function
\begin{enumerate}
    \item cannot be $0$-unstable,
    \item  is exactly $1$-unstable if and only if it is substitutable,
    \item can be tightly $d$-unstable for every $1\leq d \leq 2q$.
\end{enumerate}

\end{restatable}%
\end{ReviewVerbatim}


\paragraph{Expanded PaperInterface specification}
\begin{ReviewVerbatim}
∀ {α : Type u_1} [inst : DecidableEq α] {q : ℕ} {C : DGD26AdmissionsPredictability.PaperChoiceRule α},
  DGD26AdmissionsPredictability.paper_choice_function_feasible C →
    DGD26AdmissionsPredictability.paper_definition_q_acceptance q C →
      0 < q → ∀ {U : Finset α}, q < U.card → ¬DGD26AdmissionsPredictability.paper_definition_zero_instability C
\end{ReviewVerbatim}

\paragraph{Lean proof endpoint (built separately)}\begin{ReviewVerbatim}
DGD26AdmissionsPredictability.PaperInterface.paper_no_zero_instability_under_capacity_statement
\end{ReviewVerbatim}

\paragraph{Recorded source-to-Spec screening}
\textbf{Verdict:} \texttt{not recorded}
\\
\textbf{Reason:} not recorded
\reviewmatch{Matches source input}
\noindent\textbf{Reviewer annotation}\par
\reviewerbox
\clearpage
\hypertarget{source-claim-4efb65f58e4bc480}{}
\section*{Review row 54}
\textbf{Source locator:} {\footnotesize\raggedright cited publication, lines 280-\allowbreak{}289\par}
\paragraph{Verbatim source input}
\begin{ReviewVerbatim}
\begin{restatable}{theorem}{thminstability}
\label{thm:instability}
A $q$-acceptant choice function
\begin{enumerate}
    \item cannot be $0$-unstable,
    \item  is exactly $1$-unstable if and only if it is substitutable,
    \item can be tightly $d$-unstable for every $1\leq d \leq 2q$.
\end{enumerate}

\end{restatable}%
\end{ReviewVerbatim}


\paragraph{Expanded PaperInterface specification}
\begin{ReviewVerbatim}
∀ {α : Type u_1} [inst : DecidableEq α] {q : ℕ} {C : DGD26AdmissionsPredictability.PaperChoiceRule α},
  DGD26AdmissionsPredictability.paper_choice_function_feasible C →
    DGD26AdmissionsPredictability.paper_definition_q_acceptance q C →
      (DGD26AdmissionsPredictability.paper_definition_substitutability C ↔
        DGD26AdmissionsPredictability.paper_definition_d_instability 1 C)
\end{ReviewVerbatim}

\paragraph{Lean proof endpoint (built separately)}\begin{ReviewVerbatim}
DGD26AdmissionsPredictability.PaperInterface.paper_substitutability_one_instability_equivalence_statement
\end{ReviewVerbatim}

\paragraph{Recorded source-to-Spec screening}
\textbf{Verdict:} \texttt{not recorded}
\\
\textbf{Reason:} not recorded
\reviewmatch{Matches source input}
\noindent\textbf{Reviewer annotation}\par
\reviewerbox
\clearpage
\hypertarget{source-claim-d2442dd19c06f77a}{}
\section*{Review row 55}
\textbf{Source locator:} {\footnotesize\raggedright cited publication, lines 280-\allowbreak{}289\par}
\paragraph{Verbatim source input}
\begin{ReviewVerbatim}
\begin{restatable}{theorem}{thminstability}
\label{thm:instability}
A $q$-acceptant choice function
\begin{enumerate}
    \item cannot be $0$-unstable,
    \item  is exactly $1$-unstable if and only if it is substitutable,
    \item can be tightly $d$-unstable for every $1\leq d \leq 2q$.
\end{enumerate}

\end{restatable}%

[Next verbatim source excerpt]

\subsection{Proof of \Cref{thm:instability}} \label{sec:instability-proof}

Recall \Cref{thm:instability}:
\thminstability*

Putting things together, the first statement of \Cref{thm:instability} comes from \Cref{thm:independent-zero-unstable} and the second from \Cref{thm:sub-unstable-equiv}. The third
has roots in \Cref{lem:unstable-calc}, which allows us to calculate instability.

Rigorously, we can generate a $d$-unstable $\choice$ for any $1 < d \leq 2q$. Take any $1$-unstable, $q$-acceptant $\choice$ induced by a total order $\succ$ (i.e, is a ranking). Let $n$ be $\frac{d}{2}$, rounded down. Then let $X^* = \{x^*_1, ..., x^*_n\}$ be the $n$ \textit{minimal} elements of $\succ$ (or any $n$ elements that are not among the $n$ maximal elements, if there exist any.) Then modify $\choice$ such that if $X^* \subseteq X$ for any input $X$, $X^* \subseteq \choice(X)$. Then $\choice$ is tightly $2n-1$-unstable. To make $\choice$ $2n$-unstable, let there be some $x'$ such that if $X^* \subseteq X$ but $x' \in X$ as well, then \textit{no} element of $X^*$ is in $\choice(X)$. %

The intuition for the construction is that $X^*$ forms a highly desirable team, even though each element is not individually desirable, nor are partial subsets of $X^*$.
In the inconsistent case (even instability) $x'$ represents a candidate whose mere presence negates the value of $X^*$ as a team, even when not accepted.
\end{ReviewVerbatim}


\paragraph{Expanded PaperInterface specification}
\begin{ReviewVerbatim}
∀ (q d : ℕ),
  1 ≤ d →
    d ≤ 2 * q →
      ∃ β inst C,
        DGD26AdmissionsPredictability.paper_choice_function_feasible C ∧
          DGD26AdmissionsPredictability.paper_definition_q_acceptance q C ∧
            DGD26AdmissionsPredictability.paper_definition_tight_d_instability d C
\end{ReviewVerbatim}

\paragraph{Lean proof endpoint (built separately)}\begin{ReviewVerbatim}
DGD26AdmissionsPredictability.PaperInterface.paper_theorem1_tight_all_d_statement
\end{ReviewVerbatim}

\paragraph{Recorded source-to-Spec screening}
\textbf{Verdict:} \texttt{not recorded}
\\
\textbf{Reason:} not recorded
\reviewmatch{Matches source input}
\noindent\textbf{Reviewer annotation}\par
\reviewerbox
\clearpage
\hypertarget{source-claim-b3ee09f1e9e6b10d}{}
\section*{Review row 56}
\textbf{Source locator:} {\footnotesize\raggedright cited publication, lines 1097-\allowbreak{}1109\par}
\paragraph{Verbatim source input}
\begin{ReviewVerbatim}
\begin{restatable}[Linear Assignment Instability]{theorem}{LAPstable}
\label{thm:lap-unstable}
Linear assignment problems are 1-unstable.
\end{restatable}

\textit{Intuition}: The ranking-like structure of linear assignment slots maintains substitutability, because each slot maintains a  preference ordering over elements that is independent of the cohort.

\begin{proof}
Recall that for $q$-acceptant functions, substitutability and 1-instability are equivalent by \Cref{thm:sub-unstable-equiv}.

The proof is by contradiction; assume $\choice$ is not substitutable. As per \Cref{lem:unsubstitutable}, there exists $\choice(X), \choice(X')$ where some $x^* \in X$ is not in $\choice(X)$ but $x^* \in \choice(X')$.
By \Cref{lem:lap-ordering}, $\choice(X) = \{x_{s_1}, ..., x_{s_q}\}$ tells us we have $x_{s_i} \succ_{s_i} x^*$ for all $i \in \{1, ..., q\}$. For $x^*$ to be in $\choice(X')$, $x_{s_i}$ must also be in $\choice(X')$, meaning that $\choice(X) \subseteq \choice(X')$. However, as $x^* \in \choice(X')$, $|\choice(X')| > |\choice(X')|$, which violates $q$-acceptance, leading to contradiction.
\end{proof}
\end{ReviewVerbatim}


\paragraph{Expanded PaperInterface specification}
\begin{ReviewVerbatim}
∀ {α : Type u_1} [inst : DecidableEq α] {σ : Type u_2} [inst_1 : DecidableEq σ] [inst_2 : Fintype σ] {w : α → σ → ℝ}
  (hwell : DGD26AdmissionsPredictability.LAP.Assignment.WellPosedObjective w),
  DGD26AdmissionsPredictability.paper_definition_d_instability 1
    (DGD26AdmissionsPredictability.paperLAPChoiceRule w hwell)
\end{ReviewVerbatim}

\paragraph{Lean proof endpoint (built separately)}\begin{ReviewVerbatim}
DGD26AdmissionsPredictability.PaperInterface.paper_lap_assignment_one_instability_statement
\end{ReviewVerbatim}

\paragraph{Recorded source-to-Spec screening}
\textbf{Verdict:} \texttt{not recorded}
\\
\textbf{Reason:} not recorded
\reviewmatch{Matches source input}
\noindent\textbf{Reviewer annotation}\par
\reviewerbox
\clearpage
\hypertarget{source-claim-a0520310d6532be1}{}
\section*{Review row 57}
\textbf{Source locator:} {\footnotesize\raggedright cited publication, lines 1111-\allowbreak{}1129\par}
\paragraph{Verbatim source input}
\begin{ReviewVerbatim}
\begin{restatable}[Linear Assignment Variability]{theorem}{LAPvariable}
\label{thm:lap-variability}
The variability of a linear assignment problem is upper bounded by the number of distinct preference orderings induced by its slots.
\end{restatable}

\textit{Intuition}:
The proof follows from the observation that, if slot $s$ ranks some other element $x$ in the chosen set higher than the element $x_s$ currently assigned to $s$, $x$ cannot be in the variable set. Then, if $k$ slots (say, slots $s_1, ..., s_k$, with assigned elements $x_{s_1}, ..., x_{s_k}$) induce the same preference ordering over elements, then only the minimal element in $\{x_{s_1}, ..., x_{s_k}\}$ according to that preference ordering can be in $V_\choice(X)$.

\begin{proof}
We want to show that variability is upper bounded by the number of unique orderings. Recall that the variable set $V_\choice(X) \subseteq \choice(X)$ is at most cardinality $q$.

I claim that for any $X$ and $\choice(X)$, $x_{s_i}$ is not in $V_\choice(X)$ if $x_{s_i} \succ_{s_j} x_{s_j}$ for any $j \neq i$.
To be precise, if $x_{s_i} \in \choice(X)$, and $x_{s_i} \succ_{s_j} x_{s_j}$ for some $j \neq i$, then $x_{s_i} \in \choice(X')$ for any $X' = X \cup \{x'\}$.

The proof of this claim is by contradiction: assume there exists $\choice(X), \choice(X')$ where some $x_{s_i} \succ_{s_j} x_{s_j}$ but $x_{s_i} \not\in \choice(X')$. Firstly, because of the 1-instability proven in \Cref{thm:lap-unstable}, $|\choice(X) \setminus \choice(X')| = 1$, which means by assumption, $\choice(X) \setminus \choice(X') = \{x_{s_i}\}$, which further implies $x_{s_j} \in \choice(X')$. However,  \Cref{lem:lap-ordering} then applies to $x_{s_j}$, where $x_{s_i} \succ_{s_j} x_{s_j}$ implies $x_{s_i} \in \choice(X')$, leading to contradiction.

Secondly, note that if $s_i, s_j$ induce the same total ordering, then either $x_{s_i} \succ_{s_j} x_{s_j}$ or $x_{s_j} \succ_{s_i} x_{s_i}$ (and the same for $\succ_{s_i}$) for any $\choice(X)$, as a property of any total ordering. Then one of either $x_{s_j}$ or $x_{s_j}$ are not in $V_\choice(X)$. It is intuitive to see that this generalizes: if $s_1, ... s_j$ all induce the same ordering $\succ_{s_j}$, there is exactly one minimal element in $x_{s_1}, ..., x_{s_j}$, because any finite subset of a totally ordered set is well-ordered.
Say that $x_{s_j}$ is the minimal element, w.l.o.g.; then none of $x_{s_1}, ..., x_{s_{j-1}}$ are in $V_\choice(X)$. More generally, let there be $m$ unique orderings, which we notate $\succ_1,..., \succ_m$. Let $k \in \{1,..., m\}$ and let there be $c_k$ slots that induce $\succ_k$. Then we have
defined sets of size $c_k-1$ elements that cannot be in $V_\choice(X)$. Then $|V_\choice(X)| \leq m$, and therefore the overall variability is at most $m$.
\end{ReviewVerbatim}


\paragraph{Expanded PaperInterface specification}
\begin{ReviewVerbatim}
∀ {α : Type u_1} [inst : DecidableEq α] [inst_1 : Fintype α] {σ : Type u_2} [inst_2 : DecidableEq σ]
  [inst_3 : Fintype σ] {w : α → σ → ℝ} (hwell : DGD26AdmissionsPredictability.LAP.Assignment.WellPosedObjective w),
  (∀ (s : σ), DGD26AdmissionsPredictability.LAP.Assignment.SlotNoTies w s) →
    DGD26AdmissionsPredictability.paper_definition_variability_at_most
      (DGD26AdmissionsPredictability.LAP.Assignment.distinctSlotOrderCount w)
      (DGD26AdmissionsPredictability.paperLAPChoiceRule w hwell)
\end{ReviewVerbatim}

\paragraph{Lean proof endpoint (built separately)}\begin{ReviewVerbatim}
DGD26AdmissionsPredictability.PaperInterface.paper_lap_assignment_slot_order_class_variability_of_unique_global_optima_statement
\end{ReviewVerbatim}

\paragraph{Recorded source-to-Spec screening}
\textbf{Verdict:} \texttt{not recorded}
\\
\textbf{Reason:} not recorded
\reviewmatch{Matches source input}
\noindent\textbf{Reviewer annotation}\par
\reviewerbox
\clearpage
\hypertarget{source-claim-a9c5e5d5c616ab58}{}
\section*{Review row 58}
\textbf{Source locator:} {\footnotesize\raggedright cited publication, lines 910-\allowbreak{}923\par}
\paragraph{Verbatim source input}
\begin{ReviewVerbatim}
\begin{restatable}[$q$-Representativeness and $1$-Variability are Equivalent]{theorem}{rankingvariability}
\label{thm:ranking-variability}
$\choice$ is $q$-representative if and only if it is $q$-acceptant and $1$-unstable with variability $1$.
\end{restatable}
\begin{proof}
We prove this in two parts: first, that $q$-representativeness implies $1$-variability, and then the converse.

\paragraph{Direction 1} %
Note that $q$-representativeness implies $q$-acceptance by definition.

Next, we prove that $q$-representativeness implies $1$-instability. Assume for contradiction that some $q$-representative $\choice$ is not substitutable (and therefore not $1$-unstable). Then there exists (by \Cref{lem:unsubstitutable}) some $x^*$, $x'$, $X$ and $X' = X \cup \{x'\}$ where $x^* \in X$ and $x^* \in \choice(X')$ but $x^* \not\in \choice(X)$. By $q$-acceptance, there must be at least one $x$ in $\choice(X)$ that is not in $\choice(X')$ (which must take the place of $x^*$). However, $\choice(X')$ implies that $x^* \succ x$, while $\choice(X)$ implies that $x \succ x^*$ --- violating asymmetry and leading us to contradiction.

Lastly, proving that $q$-representative $\choice$ leads to $m=1$ is relatively intuitive. For any $X$, $\choice(X)$ are the $q$ highest-ranked elements of $X$. Denote $X = \{x_1, x_2, ..., x_k\}$ where $|X| = k$ and $x_1 \succ x_2 \succ ...\succ x_k$. Then $\choice(X) = \{x_1, ..., x_{q^*}\}$ where $q^* = \min\{q, k\}$.
Assume for contradiction that there is a $q$-representative $\choice$ with variability $m > 1$. Then there exists some $X$ where $|V_\choice(X)| > 1$; say that $\{x^*_a, x^*_b\} \subseteq V_\choice(X)$. Then there exists some $x'_a$ where $\choice(X) \setminus \choice(X \cup x'_a) = \{x^*_a\}$ and similar $x'_b$ for $x^*_b$. The first implies $x^*_b$ is ranked above $x^*_a$, but the second implies the opposite. Then choices cannot be made according to a static ranking and $\choice$ is not $q$-representative, completing the proof by contradiction. %
\end{ReviewVerbatim}


\paragraph{Expanded PaperInterface specification}
\begin{ReviewVerbatim}
∀ {α : Type u_1} [inst : DecidableEq α] [inst_1 : Fintype α] {q : ℕ}
  {C : DGD26AdmissionsPredictability.PaperChoiceRule α},
  DGD26AdmissionsPredictability.paper_choice_function_feasible C →
    DGD26AdmissionsPredictability.paper_definition_q_representativeness q C →
      0 < q →
        q < Fintype.card α →
          DGD26AdmissionsPredictability.paper_definition_q_acceptance q C ∧
            DGD26AdmissionsPredictability.paper_definition_d_instability 1 C ∧
              DGD26AdmissionsPredictability.paper_definition_variability_exactly 1 C
\end{ReviewVerbatim}

\paragraph{Lean proof endpoint (built separately)}\begin{ReviewVerbatim}
DGD26AdmissionsPredictability.PaperInterface.paper_q_representative_forward_exact_statement
\end{ReviewVerbatim}

\paragraph{Recorded source-to-Spec screening}
\textbf{Verdict:} \texttt{not recorded}
\\
\textbf{Reason:} not recorded
\reviewmatch{Matches source input}
\noindent\textbf{Reviewer annotation}\par
\reviewerbox
\clearpage
\hypertarget{source-claim-3d12c759820fa6e4}{}
\section*{Review row 59}
\textbf{Source locator:} {\footnotesize\raggedright cited publication, lines 910-\allowbreak{}950\par}
\paragraph{Verbatim source input}
\begin{ReviewVerbatim}
\begin{restatable}[$q$-Representativeness and $1$-Variability are Equivalent]{theorem}{rankingvariability}
\label{thm:ranking-variability}
$\choice$ is $q$-representative if and only if it is $q$-acceptant and $1$-unstable with variability $1$.
\end{restatable}
\begin{proof}
We prove this in two parts: first, that $q$-representativeness implies $1$-variability, and then the converse.

\paragraph{Direction 1} %
Note that $q$-representativeness implies $q$-acceptance by definition.

Next, we prove that $q$-representativeness implies $1$-instability. Assume for contradiction that some $q$-representative $\choice$ is not substitutable (and therefore not $1$-unstable). Then there exists (by \Cref{lem:unsubstitutable}) some $x^*$, $x'$, $X$ and $X' = X \cup \{x'\}$ where $x^* \in X$ and $x^* \in \choice(X')$ but $x^* \not\in \choice(X)$. By $q$-acceptance, there must be at least one $x$ in $\choice(X)$ that is not in $\choice(X')$ (which must take the place of $x^*$). However, $\choice(X')$ implies that $x^* \succ x$, while $\choice(X)$ implies that $x \succ x^*$ --- violating asymmetry and leading us to contradiction.

Lastly, proving that $q$-representative $\choice$ leads to $m=1$ is relatively intuitive. For any $X$, $\choice(X)$ are the $q$ highest-ranked elements of $X$. Denote $X = \{x_1, x_2, ..., x_k\}$ where $|X| = k$ and $x_1 \succ x_2 \succ ...\succ x_k$. Then $\choice(X) = \{x_1, ..., x_{q^*}\}$ where $q^* = \min\{q, k\}$.
Assume for contradiction that there is a $q$-representative $\choice$ with variability $m > 1$. Then there exists some $X$ where $|V_\choice(X)| > 1$; say that $\{x^*_a, x^*_b\} \subseteq V_\choice(X)$. Then there exists some $x'_a$ where $\choice(X) \setminus \choice(X \cup x'_a) = \{x^*_a\}$ and similar $x'_b$ for $x^*_b$. The first implies $x^*_b$ is ranked above $x^*_a$, but the second implies the opposite. Then choices cannot be made according to a static ranking and $\choice$ is not $q$-representative, completing the proof by contradiction. %

\paragraph{Direction 2}
Assume $q$-acceptance, variability 1, and instability 1. Now we prove $\choice$ is $q$-representative. To be precise, I take this to mean that $\choice$ is indistinguishable from some $q$-representative $\choice$; $\choice$ could be represented by some total ordering over elements of $X$. We show this strict total ordering constructively.

Define $\succ$ such that, if $x_a \in \choice(X)$ but $x_b \not\in\choice(X)$ for any $x_a, x_b \in X$, $x_a \succ x_b$. Note that this implies that the existence of some set $A$ where $x_a \in \choice(A)$ and $A_\choice(A) = \{x_b\}$ by consistency (e.g. $A = \choice(X) \cup \{x_b\}$). This also implies $A'$ where $x_a, x_b \in \choice(A')$, $V_\choice(A') = \{x_b\}$, where $A' = A \setminus \{x\}$ for any $x \neq x_a, x_b$.

Then, we need to prove this relation is asymmetric, transitive, and completeness (or rather, is indistinguishable from a complete relation when examining the induced choice function).

\paragraph{Asymmetry}
First, we show asymmetry. Assume for contradiction that $a \succ b$ and $b \succ a$. Then there exists $X_a$ where $a, b \in X_a$, $a \in \choice(X_a)$ and $b \not\in \choice(X_a)$, and a similar $X_b$ where $a \not\in \choice(X_b)$ and so on.

Let $X_a = \choice(X_a) \cup \{b\}$ (we can construct such out of any $X_a$, by consistency), where $A_\choice(X_a) = \{b\}$, and denote $\choice(X_a) = \{x^a_1, ..., x^a_{q-1}, a\}$. Construct $X_b$ similarly. Further, let $X'_a = X_a \setminus \{x^a_1\}$, and note that $V_\choice(X_a') = \{b\}$.

Then from here, we add the elements of $X_b$ to $X_a'$, constructing $X_1, X_2,$ and so forth.

Let $X_0 = X_a'$ and sequentially construct a series of $X_t$. Add $x^b_t$ to $X_{t-1}$. If $x_b \not\in \choice(X_{t-1} \cup \{x^b_t\})$, then let $X_t = X_{t-1} \cup \{x^b_t\}$, and $V_\choice(X_t) = \{b\}$ by \Cref{lem:consistent-remove}. If $x^b_t \in \choice(X_{t-1} \cup \{x^b_t\})$, then $b \not\in \choice(X_{t-1} \cup \{x^b_t\})$ and $A_\choice(X_{t-1} \cup \{x^b_t\}) = \{b\}$ (by \Cref{lem:a-v-equal}). Then let $X_t = \{x^b_t\} \cup X_a' \setminus x^a_i$ where $x^a_{i-1} \not\in X_{t-1}$ but $x^a_i \in X_{t-1}$, and $V_\choice(X_t) = \{b\}$ (again, by \Cref{lem:a-v-equal}). Note that $a \in \choice(X_t)$, $b \in X_t$, and specifically $b \in V_\choice(X_t)$, for all $t \in \{0,..., q-2\}$.

Then $X_{q-2}$ will contain $x^b_1,...x^b_{q-2}$. Let $X^* = X_{q-2} \cup \{x^b_{q-1}\}$. Note that adding $x^b_{q-1}$ to $X_{q-2}$ makes $X^* \supseteq X_b$. Then, if $a \in \choice(X^*)$, substitutability is violated, as $a \not\in \choice(X_b)$. If $x^b_{q-1} \not\in \choice(X^*)$, then $\choice(X_{q-2}) = \choice(X^*)$ by consistency, and $a \in \choice(X_{q-2})$. If instead $x^b_{q-1} \in \choice(X^*)$, recall that $V_\choice(X_{q-2}) = \{b\}$, so $a \not\in V_\choice(X_{q-2})$, and $a$ remains in $\choice(X^*)$. Therefore, either possibility leads to contradiction, proving that $\succ$ is asymmetric.

\paragraph{Transitivity}
Say that $a \succ b$ and $b \succ c$ and we want to prove $a \succ c$. If $a \succ b$, there exists some $X$ where $a, b \in X$, $a \in \choice(X)$, and $b \not\in \choice(X)$. Let $X$ be $\choice(X) \cup \{b\}$, where $a \in \choice(X), b \not\in\choice(X)$ by consistency. Then consider $X \cup \{c\}$. By asymmetry, $c \not\succ b$, so $c \not\in \choice(X\cup \{c\})$. By consistency, then, $\choice(X\cup \{c\}) = \choice(X)$, so $a \in \choice(X\cup \{c\})$ and $a \succ c$, completing the proof.

\paragraph{(Almost)-Completeness}
See that this relation must exist for all but $q$ elements --- and any arbitrary linear extension thereof generates the same choice function behavior.

For any $a, b$, let there be some $X_0$ where $a, b \not\in \choice(X_0)$. Then remove elements of $\choice(X_t)$ in a sequence $X_0, ..., X_t$ until $A_\choice(X_t) = \{a\}$ or $\{b\}$. This must happen for some value of $t$ because when $|X_t| = q+1$, at least one of $a$ or $b$ are in $\choice(X)$, and due to $1$-instability, both $a$ and $b$ cannot be added to $\choice(X_t)$ at once. The only exception is if there is no such $X_0$ where $a, b \not\in \choice(X_0)$, which can only be true if $a$ is such that $a \in \choice(X)$ for all $X$ (and similarly for $b$). (If $a \not\in \choice(X)$ for some $X$, then substitutability implies $a \not\in \choice(X \cup \{b\})$). However, there can only be $q$ such elements (by $q$-acceptance). Note that any linear extension of this partial order will produce the same choice function behavior; that is, any constructed total ordering that otherwise aligns with $\succ$ will produce the same choice behavior, and thus, $\choice$ can be represented by any such strict total ordering.
\end{proof}
\end{ReviewVerbatim}


\paragraph{Expanded PaperInterface specification}
\begin{ReviewVerbatim}
∀ {α : Type u_1} [inst : DecidableEq α] [inst_1 : Fintype α] {q : ℕ}
  {C : DGD26AdmissionsPredictability.PaperChoiceRule α},
  DGD26AdmissionsPredictability.paper_choice_function_feasible C →
    DGD26AdmissionsPredictability.paper_definition_q_acceptance q C →
      DGD26AdmissionsPredictability.paper_definition_d_instability 1 C →
        DGD26AdmissionsPredictability.paper_definition_variability_exactly 1 C →
          0 < q → q < Fintype.card α → DGD26AdmissionsPredictability.paper_definition_q_representativeness q C
\end{ReviewVerbatim}

\paragraph{Lean proof endpoint (built separately)}\begin{ReviewVerbatim}
DGD26AdmissionsPredictability.PaperInterface.paper_q_representative_converse_exact_statement
\end{ReviewVerbatim}

\paragraph{Recorded source-to-Spec screening}
\textbf{Verdict:} \texttt{not recorded}
\\
\textbf{Reason:} not recorded
\reviewmatch{Matches source input}
\noindent\textbf{Reviewer annotation}\par
\reviewerbox
\clearpage
\hypertarget{source-claim-5fcefe7f9430bc35}{}
\section*{Review row 60}
\textbf{Source locator:} {\footnotesize\raggedright cited publication, lines 962-\allowbreak{}996\par}
\paragraph{Verbatim source input}
\begin{ReviewVerbatim}
\paragraph{Sequentially Composed Substitutable Functions Are Substitutable}
Similarly, we show sequential compositions of substitutable functions are also substitutable.
\begin{restatable}[Composition Preserves Substitutability]{theorem}{seqsub}
\label{thm:composition-substitutable}
Sequential composition of substitutable functions yields a substitutable function.
\end{restatable}

\textit{Intuition}:
``Appending'' a substitutable choice function $\choice_2$ before another substitutable function $\choice_1$ affects the inputs to $\choice_1$ in a structured way.
\begin{proof}
    Proof by induction.
    In the base case, for $n = 1$, $\choice(X) = \choice_1(X)$, and as $\choice_1$ is given to be substitutable, $\choice$ is as well.

    For the inductive step, show that $\choice^{n+1}$ is substitutable when comprised of $\choice_{n+1}, \choice_n, ..., \choice_1$, assuming that $\choice^n$ comprised of $\choice_n,...,\choice_1$ is substitutable.

    The definition of substitutability is that $X \subseteq X'$ implies $\choice(X') \cap X \subseteq \choice(X)$. So we need to show $\choice^{n+1}(X') \cap X \subseteq \choice^{n+1}(X)$.

    First, note that $\choice^{n+1}(X) = \choice_{n+1}(X) \cup \choice^n(X \setminus \choice_{n+1}(X))$.
    Let $S = X \setminus \choice_{n+1}(X)$ and $S' = X' \setminus \choice_{n+1}(X')$. As $S \subseteq X'$, we only need to show that no elements of $S$ are elements of $\choice_{n+1}(X')$. As $\choice_{n+1}$ is substitutable, $\choice_{n+1}(X') \cap X = \choice_{n+1}(X)$. So, any element of $X$ that is in $\choice_{n+1}(X') $ must be in $\choice_{n+1}(X)$, and therefore cannot be in $S = X \setminus \choice_{n+1}(X)$. Therefore, $S \subseteq S'$.

    Now, we can apply the induction hypothesis, so we assume $\choice^n$ is substitutable to get

    \begin{align*}
        \choice^n(S') \cap S & \subseteq \choice^n(S) && \text{induction hypothesis} \\
        \choice^n(S') \cap \left(X \setminus \choice_{n+1}(X)\right) & \subseteq \choice^n(S) &&  \text{substitute out $S$}\\      \left(\choice^n(S') \cap X\right) \setminus \left(\choice^n(S') \cap  \choice_{n+1}(X)\right) & \subseteq \choice^n(S) && \text{intersection commutes}\\
        \choice_{n+1}(X) \cup \left[\left(\choice^n(S') \cap X\right) \setminus \left(\choice^n(S') \cap  \choice_{n+1}(X)\right)\right] & \subseteq \choice_{n+1}(X) \cup \choice^n(S) && \text{add $\choice_{n+1}(X)$} \\ \intertext{Note that $\choice_{n+1}(X) \supseteq\left(\choice^n(S') \cap  \choice_{n+1}(X)\right)$, so the set difference can be simplified out:}
        \choice_{n+1}(X) \cup \left(\choice^n(S') \cap X\right) & \subseteq \choice_{n+1}(X) \cup \choice^n(S) && \text{equivalent} \\ \intertext{Note that for any sets $A,B,C$, $(A \cup B) \cap C = (A \cap C)\cup (B \cap C) \subseteq A \cup (B \cap C)$, so:
        }
        \left(\choice_{n+1}(X) \cup \choice^n(S') \right) \cap X & \subseteq \choice_{n+1}(X) \cup \choice^n(S) && \text{} \\
        \left(\choice_{n+1}(X) \cup \choice^n(X' \setminus \choice_{n+1}(X')) \right) \cap X & \subseteq \choice_{n+1}(X) \cup \choice^n(X \setminus \choice_{n+1}(X)) &&  \text{substitute out $S$, $S'$}  \\
        \choice^{n+1}(X') \cap X & \subseteq \choice^{n+1}(X) &&  \text{definition of $\choice^{n+1}$}
    \end{align*}

    Which is the definition of substitutability, completing the proof.
\end{proof}
\end{ReviewVerbatim}


\paragraph{Expanded PaperInterface specification}
\begin{ReviewVerbatim}
∀ {α : Type u_1} [inst : DecidableEq α] {Cs : List (DGD26AdmissionsPredictability.PaperChoiceRule α)},
  (∀ C ∈ Cs, DGD26AdmissionsPredictability.paper_choice_function_feasible C) →
    (∀ C ∈ Cs, DGD26AdmissionsPredictability.paper_definition_substitutability C) →
      DGD26AdmissionsPredictability.paper_definition_substitutability
        (DGD26AdmissionsPredictability.paper_definition_sequential_composition Cs)
\end{ReviewVerbatim}

\paragraph{Lean proof endpoint (built separately)}\begin{ReviewVerbatim}
DGD26AdmissionsPredictability.PaperInterface.paper_sequential_composition_substitutable_statement
\end{ReviewVerbatim}

\paragraph{Recorded source-to-Spec screening}
\textbf{Verdict:} \texttt{not recorded}
\\
\textbf{Reason:} not recorded
\reviewmatch{Matches source input}
\noindent\textbf{Reviewer annotation}\par
\reviewerbox
\clearpage
\hypertarget{source-claim-b564e283fe9ce18d}{}
\section*{Review row 61}
\textbf{Source locator:} {\footnotesize\raggedright cited publication, lines 733-\allowbreak{}745\par}
\paragraph{Verbatim source input}
\begin{ReviewVerbatim}
\begin{restatable}{theorem}{BoundedDistance}
\label{thm:sub-accept-consistent}
Any $q$-acceptant, substitutable $\choice$ is consistent.
\end{restatable}

\begin{proof}
Consider $X_1, X_2$ where $\choice(X_2) \subseteq X_1 \subseteq X_2$ for substitutable, $q-$acceptant $\choice$. Substitutability guarantees that if $x \in \choice(X_2)$ and $x \in X_1$, $x \in \choice(X_1)$. By construction, this applies to all $x \in \choice(X_2)$, because $\choice(X_2) \subseteq X_1$; therefore, $\choice(X_2) \subseteq \choice(X_1)$. As $|\choice(X_2)| = |\choice(X_1)|$, and $\choice(X_1),\choice(X_2)$ are finite sets, this means that $\choice(X_2) = \choice(X_1)$, which is the definition of consistency.
\end{proof}

\subsection{\textit{d}-Instability} \label{sec:d-unstable-appendix}

\subsubsection{Proof of \Cref{lem:unstable-calc}} \label{sec:calc-proof}
\end{ReviewVerbatim}


\paragraph{Expanded PaperInterface specification}
\begin{ReviewVerbatim}
∀ {α : Type u_1} [inst : DecidableEq α] {q : ℕ} {C : DGD26AdmissionsPredictability.PaperChoiceRule α},
  DGD26AdmissionsPredictability.paper_definition_q_acceptance q C →
    DGD26AdmissionsPredictability.paper_definition_substitutability C →
      DGD26AdmissionsPredictability.paper_definition_consistency C
\end{ReviewVerbatim}

\paragraph{Lean proof endpoint (built separately)}\begin{ReviewVerbatim}
DGD26AdmissionsPredictability.PaperInterface.paper_q_acceptant_substitutable_consistent_statement
\end{ReviewVerbatim}

\paragraph{Recorded source-to-Spec screening}
\textbf{Verdict:} \texttt{not recorded}
\\
\textbf{Reason:} not recorded
\reviewmatch{Matches source input}
\noindent\textbf{Reviewer annotation}\par
\reviewerbox
\clearpage
\hypertarget{source-claim-7c749451f4caf111}{}
\section*{Review row 62}
\textbf{Source locator:} {\footnotesize\raggedright cited publication, lines 315-\allowbreak{}322\par}
\paragraph{Verbatim source input}
\begin{ReviewVerbatim}
\begin{restatable}{theorem}{thmvariability}
\label{thm:variability}
Consider a $q$-acceptant, 1-unstable choice function $\mathcal C$ which
can be represented as the sequential composition of $n$ choice functions, each characterized by a total order.
Then $\mathcal C$ has a variability  $m$, where $1 \leq m \leq n$. Furthermore, $\mathcal C$ has variability $1$ if and only if it can be characterized by a single total order, $n=1$.
\end{restatable}
In the appendix, we discuss when this upper bound is tight. We also extend this result to choice functions which are not explicitly defined as a fixed sequential composition, such as linear assignment problems.
\Cref{thm:lap-variability} addresses the context-dependent sequential queues which result from such optimization based admissions.

[Next verbatim source excerpt]

\subsection{Proof of \Cref{thm:variability}} \label{sec:variability-proof}

Recall \Cref{thm:variability}:
\thmvariability*

We see that the second claim comes exactly from \Cref{thm:ranking-variability}, and that combining it with \Cref{thm:additive-variability} directly gives us the first claim.
\end{ReviewVerbatim}

\paragraph{Approved corrected review target}
The archival source above is not asserted equivalent to this replacement.
\begin{ReviewVerbatim}
For a sequential composition of total-order q-representative queues with positive total capacity on a universe larger than that capacity, there exists m with 1≤m≤n such that the composition has exact variability m.
\end{ReviewVerbatim}

\textbf{Archival source anchor:} {\footnotesize\raggedright cited publication, lines 315-\allowbreak{}322\par}
\paragraph{Recorded basis}
DGD source clarification:\allowbreak{} Theorem 2's range is read on the positive-\allowbreak{}capacity, nontrivial queue domain.\allowbreak{}
\paragraph{Expanded PaperInterface specification}
\begin{ReviewVerbatim}
∀ {α : Type u_1} [inst : DecidableEq α] [inst_1 : Fintype α] {qs : List ℕ}
  {Cs : List (DGD26AdmissionsPredictability.PaperChoiceRule α)},
  List.Forall₂
      (fun q C =>
        DGD26AdmissionsPredictability.paper_choice_function_feasible C ∧
          DGD26AdmissionsPredictability.paper_definition_q_representativeness q C)
      qs Cs →
    0 < qs.sum →
      qs.sum < Fintype.card α →
        ∃ m,
          1 ≤ m ∧
            m ≤ Cs.length ∧
              DGD26AdmissionsPredictability.paper_definition_variability_exactly m
                (DGD26AdmissionsPredictability.paper_definition_sequential_composition Cs)
\end{ReviewVerbatim}

\paragraph{Lean proof endpoint (built separately)}\begin{ReviewVerbatim}
DGD26AdmissionsPredictability.PaperInterface.paper_sequential_q_representative_variability_range_statement
\end{ReviewVerbatim}

\paragraph{Recorded source-to-Spec screening}
\textbf{Verdict:} \texttt{not recorded}
\\
\textbf{Reason:} not recorded
\reviewmatch{Matches approved corrected target}
\noindent\textbf{Reviewer annotation}\par
\reviewerbox
\clearpage
\hypertarget{source-claim-26bfa3509b616c5d}{}
\section*{Review row 63}
\textbf{Source locator:} {\footnotesize\raggedright cited publication, lines 315-\allowbreak{}323\par}
\paragraph{Verbatim source input}
\begin{ReviewVerbatim}
\begin{restatable}{theorem}{thmvariability}
\label{thm:variability}
Consider a $q$-acceptant, 1-unstable choice function $\mathcal C$ which
can be represented as the sequential composition of $n$ choice functions, each characterized by a total order.
Then $\mathcal C$ has a variability  $m$, where $1 \leq m \leq n$. Furthermore, $\mathcal C$ has variability $1$ if and only if it can be characterized by a single total order, $n=1$.
\end{restatable}
In the appendix, we discuss when this upper bound is tight. We also extend this result to choice functions which are not explicitly defined as a fixed sequential composition, such as linear assignment problems.
\Cref{thm:lap-variability} addresses the context-dependent sequential queues which result from such optimization based admissions.
Importantly, Theorem~\ref{thm:variability} means that a choice function can be faithfully represented by a machine learning model if and \emph{only if} it is characterized by a total ordering over $\calX$.

[Next verbatim source excerpt]

\subsection{Proof of \Cref{thm:variability}} \label{sec:variability-proof}

Recall \Cref{thm:variability}:
\thmvariability*

We see that the second claim comes exactly from \Cref{thm:ranking-variability}, and that combining it with \Cref{thm:additive-variability} directly gives us the first claim.
\end{ReviewVerbatim}

\paragraph{Approved corrected review target}
The archival source above is not asserted equivalent to this replacement.
\begin{ReviewVerbatim}
For a feasible q-acceptant 1-unstable choice rule on a positive-capacity universe larger than q, exact variability one holds if and only if the rule is q-representative, equivalently characterized by one strict total order.
\end{ReviewVerbatim}

\textbf{Archival source anchor:} {\footnotesize\raggedright cited publication, lines 315-\allowbreak{}323\par}
\paragraph{Recorded basis}
DGD source clarification:\allowbreak{} variability one is the one-\allowbreak{}total-\allowbreak{}order characterization on the nondegenerate queue domain.\allowbreak{}
\paragraph{Expanded PaperInterface specification}
\begin{ReviewVerbatim}
∀ {α : Type u_1} [inst : DecidableEq α] [inst_1 : Fintype α] {q : ℕ}
  {C : DGD26AdmissionsPredictability.PaperChoiceRule α},
  DGD26AdmissionsPredictability.paper_choice_function_feasible C →
    DGD26AdmissionsPredictability.paper_definition_q_acceptance q C →
      DGD26AdmissionsPredictability.paper_definition_d_instability 1 C →
        0 < q →
          q < Fintype.card α →
            (DGD26AdmissionsPredictability.paper_definition_variability_exactly 1 C ↔
              DGD26AdmissionsPredictability.paper_definition_q_representativeness q C)
\end{ReviewVerbatim}

\paragraph{Lean proof endpoint (built separately)}\begin{ReviewVerbatim}
DGD26AdmissionsPredictability.PaperInterface.paper_variability_one_iff_single_order_statement
\end{ReviewVerbatim}

\paragraph{Recorded source-to-Spec screening}
\textbf{Verdict:} \texttt{not recorded}
\\
\textbf{Reason:} not recorded
\reviewmatch{Matches approved corrected target}
\noindent\textbf{Reviewer annotation}\par
\reviewerbox

\end{document}
